% Stage H: the search region: the search intervals of the caller with their default, the range of
% the denominators, the Sturm step that narrows the intervals, RATPOINTS_NO_REVERSE_AUTO and
% RATPOINTS_BAD_ARGS; the doubles of the C code and the exact bounds of the numerators; and the
% options -l -u -dl -du of the command-line program with messages 3 and 8. Proposed names are set
% with \lname and become \lean once the declarations exist.
\chapter{The search region}\label{chap:region}

Stage~H models the search region of ratpoints. The caller of \verb|find_points| may give search
intervals for $x/z$, closed intervals whose ends are C \verb|double|s, and a range of the
denominators $z$, from $b_{\mathrm{low}}$ to $b_{\mathrm{high}}$; when it gives no intervals, the
library searches the interval $[-h, h]$, where $h$ is the double of $H$, moved up to the next double
when it lies below $H$. Before the sieve, the Sturm step of \verb|sturm.c| narrows the search
intervals: from the remainder chain of $f$ and $f'$, whose test for a repeated factor stage~D
modelled (Chapter~\ref{chap:setup}), it counts the sign changes of the chain at $\pm\infty$, bisects
$[-\infty, +\infty]$ to a depth of at most $55$, keeping the intervals on which $f$ may take a value
$\ge 0$, and intersects them with the search intervals. The sieve of a denominator $b$ then runs
over the numerators $\lceil b l \rceil, \dots, \lfloor b u \rfloor$ within the height bound of each
interval $[l, u]$, computed exactly from the doubles of the ends; the range of the denominators
bounds the loop over the denominators and the divisors of the leading coefficient; a search region
other than the whole box forbids the reversal (\verb|RATPOINTS_NO_REVERSE_AUTO|); and search
intervals $[l, u]$ without $l \le u$, or not increasing and strictly apart (each upper end below the
next lower end), or with a NaN end, are refused (\verb|RATPOINTS_BAD_ARGS|). The command-line
program reads the intervals with \verb|-l| and \verb|-u| and the range of the denominators with
\verb|-dl| and \verb|-du|, passes over a point that the library finds outside the exact region of
its command line, and prints the search region of its command line, when there is one, before the
search and the search intervals as the library leaves them after it (messages 3 and~8).

Stage~H models all of this as the C code has it. A double that holds an end of a search interval (or
that the program reads with \verb|strtod|) is a value of the new type $\mathrm{Double}$ of the
model, a structure over its IEEE~754 binary64 bits, with a model function for each operation that
the code modelled here applies to such a double: the comparisons of the C language, in which a NaN
is unordered and $-0.0$ equals $0.0$, the negation and \verb|fabs|, the conversions between
\verb|long|s and doubles, \verb|frexp|, \verb|ldexp|, \verb|nextafter| towards $\pm$\verb|HUGE_VAL|,
\verb|strtod| and \verb|printf| under \verb|%f|. Apart from the estimate of the work in the choice
of the moduli (\verb|numerators_for| and \verb|run_shape|, which multiply the ends by $b$ and
subtract them in floating-point arithmetic, and are stage~I's), ratpoints does no other arithmetic
on the ends of its intervals, so that everything the model does with them is computed exactly, on
integers and rational numbers. The options of the search (Definition~\ref{def:Options}) gain the
search intervals as a list of pairs of doubles, empty by default, and the least and the largest
denominator as integers, $1$ and $0$ by default, a largest denominator below $1$ meaning $H$, as for
the library; the search with a callback normalizes them as \verb|find_points_work_1| does, refuses
intervals that are not well formed with the new refusal $\mathrm{badArgs}$ (error~9 of the program,
which no command line whose words have fewer than $2^{56}$ characters reaches,
Theorem~\ref{thm:cliRegion}), runs the Sturm step on the search intervals, and returns the intervals
that it leaves, for message~8. The command-line program now reads \verb|-l|, \verb|-u|, \verb|-dl|
and \verb|-du|, the last options of ratpoints that it did not read, so that the exit status~12 for
an option not supported yet goes; as before, the arguments of \verb|-t|, the number of threads, and
of the tuning options that stage~I models are checked and then ignored.

For the default options the output of $\mathrm{findPoints}$ does not change: the default interval
gives exactly the numerators $-H, \dots, H$ for every denominator, and the Sturm step leaves out
only numerators $a$ for which $F(a, b) < 0$, which give no point. Under \verb|RATPOINTS_NO_CHECK|,
where the list holds the survivors of the sieve, the survivors now lie in the intervals that the
Sturm step keeps, as in the C program, so that they depend on the depth of the step. \verb|Params|
and \verb|Params.WF| do not change. The specification gains two files, which the challenge restates
verbatim: \verb|Ratpoints/Spec/Region.lean|, with the value of a double, the search intervals that
ratpoints accepts and the predicate that a point lies in the search region, and
\verb|Ratpoints/Spec/Sturm.lean|, with Sturm's theorem as a proposition; and the points in the order
of ratpoints (Definition~\ref{def:orderedPoints}) take the largest denominator of the search. Of the
twelve statements that comparator certifies, six change: membership and the order theorem hold for
the points in the search region, and so does the coverage under \verb|RATPOINTS_NO_CHECK|, all three
assuming Sturm's theorem when the Sturm step is on; the range statement under
\verb|RATPOINTS_NO_CHECK| adds that every member lies in the search region, without that assumption;
and both refusal theorems add the search intervals that are not well formed, while in the one with
the Sturm step off the refusal for an allowed reversal with $d$ even and $c_0 = c_1 = 0$ now also
needs the search region to be the whole box: any other search region forbids the reversal
(\verb|RATPOINTS_NO_REVERSE_AUTO|), and with it the check that makes this refusal. The other six
keep their text: the two of the brute-force reference, no duplicates, which now covers the search
region too, the well-formedness of the parameters of the program, the sublist under
\verb|RATPOINTS_NO_CHECK|, whose two searches share the search region and the intervals of the Sturm
step, and the protocol of the delivery.

The Sturm step keeps every point because of Sturm's theorem: the number of roots of a squarefree
polynomial in an interval $(a, b]$ is the difference of the numbers of sign changes of its Sturm
sequence at $a$ and at $b$. Mathlib has the Sturm sequence (\verb|Polynomial.sturmSeq|) and the sign
changes of a list (\verb|List.signVariations|), but not yet the theorem. The specification states it
as a proposition, $\mathrm{SturmTheorem}$, and the three statements that need it take the hypothesis
$0 \le \mathrm{sturm} \to \mathrm{SturmTheorem}$; it is not an axiom, and the proofs and comparator
check that the standard axioms are the only ones. The hypothesis stays until Mathlib has the
theorem; proving it in the project remains possible. Everything else that the step needs is proved:
that the remainder chain of the C code is the Sturm sequence of $f$ and $f'$ up to positive factors,
the counts at $\pm\infty$ and at the roots of $f$ that the C code reads, the coverage of the
bisection and of the intersection, the exactness of the ends as doubles up to the depth $55$, and
the bounds of the arrays of the step, which the C code relies on.

As in the other stages, since the aim of the project is to verify the C code as far as that can be
done, the model follows the C code where the C code has its data and its tests: the ends of the
intervals are the C code's doubles, compared as the C language compares them, the bounds of the
numerators are computed in the C code's two ways, on the words of a \verb|long| where that is exact
and in 128 bits otherwise, the bisection keeps its \verb|long| numerators and exponents, and the
program reads its ends with \verb|strtod| and the exact rational number of the word, as
\verb|main.c| does. The C code stays at release 3.1.1 while the project runs, so that the project
certifies that release. The statements about the points, and the contracts above the bounds of the
numerators, speak about the values of the ends, rational numbers or $\pm\infty$
($\mathrm{InRegion}$, $\mathrm{InIntervals}$ and the list of the search over $\mathrm{doubleValue}$
of the ends). Only the value of a double, which decodes its bits, and well-formedness read the bits:
well-formedness is the test of \verb|RATPOINTS_BAD_ARGS|, that no end is a NaN and that the values
form search intervals ($\mathrm{IntervalsWF}$ on the $64$-bit words), and it appears in the two
refusal statements, as the hypothesis of those contracts, and in the conclusions that preserve it
(the Sturm step, the join of the program's intervals). A later release whose ends are of another
kind would replace these two, $\mathrm{doubleValue}$ by the value of its own ends and
$\mathrm{IntervalsWF}$ by a test of them, such as $\mathrm{AreSearchIntervals}$ of their values; the
rest would serve it unchanged.

\section{What the C code does}\label{sec:regionC}

\paragraph{The search region.} The arguments of the library hold the search intervals, an array
\verb|domain| of the type \verb|ratpoints_interval| (\csource{ratpoints.h}{359}{359}) with the
number \verb|num_inter| of its intervals, the range of the denominators \verb|b_low| and
\verb|b_high|, and the depth \verb|sturm| of the Sturm step (\csource{ratpoints.h}{363}{364},
\csource{ratpoints.h}{369}{369}, lines that they share with fields of the parameters and of the
choice of the moduli). The manual asks the caller for room in \verb|domain| for
$\mathrm{num\_inter} + d$ intervals, which the Sturm step may write, and says that \verb|b_low| and
\verb|b_high| at most $0$ stand for $1$ and $H$. \verb|find_points_work_1| normalizes them: a
negative \verb|num_inter| becomes $0$ (\csource{find_points.c}{4641}{4641}), \verb|b_low| below $1$
becomes $1$, \verb|b_high| below $1$ or above $H$ becomes $H$ (\csource{find_points.c}{4659}{4661};
$b_{\mathrm{low}} > b_{\mathrm{high}}$ is not refused and leaves no denominator), and a depth above
\verb|DBL_MANT_DIG + 2| $= 55$ becomes $55$ (\csource{find_points.c}{4682}{4682}), beyond which the
ends of the intervals of the Sturm step would not all be doubles. Then it sets the private bit
\verb|RATPOINTS_NO_REVERSE_AUTO| of the flags (\csource{ratpoints.h}{421}{421}) when the caller
gives search intervals, $b_{\mathrm{low}} > 1$ or $b_{\mathrm{high}} < H$, after the normalization
(\csource{find_points.c}{4686}{4687}); the reversal check runs only when neither this bit nor
\verb|RATPOINTS_NO_REVERSE| is set (line 4739), and the mask of the bits of the input clears it on
the next call (line 4606). \verb|find_points_work| saves the range and the depth of the caller
before and restores them after (lines 4542-4545 and 4560-4563), so that the caller does not see the
normalization; \verb|num_inter| and \verb|domain| are not restored, as they describe the region
searched.

\paragraph{The search intervals.} After the reversal check and the test of the \verb|long|s, the
search intervals (\csource{find_points.c}{4829}{4855}): a null pointer \verb|domain| is refused as
\verb|RATPOINTS_BAD_ARGS| (line 4830, a check of the interface); without intervals the library takes
$[-h, h]$, where $h$ is the double of $H$, moved up by \verb|nextafter| when $h < 2^{63}$ and the
\verb|long| of $h$ is below $H$, so that the interval holds every $x = a/b$ of the search; and
otherwise it refuses as \verb|RATPOINTS_BAD_ARGS| an interval without $l \le u$ or one that does not
begin above the end of the interval before it, comparisons that are false for a NaN, so that a NaN
end is refused too; infinite ends, degenerate intervals $[a, a]$ and any number of intervals are
accepted. So the refusals come in the order: the length and the height (the program), the leading
zeros, the reversal check, the search intervals, and the remainder chain of the Sturm step.

\paragraph{The Sturm step.} \verb|_ratpoints_compute_sturm| (\csource{sturm.c}{86}{332}) takes the
polynomial searched, the depth and the search intervals. After the remainder chain
(\csource{sturm.c}{99}{163}, stage~D), whose rows are positive rational multiples of the Sturm
sequence of $f$ and $f'$, it counts the sign changes of the leading coefficients of the rows, at
$+\infty$ as they are and at $-\infty$ with the signs of the rows of odd degree flipped
(\csource{sturm.c}{165}{175}). When the counts are equal, so that $f$ has no real root, and $c_0 <
0$, $f$ is negative everywhere: it sets \verb|num_inter| to $0$ and returns $0$
(\csource{sturm.c}{177}{185}); this happens only for even degree with $c_d < 0$. Otherwise it
bisects $[-\infty, +\infty]$, the ends $\mp$\verb|LONG_MAX| standing for $\mp\infty$
(\csource{sturm.c}{189}{194}), by the nested function \verb|iterate| of gcc
(\csource{sturm.c}{197}{267}), on an interval with the ends $n_l/2^{e_l}$ and $n_r/2^{e_r}$, the
counts at its ends and the signs of $f$ just inside them: an interval with equal counts, so without
a root, and the sign $-1$ is dropped; one with equal counts and the sign $+1$, or at the depth of
the step, is kept as the interval of the doubles \verb|ldexp((double)n, -e)| of its ends, appended
to the array \verb|ivlocal| of $1 + \lfloor d/2 \rfloor$ intervals on the stack, or extending the
last interval kept when that ends at its lower end (\verb|==| of doubles); and any other is split at
its midpoint, $0$ for $[-\infty, +\infty]$, $-1$ or $2 n_r$ for $[-\infty, n_r]$, $1$ or $2 n_l$ for
$[n_l, +\infty]$, and otherwise the half of the sum of the ends over the larger denominator
(\csource{sturm.c}{234}{251}), where the signs of $f$ and $f'$ and the sign changes of the chain are
evaluated, the zeros skipped (\csource{sturm.c}{252}{260}, with \verb|eval_sign|, Horner's rule in
integers, \csource{sturm.c}{67}{84}), and both halves are bisected, the left one first, the count at
the midpoint and the sign of $f$ there handed to both; at a root $m$ of $f$, which is simple, the
left half gets the count plus $1$ and $-\mathrm{sign}\, f'(m)$ and the right half the count and
$\mathrm{sign}\, f'(m)$, the count and the sign just left and just right of $m$. A finite end at the
depth $\delta$ is $n/2^e$ with $|n| \le 2^{\delta - 2}$ (the comment at lines 209-216), so that up
to the depth $55$ every end is a double exactly, and $\mp$\verb|LONG_MAX| become $\mp 2^{63}$. The
first call starts from the counts at $\mp\infty$ and the sign of $c_d$, negated at $-\infty$ for odd
$d$ (\csource{sturm.c}{269}{271}). Last, it intersects the intervals kept with the search intervals,
from a copy of these on the heap, with two cursors and four cases (\csource{sturm.c}{272}{325}): the
intervals noted may be degenerate, and their ends are copied, so that a $-0.0$ of the caller can
stay. It returns $-1$ when the chain finds a repeated factor, $0$ when $f$ is negative everywhere or
the intersection is empty, and otherwise the number of the intervals, which it writes into
\verb|domain| and \verb|num_inter|; there are at most $\mathrm{num\_inter} + d$ of them, the room
that the manual asks for.

\paragraph{The returns.} With $\mathrm{sturm} \ge 0$, \verb|find_points_work_1| calls the step
(\csource{find_points.c}{4857}{4892}): $-1$ gives \verb|RATPOINTS_NON_SQUAREFREE|
(\csource{find_points.c}{4879}{4879}); $0$ returns the count $0$ at once for even degree with $c_d$
not a square (\csource{find_points.c}{4889}{4889}), when there is no point at infinity either, and
otherwise sets the flag \verb|sturm_empty| (\csource{find_points.c}{4890}{4890}), by which the
search returns after the flags, the exits and the points at infinity, before the sieve
(\csource{find_points.c}{5106}{5106}). At the return of line 4889 the flags of the squares are clear
and \verb|RATPOINTS_REVERSED| is in fact clear too: the cases 1 and 3 of the reversal leave an odd
degree, and in case~2 the reversed polynomial takes the value $c_d > 0$ at $0$, in the default
interval, which the Sturm step keeps.

\paragraph{The bounds of the numerators.} Once for the run, \verb|run_setup| makes the bounds of the
search intervals (\csource{find_points.c}{3866}{3870}; the struct \verb|rp_intervals| of lines
3435-3440): for each end $x$, \verb|exact_end| (\csource{find_points.c}{3264}{3299}) writes
$|x| = m\,2^{-s}$ exactly, $m = 2^{63}$ and $s = 0$ for $|x| \ge 2^{63}$ and the infinities, which
lie beyond every numerator, and otherwise $m < 2^{53}$ from \verb|frexp| and \verb|ldexp|, its
trailing zero bits dropped while $s > 0$, $m$ shifted up for $s < 0$ and $s$ capped at $127$, with
the sign $x < 0$ (so that $-0.0$ is not negative) and the rounding $r = 2^s - 1$ when the end wants
$b|x|$ rounded up, at a lower end with $x \ge 0$ and an upper end with $x < 0$, and $r = 0$
otherwise; \verb|exact_scale| (\csource{find_points.c}{3231}{3257}) computes
$\min(2^{64} - 1, \lfloor (b m + r)/2^s \rfloor)$ in 128 bits, with \verb|unsigned __int128| in the
default build and from 32-bit halves in a build without that type or with \verb|RP_NO_INT128|
defined; \verb|end_form| (\csource{find_points.c}{3304}{3318}) gives an end the form
$w = (b M + C) \gg \mathit{sh}$ that is exact in a \verb|long| for every $b \le b_{\mathrm{high}}$,
$(0, \mp H, 0)$ for an end beyond the height bound on its side and
$(\pm m, 2^s - 1 \text{ or } 0, s)$ for a short mantissa; and \verb|exact_bounds|
(\csource{find_points.c}{3323}{3340}) takes an interval inline when $H \le 2^{63} - 65$ and both
ends have a form, and gives it the forms of an empty interval otherwise. For a denominator $b$,
\verb|interval_bits| (\csource{find_points.c}{3410}{3434}) computes $\lceil b l \rceil$ and
$\lfloor b u \rfloor$ from the forms: the code $2$ when the first is above $H$, so that no later
interval holds a numerator either; the code $1$ when the second is below $-H$ for an interval taken
inline, and otherwise the products in 128 bits of \verb|interval_bits_exact|
(\csource{find_points.c}{3381}{3406}), which splits the ends for the heights above $2^{63} - 65$;
and else the code $0$ with the bits of the numerators clipped to $[-H, H]$ in the packing of the
class (\verb|numerators_to_bits|, Definition~\ref{def:numeratorsToBits}). Since the search intervals
are increasing and strictly apart and the bounds exact,
$\lfloor b u_i \rfloor < \lceil b l_{i+1} \rceil$: the ranges of the numerators increase and are
disjoint.

\paragraph{The sieve of a denominator.} \verb|sift| runs over the bounds of the intervals in turn
(lines 3471-3613): the code $2$ returns, the code $1$ goes on to the next interval, and at the first
interval with the code $0$, also one that holds no numerator, it sets up the denominator, the rows
of the moduli and the specs of the sieve, building the rows that are missing (the flag
\verb|do_setup|, set at line 3458 and tested and cleared at lines 3490-3494; the set-up lines
3490-3561); then it sieves the bits of the interval in pieces. So a denominator none of whose
intervals has $\lceil b l \rceil \le H$ and $\lfloor b u \rfloor \ge -H$ (each gives the code $1$,
or the code $2$ returns first, after which no later interval would give anything else) neither reads
nor builds its rows. One with such an interval sets them up even when that interval holds no
numerator, or none of the class, and within a denominator the numerators come in increasing order.

\paragraph{The range of the denominators.} The model has carried the range since stages~E and~F,
with $b_{\mathrm{low}} = 1$ and $b_{\mathrm{high}} = H$: the divisor set-up keeps the divisors up to
$b_{\mathrm{high}}$ (lines 875-876), whose number decides the shape of the divisors, so that a
smaller $b_{\mathrm{high}}$ can change the order; the Jacobi tables need
$b_{\mathrm{high}} \le$ \verb|RP_MULMOD_LIMIT| (line 688); the forbidden divisors are the prime
powers up to $b_{\mathrm{high}}$ (lines 2217, 2230, 2637, 2639, 2735); and the shapes list the
denominators from $b_{\mathrm{low}}$ to $b_{\mathrm{high}}$ (lines 3816-3830). The estimate of the
work in the choice of the moduli reads the range and the intervals in floating-point arithmetic
(\verb|numerators_for| and \verb|run_shape|, lines 1836-2106), which is stage~I's.

\paragraph{The doubles.} Apart from that estimate, the C code does no floating-point arithmetic on
the ends: it converts \verb|long|s to doubles, which rounds to nearest with ties to even (the
default interval, the ends of the bisection), compares doubles, steps to the next double
(\verb|nextafter|), takes a double apart and scales it by a power of $2$ (\verb|frexp|,
\verb|ldexp|), converts doubles to \verb|long|s and \verb|unsigned long|s, and prints them. glibc's
\verb|strtod| rounds the exact value of its word correctly, to nearest with ties to even, also at
the boundaries of overflow and underflow, and its \verb|printf| prints under \verb|%f| (and
\verb|%lf|) the exact value of the double rounded to six places, ties to even, with a \verb|-|
whenever the sign bit is set.

\paragraph{The program.} \verb|main.c| keeps the intervals of the command line in two arrays of
$2\,\mathrm{maxDegree}$ entries, the doubles that it gives the library and the exact ends
(\csource{main.c}{66}{73}, the type \verb|exact_end| with $\mathrm{inf} = -1$, $0$ or $1$ and a GMP
rational $q$), with the defaults \verb|num_inter| $= 0$, \verb|b_low| $= 1$ and \verb|b_high| $= H$
(\csource{main.c}{333}{336}, \csource{main.c}{380}{380}). \verb|-l w| (\csource{main.c}{398}{413})
fails with error~7 when the word is missing, after another \verb|-l|, when $100$ intervals have been
read, when \verb|scan_end| refuses the word, and when the end is not above the upper end of the
interval before, compared exactly (\verb|cmp_end|, \csource{main.c}{309}{313}, two infinities of one
sign equal); \verb|-u w| (\csource{main.c}{414}{436}) without a \verb|-l| before it starts the first
interval at $-\infty$, with the double $-h$ of \verb|height_end| (\csource{main.c}{299}{305}), and
fails for a later interval; it makes the lower double of such a first interval $-\infty$ when the
upper double lies below $-h$, and fails when the lower end is above the upper end; after the options
a pending lower end gets the upper end $+\infty$, with the double $h$, or $+\infty$ when $h$ lies
below the lower double (\csource{main.c}{669}{679}). \verb|-dl| and \verb|-du| take a C \verb|long|,
error~6 otherwise, the last one winning (\csource{main.c}{650}{666}). \verb|scan_end|
(\csource{main.c}{798}{880}) reads the word as \verb|scan_double| does, with \verb|strtod|
(\csource{main.c}{746}{752}): an infinite double, also one of a number beyond the range of the
doubles, gives the infinite end of its sign; otherwise the end is the exact rational number of the
digits of the word, $m\,10^{e - k}$ or, for a hexadecimal word, $m\,2^{e - 4k}$, with its exponent
accumulated up to \verb|LONG_MAX / 20|, and $\pm 2^{-1075}$ when \verb|strtod| reads $0$ for a word
whose digits are not all $0$; and the double is that of \verb|outward| (\csource{main.c}{761}{776}),
which steps from the double of \verb|strtod| by \verb|nextafter| until it lies at or below the end
(a lower end) or at or above it (an upper end), stopping at an infinity: the largest double at most
the end, or the smallest at least it. Before the search \verb|join_intervals|
(\csource{main.c}{891}{905}, called at line 248) joins intervals whose doubles meet or overlap, so
that the library never refuses them; the data of \verb|process| keep the number of the intervals
before the join and the exact ends (\csource{main.c}{246}{247}), and \verb|process| passes an affine
point over, with the weight $0$ and the flag of the stop clear, when it lies in none of the exact
intervals (\csource{main.c}{150}{152}, with \verb|in_region|, a binary search over the exact ends,
\csource{main.c}{116}{137}): the library may find a point just beyond an end of the command line,
since its intervals are outward of them, and such a point is neither counted nor stops \verb|-1|.
Without \verb|-q|, and only when the command line gives search intervals (\verb|-l| or \verb|-u|, so
that \verb|num_inter| $> 0$; \verb|-dl| and \verb|-du| alone give none), message~3 prints the search
region of the command line, the doubles before the join, as \verb|[%f, %f]| joined by \verb| U |;
unlike message~8 it does not depend on \verb|-v|, and it comes after the height bound and its blank
line, or under \verb|-v|, which prints no height bound, directly after the curve equation
(\csource{main.c}{224}{224}, \csource{main.c}{984}{991}). Message~8 prints, after message~7 and not
under \verb|-v|, the search intervals as the library leaves them, under \verb|%lf|, which is
\verb|%f| (\csource{main.c}{1022}{1030}): $[-h, h]$ narrowed by the Sturm step without \verb|-l| and
\verb|-u|, in the coordinates of the polynomial searched, $\mp 2^{63}$ for the infinite ends of the
step, and no interval after a return of the step. Error~7 has its message at
\csource{main.c}{1044}{1046}; error~9, which \verb|main.c| gives for a code of the library other
than \verb|RATPOINTS_NON_SQUAREFREE| (error~8) and \verb|RATPOINTS_NO_MEMORY| (error~10, which the
model, having no failure for lack of memory, never gives), that is for \verb|RATPOINTS_BAD_ARGS| and
\verb|RATPOINTS_WORK_LENGTH_TOO_SMALL| (\csource{main.c}{256}{261}), at
\csource{main.c}{1048}{1048}.

\paragraph{Why the points do not change.} Each row of the chain of the C code is a positive rational
multiple of the row of the Sturm sequence of $f$ and $f'$, so that the sign changes of the chain are
those of the Sturm sequence at every point; by Sturm's theorem a dropped interval holds no root of
$f$ and has $f < 0$ just inside its left end, so $f < 0$ on it, and every $x$ with $f(x) \ge 0$ lies
in an interval kept, as long as its ends are exact, which they are up to the depth $55$. So the
intersection with the search intervals keeps every $x$ of the search intervals with $f(x) \ge 0$ and
$|x| \le 2^{63}$, and the Sturm step changes no point: the points found are those of height at most
$H$ in the caller's search region (Definition~\ref{def:regionSpec}): the points at infinity,
whatever the region, and the affine points whose $z$ lies in the range of the caller and whose $x/z$
lies in one of the caller's intervals, or anywhere when there are none. The survivors under
\verb|RATPOINTS_NO_CHECK| are those of the sieve over the intervals that the step keeps, and so
depend on its depth.

\section{The specification of the search region}\label{sec:regionSpec}

The specification gains two files, which use Mathlib and nothing of the model and which the
challenge restates verbatim. \verb|Ratpoints/Spec/Region.lean| speaks about the doubles by their
bits, $64$-bit words, and about their values in $\QQ$ extended by $\bot$ and $\top$, Mathlib's
\verb|WithBotTop ℚ|: the bits serve the value of a double and the test of a NaN in well-formedness,
and every other predicate reads the values; it imports \verb|Mathlib.Order.WithBotTop| and
\verb|Mathlib.Algebra.Order.Field.Rat|. \verb|Ratpoints/Spec/Sturm.lean| imports the Sturm sequence,
the sign changes of a list, the derivative, squarefreeness, the real numbers and the cardinality of
sets (\verb|Mathlib.Algebra.Polynomial.Sturm.Sequence|, \verb|Mathlib.Data.List.SignVariations|,
\verb|Mathlib.Algebra.Polynomial.Derivative|, \verb|Mathlib.Algebra.Squarefree.Basic|,
\verb|Mathlib.Basic.Real.Basic| and \verb|Mathlib.Data.Set.Card|).

\begin{definition}[The search region]\label{def:regionSpec}
  Proposed names: \lname{Ratpoints.IsNaN}, \lname{Ratpoints.doubleValue},
  \lname{Ratpoints.AreSearchIntervals}, \lname{Ratpoints.IntervalsWF}, \lname{Ratpoints.InRegion},
  \lname{Ratpoints.IsFullBox}, \lname{Ratpoints.denominatorBound},
  \lname{Ratpoints.instDecidableIsNaN}, \lname{Ratpoints.instDecidableAreSearchIntervals},
  \lname{Ratpoints.instDecidableIntervalsWF}, \lname{Ratpoints.instDecidablePredInRegion},
  \lname{Ratpoints.instDecidableIsFullBox}. For the bits $x < 2^{64}$ of an IEEE binary64 double,
  with the exponent field $e = \lfloor x / 2^{52} \rfloor \bmod 2^{11}$, the fraction
  $f = x \bmod 2^{52}$ and the sign $\sigma = -1$ for $x \ge 2^{63}$ and $\sigma = 1$ otherwise:
  $x$ is a NaN when $e = 2047$ and $f \ne 0$, and the value of the double is
  \[
    \mathrm{doubleValue}(x) =
    \begin{cases}
      \sigma f\, 2^{-1074} & \text{if } e = 0 \text{ (the zeros and the subnormals)}, \\
      \sigma (2^{52} + f)\, 2^{e - 1075} & \text{if } 1 \le e \le 2046, \\
      \bot \text{ for } \sigma = -1, \ \top \text{ for } \sigma = 1 & \text{if } e = 2047,
    \end{cases}
  \]
  so that both zeros have the value $0$ and the value of a NaN is junk, an infinity, and nothing
  depends on it: a well-formed list has no NaN end, and the statements that read the values of the
  ends (statements 3, 8, 9 and~10 of Section~\ref{sec:certifiedH}) assume a successful search, which
  by statements 5 and~6 refuses search intervals that are not well formed. The intervals of a list
  over the values \emph{are search intervals} when each $[l, u]$ has $l \le u$ and each upper end
  lies below the lower end of the next interval; a list of pairs of bits is \emph{well formed}, the
  test of \verb|RATPOINTS_BAD_ARGS|, when no end is a NaN and the values of the ends are search
  intervals. A triple $(x, y, z)$ lies \emph{in the search region} of the intervals $D$ (over the
  values) and the range from $b_{\mathrm{low}}$ to $b_{\mathrm{high}}$ when $z = 0$, or
  $b_{\mathrm{low}} \le z$, $z \le b_{\mathrm{high}}$ unless $b_{\mathrm{high}} < 1$, and $x/z$ lies
  in one of the intervals, or anywhere when $D$ is empty:
  \[
    z = 0 \ \lor\ \bigl(b_{\mathrm{low}} \le z \ \land\ (b_{\mathrm{high}} < 1 \lor
      z \le b_{\mathrm{high}}) \ \land\ (D = [\,] \lor
      \exists [l, u] \in D,\ l \le x/z \le u)\bigr).
  \]
  The search region is the \emph{whole box} of the height bound $H$ when there are no intervals,
  $b_{\mathrm{low}} \le 1$, and $b_{\mathrm{high}} < 1$ or $H \le b_{\mathrm{high}}$; only then may
  ratpoints reverse the polynomial. The \emph{largest denominator} of the search is
  $b_{\mathrm{high}}$ when $1 \le b_{\mathrm{high}} \le H$, and $H$ otherwise. Each predicate has
  its decidability instance, which the challenge restates with it.
\end{definition}

The file checks its definitions by \verb|decide +kernel|, which adds no axiom: the values of $1$,
$-0.0$, $\pm\infty$, the least subnormal, $2^{63}$ and $-1.5$; a NaN end refused, $[-0.0, 0.0]$ and
$[0.0, -0.0]$ accepted, $[-1, -0.0]$ followed by $[0.0, 1]$ refused, since the two are not strictly
apart, and $[-\infty, -\infty]$ accepted, as ratpoints accepts and refuses them; and a few points in
and outside regions. The region of the specification reads the values of the caller's doubles, not
their bits, so that $-0.0$ and $0.0$ give the same region.

\begin{definition}[Sturm's theorem]\label{def:sturmTheorem}
  Proposed name: \lname{Ratpoints.SturmTheorem}. The proposition that for every squarefree real
  polynomial $p$ and all reals $a \le b$
  \[
    \#\{x \in \mathbb{R} : a < x \le b,\ p(x) = 0\} + V_p(b) = V_p(a),
  \]
  where $V_p(t)$ is the number of sign changes (\verb|List.signVariations|) of the values at $t$ of
  the Sturm sequence of $p$ and its derivative (\verb|Polynomial.sturmSeq p (derivative p)|), the
  number of roots being \verb|Set.ncard|. It is the textbook theorem for an interval $(a, b]$, a
  hypothesis of the theorems that need it, not an axiom.
\end{definition}

The form with $(a, b]$ and ends that may be roots is the one that the bisection needs, which
evaluates the chain at midpoints that may be roots of $f$; the ends at $\pm\infty$ come from the
asymptotics of polynomials, proved in the project (Section~\ref{sec:coverage}). The points in the
order of ratpoints (Definition~\ref{def:orderedPoints}) gain the largest denominator $B$ of the
search: the denominators of the search of the specification are taken up to $B$, while the points
keep the height bound $H$; $B$ is the height bound, or the bound of the range of the denominators
($\mathrm{denominatorBound}$), and the other parts of a search region, the intervals and the least
denominator, only filter the list, while a search region other than the whole box forbids the
reversal. The file of the order imports the file of the region, for four examples with a search
region that are checked by \verb|decide +kernel| against the output of the C program:
\verb|-6 5 -3 30| at the height $12$ with \verb|-dl 5 -du 6|, which keeps the divisor order, $6$
before $5$; \verb|-1 0 9| at $10$ with \verb|-du 9|, which forbids the reversal, and with
\verb|-dl 0 -du 11|, the whole box, which allows it; and \verb|-1 0 9| at $10$ with \verb|-du 9| and
the interval that the program passes the library for \verb|-l -0.2 -u 0.6|, as \verb|outward| gives
its ends: the double nearest to $-0.2$, which lies below it, and the smallest double above $0.6$
(the double nearest to $0.6$ lies below it); the program's filter by the exact ends drops no point
here.

\section{The doubles}\label{sec:doubles}

\begin{definition}[The doubles of the C code]\label{def:Double}
  Proposed names: \lname{Ratpoints.Double}, \lname{Ratpoints.Double.mantDig},
  \lname{Ratpoints.Double.zero}, \lname{Ratpoints.Double.huge}, \lname{Ratpoints.Double.twoPow63},
  \lname{Ratpoints.Double.signBit}, \lname{Ratpoints.Double.expField},
  \lname{Ratpoints.Double.fracField}, \lname{Ratpoints.Double.isNaN},
  \lname{Ratpoints.Double.isInf}, \lname{Ratpoints.Double.isZero},
  \lname{Ratpoints.Double.isFinite}, \lname{Ratpoints.Double.orderValue},
  \lname{Ratpoints.Double.lt}, \lname{Ratpoints.Double.le}, \lname{Ratpoints.Double.beq},
  \lname{Ratpoints.Double.neg}, \lname{Ratpoints.Double.fabs}, \lname{Ratpoints.Double.toDyadic},
  \lname{Ratpoints.Double.toRat}, \lname{Ratpoints.Double.roundShift},
  \lname{Ratpoints.Double.ofDyadic}, \lname{Ratpoints.Double.ofRatio},
  \lname{Ratpoints.Double.ofInt64}, \lname{Ratpoints.Double.toInt64},
  \lname{Ratpoints.Double.toUInt64}, \lname{Ratpoints.Double.ldexp}, \lname{Ratpoints.Double.frexp},
  \lname{Ratpoints.Double.nextUp}, \lname{Ratpoints.Double.nextDown}. A C \verb|double| by its $64$
  bits, a structure with no equality of its own, so that \verb|==| never means equality of the bits:
  the C language's \verb|==| is $\mathrm{beq}$. Its fields (the sign bit, the exponent field, the
  fraction) and its kinds, as \verb|isnan|, \verb|isinf| and \verb|isfinite| of \verb|math.h| decide
  them, and the zeros; \verb|DBL_MANT_DIG| $= 53$, a fact of the format, not a parameter of
  ratpoints; $0.0$, \verb|HUGE_VAL| $= +\infty$ and the literal $2^{63}$ of the C code. The
  comparisons \verb|<|, \verb|<=| and \verb|==| of the C language, false when either side is a NaN
  and otherwise the order of the values, by a key of each double that is not a NaN, the bits without
  the sign bit, negated when it is set, which increases with the value and is $0$ for both zeros;
  the negation and \verb|fabs|, which flip and clear the sign bit, also of a zero, an infinity and a
  NaN. The exact value of a finite double, $m\,2^k$ as a pair of integers and as a rational number,
  as \verb|mpq_set_d| of GMP gives it; the double nearest to $\pm m\,2^k$ and to $\pm n/d$, with
  ties to even, as IEEE~754 rounds an exact result to nearest, with a signed zero, the subnormal
  range and an infinity beyond the largest finite double; the conversion of a \verb|long| to a
  double, the double nearest to it, exact for $|n| \le 2^{53}$ and $\pm 2^{63}$ for
  $\pm(2^{63} - 1)$, and of a finite double to a \verb|long| and to an \verb|unsigned long|,
  truncated, for the doubles for which the C language defines the conversion (for the others the
  value is of no meaning); \verb|ldexp| and \verb|frexp|, as glibc computes them, exact in the range
  that ratpoints uses; and \verb|nextafter| towards $+\infty$ and towards $-\infty$, the neighbours
  of a double, with the least subnormals next to the zeros. These model the C language and the C
  library on the platforms of ratpoints, IEEE~754 binary64 with rounding to nearest, as gcc and
  glibc implement them, and hold no C range.
\end{definition}

\begin{lemma}[The doubles and their values]\label{lem:Double}
  \uses{def:Double, def:regionSpec}
  \lean{Ratpoints.Double.isNaN_iff, Ratpoints.Double.maxOrderValue, Ratpoints.Double.magnitudeUnits, Ratpoints.Double.orderValueUnits, Ratpoints.Double.valueOfOrderValue, Ratpoints.Double.strictMonoOn_valueOfOrderValue, Ratpoints.Double.doubleValue_eq_valueOfOrderValue, Ratpoints.Double.lt_eq_true_iff, Ratpoints.Double.le_eq_true_iff, Ratpoints.Double.beq_eq_true_iff, Ratpoints.Double.lt_eq_true_iff', Ratpoints.Double.le_eq_true_iff', Ratpoints.Double.doubleValue_eq_toRat, Ratpoints.Double.toRat_eq, Ratpoints.Double.toRat_neg, Ratpoints.Double.toRat_fabs, Ratpoints.Double.ofDyadic_exact, Ratpoints.Double.spec_ofDyadic, Ratpoints.Double.spec_ofSignificand, Ratpoints.Double.ofInt64_exact, Ratpoints.Double.ofInt64_maxValue, Ratpoints.Double.spec_ofInt64_natCast, Ratpoints.Double.ldexp_exact, Ratpoints.Double.frexp_spec, Ratpoints.Double.toInt_toInt64, Ratpoints.Double.toNat_toUInt64, Ratpoints.Double.orderValue_nextUp, Ratpoints.Double.nextUp_le_of_lt, Ratpoints.Double.lt_nextUp, Ratpoints.Double.orderValue_nextDown, Ratpoints.Double.le_nextDown_of_lt, Ratpoints.Double.nextDown_lt, Ratpoints.Double.toRat_nextUp_of_normal, Ratpoints.Double.le_doubleValue_or_of_lt_eq_false, Ratpoints.Double.isFinite_and_abs_lt_of_lt, Ratpoints.Double.lt_zero_eq_true_iff_toRat, Ratpoints.Double.lt_twoPow63_eq_true_iff_toRat}
  \leanok
  The test of the model is the NaN of the specification, and the value of a double that is not a NaN
  is a strictly increasing function of its key on the range of the keys, so that the comparisons of
  the C language are the order and the equality of the values (Definition~\ref{def:regionSpec}),
  with or without the hypotheses that the two sides are not NaNs; the value of a finite double is
  its exact value, and the negation and \verb|fabs| negate it and take its absolute value. The
  double nearest to $\pm m\,2^k$ is exact when $0 < m < 2^{53}$, $k \ge -1074$ and
  $m\,2^k < 2^{1024}$, and in the normal range it is $\pm M 2^q$ with $2^{52} \le M < 2^{53}$ and
  $m\,2^k < (M + 1) 2^q$; the conversion of a \verb|long| is exact for $|n| \le 2^{53}$ and gives
  $\pm 2^{63}$ for $\pm(2^{63} - 1)$, and the next double above the conversion of $1 \le H < 2^{63}$
  lies above $H$; \verb|ldexp| scales the value exactly while it stays in the normal range,
  \verb|frexp| gives $x = g\,2^{\mathit{ex}}$ with $1/2 \le |g| < 1$, and the conversions to
  \verb|long| and \verb|unsigned long| are the floor of a value in their range. The neighbours of a
  double that is not a NaN and not at the end of the order have the next keys, so that no double
  lies strictly between a double and its neighbour; the step above a positive normal double of
  exponent field $e$ (Definition~\ref{def:regionSpec}) is $2^{e - 1075}$. And the facts on the ends
  that the bounds and the default interval read: a comparison that fails puts the values in the
  other order, a double below $2^{63}$ in absolute value is finite, and the tests $x < 0$ and
  $x < 2^{63}$ are those of the exact value.
\end{lemma}

\begin{proof}\leanok
  From the fields: the key of a double is its number of steps from $0$ in the order of the values,
  the zeros and the subnormals first, and the value of a key is a sum of steps that grow with the
  exponent; the rounding by $\mathrm{roundShift}$ and the carry at $2^{53}$ have their bounds; the
  rest by cases on the fields and by \verb|decide| on the constants.
\end{proof}

\begin{lemma}[The roundings are faithful]\label{lem:rounding}
  \uses{def:Double, lem:Double}
  Proposed names: \lname{Ratpoints.Double.lt\_nextUp\_and\_nextDown\_lt\_ofRatio},
  \lname{Ratpoints.Double.lt\_nextUp\_and\_nextDown\_lt\_ofDyadic},
  \lname{Ratpoints.Double.isNaN\_ofRatio}, \lname{Ratpoints.Double.isNaN\_ofDyadic},
  \lname{Ratpoints.Double.ofRatio\_eq}, \lname{Ratpoints.Double.ofDyadic\_eq},
  \lname{Ratpoints.Double.ofRatio\_rounding}, \lname{Ratpoints.Double.ofDyadic\_rounding},
  \lname{Ratpoints.Double.assemble}, \lname{Ratpoints.Double.ofSignificand},
  \lname{Ratpoints.Double.roundDiv}, \lname{Ratpoints.Double.ratioExponent},
  \lname{Ratpoints.Double.ratioLast}, \lname{Ratpoints.Double.ratioPair},
  \lname{Ratpoints.Double.lt\_nextUp\_and\_nextDown\_lt},
  \lname{Ratpoints.Double.le\_of\_lt\_nextUp},
  \lname{Ratpoints.Double.le\_nextDown\_of\_doubleValue\_lt},
  \lname{Ratpoints.Double.nextUp\_isFinite\_or}, \lname{Ratpoints.Double.nextDown\_isFinite\_or},
  \lname{Ratpoints.Double.nextUp\_of\_beq\_zero}, \lname{Ratpoints.Double.doubleValue\_one}. The
  double nearest to $\pm n/d$, for $d \ne 0$, and the double nearest to $\pm m\,2^k$ are never a
  NaN. For $n \ne 0$ and $m \ne 0$, when the double is finite, it is a faithful rounding of the
  value: the value lies strictly between the two neighbours of the double, which are finite or
  infinite. A double whose next double lies above a value $v$ is at least every double at most $v$,
  and dually.
\end{lemma}

\begin{proof}
  The last steps of both roundings, the exponent, the quotient rounded to the nearest integer and
  the assembly of the bits, are restated in the proofs with an equation by \verb|rfl|; the rounded
  quotient lies within one unit of the exact quotient, and the neighbours of the result are one unit
  of the last place away, also across a change of the exponent.
\end{proof}

The model computes the nearest double with ties to even, as glibc's \verb|strtod| does; the proofs
show only that the rounding is faithful, which is all that the theorems need, and the guards and the
differential tests check the nearest double against glibc (Section~\ref{sec:regionTests}).

\section{The options and the search intervals}\label{sec:regionOptions}

The options of the search (Definition~\ref{def:Options}) gain the search intervals, a list of pairs
of doubles, empty by default, for the array \verb|domain| with its number \verb|num_inter|, and the
least and the largest denominator, integers like the \verb|long|s \verb|b_low| and \verb|b_high|,
with the defaults $1$ and $0$, the latter meaning $H$, as for the library; the depth of the Sturm
step keeps its type. The lines of \verb|ratpoints.h| that hold these fields hold fields of the
parameters and of the choice of the moduli too, and the options own them with a deviation that names
the others: \verb|sp1| and \verb|sp2| of the parameters, \verb|sp3| of stage~I, and
\verb|num_primes| and \verb|max_forbidden| of the options of the program. The proofs read the
intervals through their values and their bits.

\begin{definition}[The values and the bits of the intervals]\label{def:inIntervals}
  \uses{def:Double, def:regionSpec}
  \lean{Ratpoints.domainValues, Ratpoints.domainBits, Ratpoints.InIntervals, Ratpoints.instDecidablePredInIntervals}
  \leanok
  The intervals of a list of pairs of doubles over the values of their ends
  (Definition~\ref{def:regionSpec}) and over their bits; and the predicate that a rational number
  lies in one of the intervals over the values. The first two are abbreviations, so that the
  statements of the proofs read $\mathrm{doubleValue}$ and the bits of the doubles as the challenge
  does.
\end{definition}

\begin{definition}[The normalized options]\label{def:normalize}
  \uses{def:Options, def:Double}
  Proposed name: \lname{Ratpoints.Options.normalize}. The options as \verb|find_points_work_1|
  normalizes the fields that hold them, for the height bound $H$: the least denominator $1$ for one
  below $1$, the largest denominator $H$ for one below $1$ or above $H$, the depth of the Sturm step
  capped at $\mathrm{mantDig} + 2 = 55$, and the switch of the reversal cleared, for the bit
  \verb|RATPOINTS_NO_REVERSE_AUTO|, when the caller gives search intervals or the range of the
  denominators is not $1, \dots, H$ (\csource{find_points.c}{4640}{4641},
  \csource{find_points.c}{4659}{4661}, \csource{find_points.c}{4678}{4689},
  \csource{ratpoints.h}{421}{424}). The reversal check reads the normalized switch, which folds the
  two bits of the flags that the C code tests (line 4739) into one, and the search the normalized
  range and depth.
\end{definition}

\begin{lemma}[The normalized options]\label{lem:normalize}
  \uses{def:normalize, def:regionSpec, def:inIntervals}
  Proposed names: \lname{Ratpoints.Options.reverse\_normalize},
  \lname{Ratpoints.Options.domain\_normalize}, \lname{Ratpoints.Options.noCheck\_normalize},
  \lname{Ratpoints.Options.noY\_normalize}, \lname{Ratpoints.Options.toNat\_bHigh\_normalize},
  \lname{Ratpoints.Options.one\_le\_toNat\_bLow\_normalize},
  \lname{Ratpoints.Options.one\_le\_toNat\_bHigh\_normalize},
  \lname{Ratpoints.Options.toNat\_bHigh\_normalize\_le},
  \lname{Ratpoints.Options.lt\_one\_or\_le\_bHigh\_of\_le},
  \lname{Ratpoints.Options.le\_toNat\_bHigh\_normalize\_of\_le},
  \lname{Ratpoints.Options.bLow\_normalize\_le\_iff},
  \lname{Ratpoints.Options.sturm\_normalize\_nonneg\_iff},
  \lname{Ratpoints.Options.sturm\_normalize\_neg\_iff},
  \lname{Ratpoints.Options.sturm\_normalize\_le}, \lname{Ratpoints.inRegion\_iff},
  \lname{Ratpoints.inRegion\_of\_isFullBox}. The normalized switch of the reversal is the switch and
  the whole box (Definition~\ref{def:regionSpec}); the normalization keeps the intervals and the
  options of the output; the normalized largest denominator is the largest denominator of the
  search, from $1$ to $H$ for $H \ge 1$, and the normalized least one is at least $1$; a
  denominator $1 \le z \le H$ lies in the normalized range exactly when it lies in the range of the
  caller as the search region reads it; the normalized depth is at most $55$ and has the sign of the
  depth. The search region of intervals over the values is that of $\mathrm{InIntervals}$, and every
  triple with $0 \le z \le H$ lies in the whole box.
\end{lemma}

\begin{proof}
  From the definitions, by cases on the clamps.
\end{proof}

\begin{definition}[The search intervals of a search]\label{def:searchIntervals}
  \uses{def:Double, def:reversal}
  \lean{Ratpoints.searchIntervals}
  Proposed name: \lname{Ratpoints.Refusal.badArgs}. For the height bound $H \ge 1$ and the
  intervals of the caller: for none, the default interval $[-h, h]$, where $h$ is the double of $H$,
  moved up to the next double when $h < 2^{63}$ and $h < H$, so that it holds every $x = a/b$ of the
  search (\csource{find_points.c}{4829}{4829}, \csource{find_points.c}{4831}{4841}); otherwise the
  intervals given, or none, which the search turns into the refusal $\mathrm{badArgs}$, the code
  \verb|RATPOINTS_BAD_ARGS|, when an interval has not $l \le u$ or does not begin above the end of
  the one before it, compared as doubles, so that a NaN end is refused too (lines 4842-4855; the
  test of a null pointer, line 4830, is excluded as the interface of the library). The program turns
  the refusal into its error~9, which it gives for a code of the library other than
  \verb|RATPOINTS_NON_SQUAREFREE| and \verb|RATPOINTS_NO_MEMORY|. Until H.5 the model has the
  default interval only, which the sieving takes for a search without intervals, and returns the
  intervals given as they are, never none.
\end{definition}

\begin{lemma}[The search intervals of a search]\label{lem:searchIntervals}
  \uses{def:searchIntervals, lem:Double, def:regionSpec, def:inIntervals}
  Proposed names: \lname{Ratpoints.run\_searchIntervals\_eq},
  \lname{Ratpoints.intervalsWF\_domainBits\_iff},
  \lname{Ratpoints.run\_searchIntervals\_eq\_none\_iff},
  \lname{Ratpoints.run\_searchIntervals\_of\_ne\_nil},
  \lname{Ratpoints.intervalsWF\_of\_run\_searchIntervals\_eq\_some},
  \lname{Ratpoints.wf\_of\_run\_searchIntervals\_eq\_some},
  \lname{Ratpoints.le\_toRat\_defaultBound}, \lname{Ratpoints.intervalsWF\_and\_inIntervals\_neg},
  \lname{Ratpoints.exists\_run\_searchIntervals\_nil}. The loop of the validation is a test of each
  index; well-formedness (Definition~\ref{def:regionSpec}) is that test in the comparisons of the C
  language; so the search intervals are none exactly when the caller's intervals are not well
  formed, they are the caller's intervals when these are well formed and not empty, and the result
  is well formed for $1 \le H < 2^{63}$. The end $h$ of the default interval is finite and at least
  $H$, so that $[-h, h]$ is well formed and holds $a/b$ for every $b > 0$ and $|a| \le H$.
\end{lemma}

\begin{proof}
  Lemma~\ref{lem:Double} turns the comparisons into the order of the values, and the chain of the
  specification is read index by index (\verb|List.isChain_iff_getElem|), each index being the test
  of the loop; the end of the default interval by the conversion of $H$ and its next double.
\end{proof}

\section{The exact bounds of the numerators}\label{sec:bounds}

The bounds of the search intervals are made once for the run, in the new model file
\verb|Ratpoints/Model/Bounds.lean|, and the two functions that turn the bounds of an interval into
the bits of the numerators of a class sit in \verb|Ratpoints/Model/Sieve.lean|, next to the
conversion of numerators to bits that they end with (Definition~\ref{def:numeratorsToBits}). They
compute on the C code's words: \verb|UInt64| and \verb|Int64| for its \verb|unsigned long|s and
\verb|long|s, with the range of each \verb|long| stated as a lemma, as for the packing.

\begin{definition}[An end of a search interval]\label{def:exactEnd}
  \uses{def:Double, def:ctz}
  \lean{Ratpoints.Sieve.ExactEnd, Ratpoints.Sieve.exactMantissa, Ratpoints.Sieve.exactEnd}
  \leanok
  An end $x = \pm m\,2^{-s}$ of a search interval as \verb|exact_end| makes it, the type
  \verb|rp_end| (\csource{find_points.c}{3217}{3218}): the mantissa $m$, odd or $s = 0$, the
  rounding $r = r_h 2^{64} + r_l$ in two words, the exponent $s$ from $0$ to $127$, and whether
  $x < 0$. For a double $0 \le a < 2^{63}$, the absolute value of an end, its mantissa and exponent:
  $a = g\,2^{\mathit{ex}}$ by \verb|frexp|, $m = g\,2^{53}$ by \verb|ldexp| and the conversion to an
  \verb|unsigned long|, $s = 53 - \mathit{ex}$; $s = 0$ for $m = 0$, and else, while $s > 0$, the
  trailing zero bits of $m$ dropped, at most $s$ of them (by the loop of \verb|RP_CTZL|,
  Definition~\ref{def:ctz}); $m$ shifted up by $-s$ when $s < 0$; and $s$ capped at $127$, which
  changes neither $\lfloor b m\,2^{-s} \rfloor = 0$ nor $\lceil b m\,2^{-s} \rceil = 1$ for
  $b < 2^{63}$ (\csource{find_points.c}{3272}{3291}). The end itself, at the lower or the upper end
  of its interval: the sign by the C language's $x < 0$, so that $-0.0$ is not negative;
  $m = 2^{63}$ and $s = 0$ when $|x| < 2^{63}$ fails, an infinity included, which lies beyond every
  numerator, and otherwise the mantissa and the exponent of $|x|$; and the rounding $2^s - 1$, in
  two words for $s \ge 64$, when the end wants $b|x|$ rounded up, which is when the end is lower and
  $x \ge 0$ or upper and $x < 0$, and $0$ otherwise (\csource{find_points.c}{3264}{3271},
  \csource{find_points.c}{3292}{3299}).
\end{definition}

\begin{definition}[The scaled end]\label{def:exactScale}
  \uses{def:exactEnd, def:mulHi}
  \lean{Ratpoints.Sieve.exactScale}
  \leanok
  For a denominator $b$ and an end $\pm m\,2^{-s}$ with its rounding $r$:
  $\min(2^{64} - 1, \lfloor (b m + r)/2^s \rfloor)$, which is $b|x|$ rounded as the end wants it
  unless that does not fit $64$ bits; the product $b m$ in two words, the low one the product of the
  words and the high one by $\mathrm{mulHi}$ (Definition~\ref{def:mulHi}), $r$ added with the carry
  of the low word, the sum shifted right by $s$ in its two words, and $2^{64} - 1$ when the high
  word of the quotient is not $0$, the arm of \verb|unsigned __int128| that the builds of the model
  compile (\csource{find_points.c}{3231}{3238}, \csource{find_points.c}{3257}{3257}). As
  $b m < 2^{126}$ and $r < 2^{127}$, nothing is lost.
\end{definition}

\begin{definition}[The bounds of the search intervals]\label{def:exactBounds}
  \uses{def:exactEnd, def:exactScale}
  \lean{Ratpoints.Sieve.IntervalBounds, Ratpoints.Sieve.endForm, Ratpoints.Sieve.exactBounds}
  \leanok
  The bounds of an interval as the sieve of a denominator reads them, the type \verb|rp_bounds|
  (\csource{find_points.c}{3219}{3221}): the forms $(b\,l_M + l_C) \gg l_{\mathit{sh}}$ and
  $(b\,u_M + u_C) \gg u_{\mathit{sh}}$ of its ends, whether it is taken inline, and both ends. The
  form of an end $e$, lower or upper, for the height bound $H$ and the denominators up to
  $b_{\mathrm{high}} \ge 1$, or none, as \verb|end_form| makes it: for an end that can lie beyond
  the bound on its side ($x < 0$ at the lower end, $x \ge 0$ at the upper end), whose $|x|$ the
  scaled end at $b = 1$ shows to be at least $H$, the form $(0, -H, 0)$ at the lower end and
  $(0, H, 0)$ at the upper end; and for $s < 64$ and
  $m \le \lfloor (2^{63} - 1 - C)/b_{\mathrm{high}} \rfloor$ the form $(\pm m, C, s)$, the sign of
  $x$, with $C = 2^s - 1$ at the lower end and $0$ at the upper end, for which $b M + C$ is a
  \verb|long| for every $b \le b_{\mathrm{high}}$ (\csource{find_points.c}{3304}{3318}). The bounds
  of the search intervals, increasing, disjoint and without a NaN, as \verb|exact_bounds| makes
  them: for each interval $[l, u]$ in turn its two ends, $l$ as a lower end and $u$ as an upper end,
  and the interval taken inline, with its forms, when $H \le 2^{63} - 65$, so that the conversion to
  bits needs no split form, and both ends have a form; otherwise the forms of an empty interval, $0$
  at the lower end and $-2^{63}$ at the upper end (\csource{find_points.c}{3323}{3340}).
\end{definition}

\begin{definition}[The bits of an interval]\label{def:intervalBits}
  \uses{def:exactBounds, def:numeratorsToBits, def:numClasses}
  \lean{Ratpoints.Sieve.intervalBitsExact, Ratpoints.Sieve.intervalBits}
  \leanok
  The bits of the numerators of the interval with the bounds $\mathit{bd}$ at the denominator $b$,
  for the height bound $H$ and the class of $b$ (Definition~\ref{def:numClasses}): a code and the
  bits of the first and the last numerator. From the products in 128 bits, as
  \verb|interval_bits_exact| computes them for an interval that is not taken inline:
  $\lceil b l \rceil$, which is $-\lfloor b|l| \rfloor$ for $l < 0$, clipped to $-H$, and the code
  $2$ when it lies beyond $H$ for $l \ge 0$, so that neither this interval nor a later one holds a
  numerator within the bound; then $\lfloor b u \rfloor$, and the code $1$ when it lies below $-H$,
  so that this interval holds none, and otherwise clipped to $H$; and the code $0$ with the bits of
  $\lceil b l \rceil, \dots, \lfloor b u \rfloor$ (Definition~\ref{def:numeratorsToBits}) in the
  form that splits the ends for the heights above $2^{63} - 65$
  (\csource{find_points.c}{3381}{3406}). From the forms, as \verb|interval_bits| computes them:
  $\mathit{lo} = \lceil b l \rceil$ and $\mathit{hi} = \lfloor b u \rfloor$ by the arithmetic shift
  of the forms; the code $2$ for $\mathit{lo} > H$; for $\mathit{hi} < -H$ the code $1$ if the
  interval is taken inline and otherwise what the products in 128 bits give, since the forms of an
  interval not taken inline are those of an empty one; and otherwise the code $0$ with the bits of
  $\mathit{lo}, \dots, \mathit{hi}$ clipped to $-H, \dots, H$, in the form that does not split the
  ends (\csource{find_points.c}{3410}{3434}). The code $0$ comes also for an interval that holds no
  numerator, when $\mathit{lo} > \mathit{hi}$.
\end{definition}

The contracts speak about the ranges of the numerators that the values of the ends give, so that a
range is the same for the bits and for the values.

\begin{definition}[The ranges of the numerators]\label{def:numeratorRanges}
  \uses{def:intervalBits, def:numClasses}
  \lean{Ratpoints.Sieve.lowNumerator, Ratpoints.Sieve.highNumerator, Ratpoints.Sieve.numeratorRanges, Ratpoints.Sieve.IntervalBounds.Represents, Ratpoints.Sieve.RangesFor, Ratpoints.Sieve.inRanges}
  \leanok
  The least numerator of the denominator $b$ in an interval with the lower end $x$,
  $\lceil b x \rceil$ clamped to $-H, \dots, H + 1$, with $-H$ for $\bot$ and $H + 1$ for $\top$;
  the largest in an interval with the upper end $x$, $\lfloor b x \rfloor$ clamped to
  $-H - 1, \dots, H$, with $-H - 1$ for $\bot$ and $H$ for $\top$; and the ranges of the numerators
  of $b$ in the intervals $D$ over the values, one for each interval in turn. The clamps are
  asymmetric so that each range has $\mathit{lo} \le \mathit{hi} + 1$, the ranges stay strictly
  apart, and a range above $H$ is seen as such. The bounds of an interval \emph{represent} the range
  $\mathit{lo}, \dots, \mathit{hi}$ at $b$ when, for every class with $k \le 6$ and $a_0 < 2^k$, the
  code $2$ comes only when $\mathit{lo} > H$, the code $1$ only when $\mathit{hi} < -H$, and any
  other code with the bits of the numerators from $\max(\mathit{lo}, -H)$ to $\min(\mathit{hi}, H)$
  in the packing of the class; the bounds of the search intervals give the ranges $\mathit{rs}$ at
  $b$ when each represents its range, each range has $\mathit{lo} \le \mathit{hi} + 1$, and the
  ranges increase and are strictly apart, which is what the loop of \verb|sift| needs; and whether
  an integer lies in one of a list of ranges.
\end{definition}

\begin{lemma}[The ends and the scaled ends]\label{lem:exactEnd}
  \uses{def:exactEnd, def:exactScale, lem:mulHi, lem:Double, def:numeratorRanges}
  \lean{Ratpoints.Sieve.toNat_exactScale, Ratpoints.Sieve.toNat_hi_mul_add_lo, Ratpoints.Sieve.toNat_shift_words_of_lt, Ratpoints.Sieve.toNat_shift_words_of_le, Ratpoints.Sieve.IsExactMantissa, Ratpoints.Sieve.isExactMantissa_run_exactMantissa, Ratpoints.Sieve.ExactEnd.value, Ratpoints.Sieve.IsEnd, Ratpoints.Sieve.isEnd_run_exactEnd, Ratpoints.Sieve.numerators_run_exactEnd, Ratpoints.Sieve.toNat_exactScale_of_isEnd}
  \leanok
  For $s \le 127$ and $b m + r < 2^{128}$, the scaled end is
  $\min(2^{64} - 1, \lfloor (b m + r)/2^s \rfloor)$, from the two-word product and the shift of two
  words. The mantissa and the exponent of a double $a > 0$ are exact, $a = m\,2^{-s}$ with
  $0 < m < 2^{63}$ and $s \le 127$, or, for a tiny $a < 2^{-75}$, capped at $s = 127$; what
  $\mathrm{exactEnd}$ makes of every double is an \emph{end}, with $s \le 127$, $m \le 2^{63}$ and
  the rounding of its side, and its value $\pm m\,2^{-s}$ is the value of the double as far as the
  numerators go: for $1 \le b < 2^{63}$ and $H < 2^{63}$, the least and the largest numerators of
  $b$ at the value of the double and at the value of the end agree
  (Definition~\ref{def:numeratorRanges}), for a NaN-free double, also an infinite one and one beyond
  $2^{63}$, whose numerators lie beyond the bound either way.
\end{lemma}

\begin{proof}\leanok
  The two-word arithmetic by the lemma of the high word of a product
  (Lemma~\ref{lem:mulHi}) and the carry; the exactness of the mantissa from \verb|frexp| and
  \verb|ldexp| (Lemma~\ref{lem:Double}) and the dropped zero bits, $m\,2^{-s}$ unchanged; the
  numerators of an end beyond $2^{63}$ or tiny are those of the clamps.
\end{proof}

\begin{lemma}[The bits of an interval]\label{lem:intervalBits}
  \uses{def:intervalBits, lem:exactEnd, lem:numeratorsToBits, def:numeratorRanges, def:exactBounds}
  \lean{Ratpoints.Sieve.spec_intervalBitsExact, Ratpoints.Sieve.spec_intervalBits_of_fast, Ratpoints.Sieve.intervalBits_eq_intervalBitsExact, Ratpoints.Sieve.lowNumerator_of_endForm, Ratpoints.Sieve.highNumerator_of_endForm}
  \leanok
  For the ends of an interval, $b \le H < 2^{63}$ and a class with $k \le 6$ and $a_0 < 2^k$: the
  bits from the products in 128 bits have the code $2$ only when the least numerator lies above $H$,
  the code $1$ only when the largest lies below $-H$, and otherwise the bits of the numerators from
  the least to the largest, clipped to $[-H, H]$. For an interval taken inline and
  $H \le 2^{63} - 65$ the bits from the forms are the same with the clamped values of the forms, and
  for one not taken inline they are those from the products, the forms being those of an empty
  interval; the clamped value of a form is the least or the largest numerator of its end.
\end{lemma}

\begin{proof}\leanok
  The conversion to bits of an interval within $[-H, H]$ is Lemma~\ref{lem:numeratorsToBits}; the
  products give $\lfloor b|x| \rfloor$ and $\lceil b|x| \rceil$ by Lemma~\ref{lem:exactEnd}, and a
  form $(b M + C) \gg s$ gives $\lceil b m\,2^{-s} \rceil$ for $C = 2^s - 1$ and the floor for
  $C = 0$, a \verb|long| by the bound on $m$.
\end{proof}

\begin{lemma}[The ranges of the numerators]\label{lem:numeratorRanges}
  \uses{def:numeratorRanges, def:regionSpec, def:inIntervals}
  \lean{Ratpoints.Sieve.lowNumerator_le_iff, Ratpoints.Sieve.le_highNumerator_iff, Ratpoints.Sieve.inRanges_numeratorRanges, Ratpoints.Sieve.lowNumerator_le_highNumerator_add_one, Ratpoints.Sieve.highNumerator_lt_lowNumerator, Ratpoints.Sieve.le_add_one_of_mem_numeratorRanges, Ratpoints.Sieve.pairwise_numeratorRanges}
  \leanok
  For $b > 0$ and $|a| \le H$: $a$ is at least the least numerator of a lower end $l$ exactly when
  $l \le a/b$, and at most the largest of an upper end $u$ exactly when $a/b \le u$; so $a$ lies in
  one of the ranges of the numerators of the intervals $D$ exactly when $a/b$ lies in one of the
  intervals. For search intervals (Definition~\ref{def:regionSpec}) each range has
  $\mathit{lo} \le \mathit{hi} + 1$ and the ranges are strictly apart, since the upper end of each
  interval lies below the lower end of the next.
\end{lemma}

\begin{proof}\leanok
  The ceiling and the floor of $b x$ against an integer, with the clamps; the strict separation
  through $\lfloor b u_i \rfloor < \lceil b l_{i+1} \rceil$ for $u_i < l_{i+1}$, after the chain of
  the search intervals is turned into its pairwise form: a chain of intervals in which each ends
  below the next begins is pairwise so (\verb|List.pairwise_of_isChain_of_le|, a fact for Mathlib's
  lists, in the mirror \verb|Data/List/Pairwise| and in no node).
\end{proof}

\begin{lemma}[The bounds give the ranges]\label{lem:exactBounds}
  \uses{def:exactBounds, lem:intervalBits, lem:numeratorRanges, lem:exactEnd, def:regionSpec,
    def:inIntervals}
  \lean{Ratpoints.Sieve.intervalBoundsOf, Ratpoints.Sieve.toList_run_exactBounds, Ratpoints.Sieve.represents_intervalBoundsOf, Ratpoints.Sieve.areSearchIntervals_domainValues, Ratpoints.Sieve.rangesFor_run_exactBounds}
  \leanok
  The bounds of the search intervals are the bounds of each interval in turn; for well-formed
  intervals (Definition~\ref{def:regionSpec}), $H < 2^{63}$, $b_{\mathrm{high}} \le H$ and
  $0 < b \le b_{\mathrm{high}}$, the bounds of each represent at $b$ the range of the numerators of
  its values, the values are search intervals, and so the bounds of the run give at $b$ the ranges
  of the numerators of the values of the intervals (Definition~\ref{def:numeratorRanges}).
\end{lemma}

\begin{proof}\leanok
  The loop of \verb|exact_bounds| is a map; Lemma~\ref{lem:intervalBits} for each interval, with the
  numerators of the ends of Lemma~\ref{lem:exactEnd}, and Lemma~\ref{lem:numeratorRanges} for the
  order of the ranges.
\end{proof}

For the one interval $[-h, h]$ of a search without intervals the ranges are the one range
$-H, \dots, H$ at every denominator, and the interval is taken inline exactly when
$H \le 2^{63} - 65$: the bits and their form are those that \verb|sift| used before stage~H, so
that the output of a search without intervals does not change.

\section{The sieve over the search intervals}\label{sec:regionSieve}

\verb|sift| (Definition~\ref{def:sift}) takes the bounds of the search intervals of the run and runs
over them as the C code does: for each interval in turn the bits of its numerators
(Definition~\ref{def:intervalBits}); the code $2$ returns, the code $1$ goes on to the next
interval, and at the first interval with the code $0$, also one without numerators, the sieve specs
are made and the rows that are missing built into the store (Definitions~\ref{def:siftBits}
and~\ref{def:buildRows}), once for the denominator, the lazy set-up of the flag \verb|do_setup|;
then the points of the bits of the interval (Definition~\ref{def:siftBits}), which returns when the
callback stops the search. Its C ranges gain the struct \verb|rp_intervals|
(\csource{find_points.c}{3435}{3440}), the field \verb|iv| of the sieving thread
(\csource{rp-private.h}{574}{574}) and its clearing by \verb|worker_init|, line 3130, which joins
the link of line 3129 (\csource{find_points.c}{3129}{3130}), and its deviation ``every numerator
from $-H$ to $H$'' goes. \verb|run_setup| (Definition~\ref{def:Shape}) takes the search intervals,
the height bound and the largest denominator and returns the bounds with the shape
(\csource{find_points.c}{3869}{3870}), the Jacobi tables reading the largest denominator; the run
(Definition~\ref{def:sieveRun}) takes the bounds and passes them to each denominator, for the field
\verb|iv| that \verb|run_block| sets for each block (\csource{find_points.c}{3913}{3913}); the
sieving (Definition~\ref{def:sieving}) runs the denominators from $b_{\mathrm{low}}$ to
$b_{\mathrm{high}}$ of the shape over the intervals; and the search
(Definition~\ref{def:searchPoints}) takes the range, the search intervals and the flag
\verb|sturm_empty| (Section~\ref{sec:regionSearch}).

The contract of one denominator (\lname{Ratpoints.Sieve.sift\_spec}) keeps its hypotheses and adds
the ranges $\mathit{rs}$ that the bounds give at $b$ (Definition~\ref{def:numeratorRanges}),
quantified in the postcondition, since the program does not mention them: the store fits the entries
before and after the denominator, also when no interval asks for the rows, and the delivery is that
of the points of the check with the flags (Definition~\ref{def:checkNumWith}) of the numerators of
$-H, \dots, H$ that lie in one of the ranges, pass the class bit of the word $w$ and survive the
entries of all three stages, in increasing order. The contract of the run
(\lname{Ratpoints.Sieve.sieveRun\_spec}) quantifies ranges for each denominator from
$b_{\mathrm{low}}$ to $b_{\mathrm{high}}$, and its invariants
(\lname{Ratpoints.Sieve.SieveRunInvariant}, \lname{Ratpoints.Sieve.SieveRunLoopInvariant}) take the
bounds and the range. With the ranges of the values of well-formed intervals
(Lemma~\ref{lem:exactBounds}), a numerator $|a| \le H$ lies in a range exactly when $a/b$ lies in an
interval (Lemma~\ref{lem:numeratorRanges}), so that each denominator delivers the points of the
numerators with $a/b$ in one of the search intervals.

\begin{lemma}[The loop over the intervals]\label{lem:siftIntervals}
  \uses{def:intervalBits, def:numeratorRanges, lem:numeratorsToBits, def:siftBits, lem:siftBits,
    def:checkNumWith, lem:deliver, lem:accepted, lem:staged, lem:storeTables}
  \lean{Ratpoints.Sieve.SiftIntervalsInvariant, Ratpoints.Sieve.SiftIntervalsLoopInvariant, Ratpoints.Sieve.IsSiftArgs, Ratpoints.Sieve.flatMap_filter_intervalBits_eq, Ratpoints.Sieve.flatMap_acceptedPoints_eq_flatMap_checkNumWith, Ratpoints.Sieve.flatMap_flatMap_filter_eq_flatMap_filter_inRanges, Ratpoints.Sieve.NumClass.map_range_eq_filter_map_range, Ratpoints.Sieve.RangesFor.represents_of_eq_append_cons, Ratpoints.Sieve.RangesFor.lt_fst_of_mem_drop, Ratpoints.Sieve.filter_numerators_eq_nil_of_lt, Ratpoints.Sieve.filter_numerators_eq_nil_of_lt_neg, Ratpoints.BruteForce.pairwise_lt_numerators, List.filter_append_filter_of_pairwise, List.flatMap_filter_eq_filter_any}
  \leanok
  Over an increasing list of integers and ranges that increase and are strictly apart, the
  concatenation, range after range, of the members in the range that pass a test is the list of the
  members in one of the ranges that pass it, so that the points of the numerators come in increasing
  order within a denominator; the bits of the class from the least to the largest bit of a range
  within $[-H, H]$ are those of $-H, \dots, H$ filtered by the range, so that the points that the
  bits of an interval with the code $0$ deliver are those of the numerators in its range that pass
  the class word and the entries, the points of the check of the numerators that survive the first
  two stages (Lemma~\ref{lem:accepted}) being those of the check with the flags of the numerators
  that survive all three (Lemma~\ref{lem:staged}); no numerator of $-H, \dots, H$ lies in a range
  above $H$ or below $-H$. The invariant of the loop: before a return, the callback has not stopped
  the search, the store fits the entries, the specs are those of the set-up once it is made, and the
  delivery is that of the points of the first ranges; after the return at a stop or at an interval
  above $H$, the store fits the entries and the delivery is that of the points of all the ranges;
  with a structure of what the contract needs of its arguments, through which the steps of the loop
  read them.
\end{lemma}

\begin{proof}
  \leanok
  By induction on the ranges, with the members of one range before those of the next; the bits of a
  range by the enumeration of a residue class (Lemma~\ref{lem:numeratorsToBits}); the invariant is
  built with the program logic of Lean core, its verification conditions closed by the step lemmas
  of the invariant and of the structure, and a stop by the delivery of a prefix
  (Lemma~\ref{lem:deliver}).
\end{proof}

The denominators of the search (Definition~\ref{def:searchDens}) take the range: for odd $d$ with
$|c_d| \ne 1$ and divisors of the leading coefficient up to $b_{\mathrm{high}}$, the divisors times
squares from $b_{\mathrm{low}}$ to $b_{\mathrm{high}}$, divisor by divisor; and otherwise
$b_{\mathrm{low}}, \dots, b_{\mathrm{high}}$. The points that the search delivers
(Definition~\ref{def:searchList}) take the range and the intervals $D$ over the values, the
numerators of a denominator $b$ being those of $-H, \dots, H$ with $a/b$ in an interval of $D$ that
pass the test, and a sound test lets through the numerator of every point over the denominators of
the search in the range; and what \verb|sieving_info| computes is sound for every height
(Lemma~\ref{lem:sievingInfo}): at its exit no triple is a point of any height, so that the exit
serves every range.

\begin{lemma}[The denominators in a range]\label{lem:searchDensRange}
  \uses{def:searchDens, def:Shape, lem:shape_dens, def:bruteForce, def:binaryDivisors,
    lem:searchDenominators}
  Proposed names: \lname{Ratpoints.le\_and\_le\_of\_mem\_searchDens},
  \lname{Ratpoints.searchDens\_eq\_filter},
  \lname{Ratpoints.searchDens\_one\_eq\_searchDenominators},
  \lname{Ratpoints.bruteForce\_eq\_append}, \lname{Ratpoints.filter\_bruteForce\_eq\_append},
  \lname{Ratpoints.filter\_bruteForce\_eq\_self}. For $1 \le b_{\mathrm{low}}$ the denominators of
  the search lie in $b_{\mathrm{low}}, \dots, b_{\mathrm{high}}$, and for $1 \le b_{\mathrm{high}}$
  they are those of the search from $1$ to $b_{\mathrm{high}}$ that are at least $b_{\mathrm{low}}$;
  those from $1$ to $b_{\mathrm{high}}$ are the denominators of the search of the specification up
  to $b_{\mathrm{high}}$ (Definition~\ref{def:binaryDivisors}, Lemma~\ref{lem:searchDenominators}).
  The points of the brute-force reference with $z = 0$ or
  $b_{\mathrm{low}} \le z \le b_{\mathrm{high}} \le H$ are its points over the candidate $(1, 0)$
  followed by those over these denominators, and its points with $z = 0$ or $1 \le z \le B$ for
  $B \ge H$ are all its points.
\end{lemma}

\begin{proof}
  \uses{thm:bruteForce}
  From the lists of the shapes and the filter of a range, the divisors times squares in a range
  being those of Lemma~\ref{lem:searchDenominators}; every point of the reference has $z \le H$ by
  its membership (Theorem~\ref{thm:bruteForce}).
\end{proof}

\begin{lemma}[The points in the search intervals]\label{lem:searchListRegion}
  \uses{def:searchList, def:inIntervals, lem:searchDensRange, def:checkNumWith}
  Proposed names: \lname{Ratpoints.mem\_region\_of\_mem\_searchList},
  \lname{Ratpoints.searchList\_nil},
  \lname{Ratpoints.Sieve.snd\_snd\_eq\_zero\_of\_mem\_pointsAtInfinityList},
  \lname{Ratpoints.Sieve.fst\_eq\_and\_snd\_snd\_eq\_of\_mem\_checkNumWith}. Without a reversal,
  every member of the list that the search delivers for the intervals $D$ and the range
  $b_{\mathrm{low}} \ge 1$ to $b_{\mathrm{high}}$ has $z = 0$, or $z$ in the range and $x/z$ in an
  interval of $D$; and for no intervals the list is the points at infinity.
\end{lemma}

\begin{proof}
  The points at infinity have $z = 0$, and the check of a numerator $a$ over $b$ gives triples with
  $x = a$ and $z = b$; the denominators lie in the range (Lemma~\ref{lem:searchDensRange}).
\end{proof}

\section{The Sturm step}\label{sec:sturmStep}

The model of the Sturm step extends \verb|Ratpoints/Model/Sturm.lean|, in the namespace
\verb|Ratpoints.Sturm|, after the remainder chain of stage~D (Definition~\ref{def:sturmChain}),
whose rows it reads as the C code reads its table. The bisection keeps the \verb|long| numerators
and exponents of the C code, as \verb|Int64|, and makes the doubles of the ends as the C code makes
them.

\begin{definition}[The signs and the counts of the chain]\label{def:evalSign}
  Proposed names: \lname{Ratpoints.Sturm.evalSign}, \lname{Ratpoints.Sturm.signsAt},
  \lname{Ratpoints.Sturm.countsAtInfinity}. The sign $-1$, $0$ or $1$ of the value at
  $n/2^e$ of the polynomial of a row, as \verb|eval_sign| computes it by Horner's rule in integers:
  from the leading coefficient $w$, for the coefficients $c_k$ below it in turn
  $w \leftarrow w n + c_k 2^{e'}$ with $e' = e, 2e, \dots$, so that $w$ is $2^{\deg \cdot e}$ times
  the value (\csource{sturm.c}{67}{84}). The signs $s_0$ and $s_1$ of $f$ and $f'$ at the midpoint
  and the number of sign changes of the rows of the chain there, zeros skipped: the first two rows
  count when their signs are opposite; then, from the sign of the last row that is not zero ($s_0$
  when $s_1 = 0$), each further row counts when its sign is the opposite, and replaces it when it is
  not zero (\csource{sturm.c}{252}{260}). And the numbers of sign changes of the rows at $+\infty$
  and at $-\infty$, from the signs of their leading coefficients, two neighbouring rows counting at
  $-\infty$ when their signs differ after the sign of each row of odd degree is flipped
  (\csource{sturm.c}{165}{175}).
\end{definition}

\begin{definition}[The bisection]\label{def:iterate}
  \uses{def:evalSign, def:Double}
  Proposed names: \lname{Ratpoints.Sturm.longMax}, \lname{Ratpoints.Sturm.midpoint},
  \lname{Ratpoints.Sturm.addInterval}, \lname{Ratpoints.Sturm.iterate}. $\mathrm{LONG\_MAX} = 2^{63}
  - 1$, as the C code computes it, the numerator of the end at $+\infty$ of the bisection, and its
  negation that of the end at $-\infty$ (\csource{sturm.c}{191}{192}). The point $n_m/2^{e_m}$ at
  which the interval from $n_l/2^{e_l}$ to $n_r/2^{e_r}$ is split: $0$ for $[-\infty, +\infty]$;
  $-1$ or $2 n_r$ for $[-\infty, n_r]$ and $1$ or $2 n_l$ for $[n_l, +\infty]$, as the finite end is
  $0$ or not, with the exponent $0$; for equal exponents, which are then $0$, the half of $n_l +
  n_r$, with the exponent $0$ when the sum is even and $1$ otherwise; and else the sum of the ends
  over the larger denominator, with the larger exponent plus $1$ (\csource{sturm.c}{234}{251}). The
  interval added to the intervals kept: its ends converted to doubles,
  $\mathrm{ldexp}(\mathrm{double}(n), -e)$, and appended, or, when the last interval kept ends at
  its lower end, compared by $\mathrm{beq}$, that interval extended to its upper end
  (\csource{sturm.c}{209}{227}). And the bisection of the nested function \verb|iterate| on an
  interval with the counts $c_l$ and $c_r$ of the chain at its ends, the signs $s_l$ and $s_r$ of
  $f$ just inside them and the depth: dropped when $c_l = c_r$ and $s_l < 0$; kept when $c_l = c_r$
  and $s_l > 0$, or at the depth of the step; and otherwise split at its midpoint, where the chain
  is evaluated, the left half bisected before the right half, both with the count and the sign of
  $f$ at the midpoint, and, when the midpoint is a root of $f$, with the count plus $1$ and
  $-\mathrm{sign}\, f'$ for the left half and $\mathrm{sign}\, f'$ for the right half
  (\csource{sturm.c}{197}{208}, \csource{sturm.c}{228}{233}, \csource{sturm.c}{261}{267}). It is a
  recursion that ends since the depth grows towards the depth of the step, the first well-founded
  recursion of the model, whose proofs go through its induction principle.
\end{definition}

\begin{definition}[The intersection]\label{def:intersect}
  \uses{def:Double}
  Proposed names: \lname{Ratpoints.Sturm.meet}, \lname{Ratpoints.Sturm.intersect}. One step of the
  intersection of an interval $I = [l_I, u_I]$ that the bisection keeps with an interval
  $J = [l_J, u_J]$ of the caller, both closed: the interval it notes, if any, and whether it goes on
  with the next interval kept or the next of the caller, the one that ends first. If $l_I \le l_J$,
  $I$ is forgotten when $u_I < l_J$; otherwise $[l_J, u_I]$ is noted with the next $I$ when
  $u_I \le u_J$, and $J$ with the next $J$ otherwise. Symmetrically when $l_J < l_I$: $J$ is
  forgotten when $u_J < l_I$, and otherwise $[l_I, u_J]$ is noted with the next $J$ when
  $u_J \le u_I$, and $I$ with the next $I$ otherwise. The ends are compared as doubles and copied
  (\csource{sturm.c}{289}{322}). The intersection: from the first interval of each list, while both
  have one left, a step notes their intersection and goes on with the list whose interval ends first
  (\csource{sturm.c}{276}{276}, \csource{sturm.c}{285}{288}, \csource{sturm.c}{323}{323},
  \csource{sturm.c}{325}{325}). The intervals noted may be degenerate.
\end{definition}

\begin{definition}[The Sturm step]\label{def:computeSturm}
  \uses{def:sturmChain, def:evalSign, def:iterate, def:intersect}
  Proposed name: \lname{Ratpoints.Sturm.computeSturm}. \verb|_ratpoints_compute_sturm| on the
  polynomial $f$ of degree $d \ge 1$ with $c_d \ne 0$, to the depth of the step, for the search
  intervals of the caller: none when the remainder chain is none (Definition~\ref{def:sturmChain}),
  for a polynomial with a repeated factor; no interval when the counts at $-\infty$ and $+\infty$
  agree, so that $f$ has no real root, and $c_0 < 0$, so that $f < 0$ everywhere; and otherwise the
  intervals that the bisection keeps from $[-\infty, +\infty]$, with the counts at $\mp\infty$ at
  its ends and the sign of $c_d$, negated at $-\infty$ for odd $d$, as the signs just inside them,
  intersected with the caller's intervals, which may leave no interval (\csource{sturm.c}{86}{97},
  \csource{sturm.c}{177}{190}, \csource{sturm.c}{193}{194}, \csource{sturm.c}{269}{273},
  \csource{sturm.c}{326}{332}). It owns the clears of the table of the chain (lines 179-182 and
  327-330) under its deviation, as the chain owns its loops, and the failure for lack of memory of
  the copy of the caller's intervals with its clears (lines 278-284) is excluded.
\end{definition}

\begin{lemma}[The signs and the counts]\label{lem:evalSign}
  \uses{def:evalSign}
  Proposed names: \lname{Ratpoints.Sturm.hornerValue}, \lname{Ratpoints.Sturm.EvalSignInvariant},
  \lname{Ratpoints.Sturm.evalSign\_spec}, \lname{Ratpoints.Sturm.run\_evalSign},
  \lname{Ratpoints.Sturm.signCount}, \lname{Ratpoints.Sturm.SignsAtInvariant},
  \lname{Ratpoints.Sturm.signsAt\_spec}, \lname{Ratpoints.Sturm.run\_signsAt},
  \lname{Ratpoints.Sturm.signAtTop}, \lname{Ratpoints.Sturm.signAtBot},
  \lname{Ratpoints.Sturm.CountsInvariant}, \lname{Ratpoints.Sturm.countsAtInfinity\_spec},
  \lname{Ratpoints.Sturm.run\_countsAtInfinity}. For $0 \le e$ and $e$ times the size of the row
  below $2^{63}$, so that the exponents are \verb|long|s, the sign that Horner's rule gives is the
  sign of the value of Horner's rule in the integers after all the steps; the signs and the count at
  the midpoint are the signs of the first two rows and the count of sign changes, zeros skipped, of
  the signs of all the rows; and the counts at $\pm\infty$ are the numbers of neighbouring rows
  whose signs at $+\infty$ and at $-\infty$ differ, each the sign of a leading coefficient, flipped
  at $-\infty$ for odd degree.
\end{lemma}

\begin{proof}
  Contracts with the program logic of Lean core, each with a named invariant of its loop,
  irreducible after its lemmas.
\end{proof}

\begin{lemma}[The bisection keeps exact intervals]\label{lem:bisection}
  \uses{def:iterate, lem:Double, def:regionSpec}
  Proposed names: \lname{Ratpoints.Sturm.endValue}, \lname{Ratpoints.Sturm.HasValue},
  \lname{Ratpoints.Sturm.hasValue\_ldexp\_ofInt64},
  \lname{Ratpoints.Sturm.hasValue\_ldexp\_ofInt64\_longMax},
  \lname{Ratpoints.Sturm.hasValue\_ldexp\_ofInt64\_neg\_longMax}, \lname{Ratpoints.Sturm.EndsWF},
  \lname{Ratpoints.Sturm.endsWF\_zero}, \lname{Ratpoints.Sturm.EndsWF.hasValue\_left},
  \lname{Ratpoints.Sturm.EndsWF.hasValue\_right}, \lname{Ratpoints.Sturm.endValue\_eq\_div},
  \lname{Ratpoints.Sturm.EndsWF.halves},
  \lname{Ratpoints.Sturm.EndsWF.lt\_midpoint\_and\_midpoint\_lt},
  \lname{Ratpoints.Sturm.IsInterval}, \lname{Ratpoints.Sturm.EndsBefore},
  \lname{Ratpoints.Sturm.addInterval\_spec}, \lname{Ratpoints.Sturm.exists\_mem\_addInterval},
  \lname{Ratpoints.Sturm.exists\_mem\_addInterval\_of\_exists\_mem},
  \lname{Ratpoints.Sturm.iterate\_wf}, \lname{Int64.toInt\_mul\_of\_lt},
  \lname{Int64.toInt\_two\_mul}, \lname{Int64.toInt\_and\_one},
  \lname{Int64.toInt\_shiftRight\_one}, \lname{Int64.toInt\_one\_shiftLeft},
  \lname{Int64.toInt\_one\_shiftLeft\_ofNat}. The value of an end of the bisection is $n\,2^{-e}$,
  and $\pm 2^{63}$ for $\pm\mathrm{LONG\_MAX}$; a double has the value $q$ when it is not a NaN and
  $\mathrm{doubleValue}$ of its bits is $q$. An interval of the bisection at the depth
  $\delta \le 55$ is \emph{well formed} when it has one of four shapes, the whole line, an infinite
  interval whose finite end, with the exponent $0$, is $0$ or on its side of $0$, or two finite
  ends, equal exponents being $0$; its exponents lie from $0$ to $\delta$; its finite ends, scaled
  to the larger exponent, are at most $2^{\delta}/4$ in absolute value from the depth $1$ on; and
  its ends increase. The interval $[-\infty, +\infty]$ at the depth $0$ is well formed; the doubles
  of the ends of a well-formed interval have the values of the ends, the conversion of a \verb|long|
  of at most $2^{53}$ and the scaling by \verb|ldexp| being exact and $\pm\mathrm{LONG\_MAX}$
  becoming $\pm 2^{63}$; both halves of a well-formed interval at a depth below $55$ are well formed
  at the next depth, the arithmetic of the \verb|long|s not overflowing, and the midpoint lies
  strictly between the ends. Adding an interval with the values $L < U$ to kept intervals with
  $l \le u$, increasing and strictly apart, that end at $L$ at the latest, keeps them so and makes
  them end at $U$, keeps what they held and holds every value of $[L, U]$; and the bisection of a
  well-formed interval at a depth $\delta \le \mathrm{iter} \le 55$, after kept intervals that end
  at its lower end at the latest, gives intervals with $l \le u$, increasing and strictly apart,
  that end at its upper end at the latest.
\end{lemma}

\begin{proof}
  The six cases of the midpoint, each with the bounds of its numerators; the exactness of the
  doubles by Lemma~\ref{lem:Double}; the bisection by the induction principle of its recursion, the
  hypotheses on the kept intervals threaded through both halves.
\end{proof}

\begin{lemma}[The intersection and the step]\label{lem:intersect}
  \uses{def:intersect, def:computeSturm, lem:bisection, def:inIntervals}
  Proposed names: \lname{Ratpoints.Sturm.intersectList}, \lname{Ratpoints.Sturm.IntersectInvariant},
  \lname{Ratpoints.Sturm.intersect\_spec}, \lname{Ratpoints.Sturm.toList\_run\_intersect},
  \lname{Ratpoints.Sturm.meet\_eq}, \lname{Ratpoints.Sturm.intervalsWF\_domainBits\_iff},
  \lname{Ratpoints.Sturm.pairwise\_intersectList},
  \lname{Ratpoints.Sturm.intervalsWF\_intersectList},
  \lname{Ratpoints.Sturm.inIntervals\_of\_inIntervals\_intersectList},
  \lname{Ratpoints.Sturm.inIntervals\_intersectList},
  \lname{Ratpoints.Sturm.length\_intersectList\_le},
  \lname{Ratpoints.Sturm.inIntervals\_domainValues\_iff},
  \lname{Ratpoints.Sturm.run\_computeSturm\_eq\_none\_iff},
  \lname{Ratpoints.Sturm.run\_computeSturm\_of\_chain\_eq\_some},
  \lname{Ratpoints.Sturm.computeSturm\_wf}, \lname{Array.drop\_toList\_eq\_getElem!\_cons}. The
  intersection is a recursion on the two lists, of which the loop computes the list; well-formed
  intervals are those without a NaN end, with $l \le u$ by the values, and pairwise strictly apart;
  for well-formed lists the intersection is well formed, a rational number lies in it exactly when
  it lies in both lists, and it has at most the sum of their lengths less one intervals. The step is
  none exactly when the chain is none; with a chain it is the exit or the intersection of the
  bisection from $[-\infty, +\infty]$; and at a depth of at most $55$, for well-formed intervals of
  the caller, its intervals are well formed and lie in the caller's intervals.
\end{lemma}

\begin{proof}
  The loop of the intersection by its named invariant; the order of the noted intervals from the
  four cases of a step, the chain of the caller's search intervals turned into its pairwise form
  (\verb|List.pairwise_of_isChain_of_le|); the step by unfolding it at the result of the remainder
  chain (Definition~\ref{def:sturmChain}), and its intervals at a depth of at most $55$ from
  Lemma~\ref{lem:bisection} and the intersection.
\end{proof}

\section{Sturm's theorem and the coverage}\label{sec:coverage}

The coverage, that the intervals of the step hold every $x$ of the caller's intervals with
$f(x) \ge 0$, is the place where Sturm's theorem enters. The general facts are proved at Mathlib's
generality, in mirrors under \verb|Ratpoints/Mathlib| that are candidates for Mathlib: on the
division of polynomials over a field (\verb|Algebra/Polynomial/FieldDivision|), on the Sturm
sequence (\verb|Algebra/Polynomial/Sturm/Sequence|), on the signs of a real polynomial near a point
and near $\pm\infty$ (\verb|Analysis/Polynomial/Basic|) and on squarefree polynomials over a perfect
field (\verb|FieldTheory/Perfect|). Mathlib after the version that the project pins has
\verb|Polynomial.signRight x p| (pull request \#44248,
\verb|Mathlib/Analysis/Polynomial/SignRight.lean|), the sign of $p$ just right of $x$, defined as a
limit for every $p$, which the project can take over at its next update of Mathlib; until then the
mirror of the signs defines no sign of its own just right or just left of a point, and its lemmas
state the sign there by the values of $p$ and $p'$ at the point (Lemma~\ref{lem:polySigns}). The
mirror on chains of pairs (\verb|Data/List/Pairwise|) gains a fact that serves the ranges of the
numerators and the intersection, not the coverage.

\begin{lemma}[Positive multiples and the Sturm sequence]\label{lem:posAssoc}
  Proposed names: \lname{Polynomial.PosAssoc}, \lname{Polynomial.PosAssoc.refl},
  \lname{Polynomial.PosAssoc.symm}, \lname{Polynomial.PosAssoc.trans},
  \lname{Polynomial.PosAssoc.neg\_mod}, \lname{Polynomial.PosAssoc.sign\_eval},
  \lname{Polynomial.PosAssoc.sign\_leadingCoeff}, \lname{Polynomial.sturmSeq\_map},
  \lname{Polynomial.forall₂\_posAssoc\_sturmSeq},
  \lname{Polynomial.forall₂\_posAssoc\_sturmSeq\_of\_getElem},
  \lname{Polynomial.signVariations\_map\_eq\_of\_forall₂},
  \lname{Polynomial.signVariations\_map\_eval\_eq\_of\_forall₂},
  \lname{Polynomial.signVariations\_map\_leadingCoeff\_eq\_of\_forall₂},
  \lname{Polynomial.signVariations\_map\_neg\_one\_pow\_mul\_leadingCoeff\_eq\_of\_forall₂},
  \lname{Polynomial.mod\_eq\_of\_dvd\_sub\_of\_degree\_lt}, \lname{Polynomial.C\_mul\_mod},
  \lname{Polynomial.mod\_C\_mul}. Over an ordered field, $q$ is a \emph{positive multiple} of $p$
  when $q = c\,p$ for a constant $c > 0$, an equivalence relation that keeps the signs of the values
  and of the leading coefficient, and the negated remainders: $-p \bmod q$ and $-p' \bmod q'$ are
  positive multiples for positive multiples $p'$ of $p$ and $q'$ of $q$. The Sturm sequences of
  positive multiples are, entry by entry, positive multiples; a list of nonzero polynomials whose
  last entry is a unit and in which each entry from the third on is a positive multiple of the
  negated remainder of the two before it is, entry by entry, a list of positive multiples of the
  Sturm sequence of its first two entries; the Sturm sequence commutes with a map of fields; and
  lists of positive multiples have the same sign changes of their values, of their leading
  coefficients and of these times $(-1)^{\deg}$. The remainder by $q$ is $r$ when $q$ divides $p -
  r$ and $\deg r < \deg q$, and a constant factor passes through it.
\end{lemma}

\begin{proof}
  By the induction principle of the Sturm sequence, with the remainder of a scaled polynomial by a
  scaled one; the signs by the signs of the factors.
\end{proof}

\begin{lemma}[The chain is a Sturm sequence]\label{lem:chainSturm}
  \uses{lem:posAssoc, lem:evalSign, def:sturmChain, lem:sturmChain_squarefree,
    lem:squarefree_binaryFormPoly, lem:sturmCounts}
  Proposed names: \lname{Ratpoints.Sturm.IsChain.posAssoc\_neg\_mod},
  \lname{Ratpoints.Sturm.rowPolys}, \lname{Ratpoints.Sturm.chainPoly},
  \lname{Ratpoints.Sturm.IsChain.forall₂\_posAssoc\_sturmSeq},
  \lname{Ratpoints.Sturm.IsChain.forall₂\_posAssoc\_sturmSeq\_map},
  \lname{Ratpoints.Sturm.hornerValue\_eq\_eval},
  \lname{Ratpoints.Sturm.run\_evalSign\_eq\_sign\_eval},
  \lname{Ratpoints.Sturm.signCount\_eq\_signVariations},
  \lname{Ratpoints.Sturm.run\_signsAt\_eq\_signVariations},
  \lname{Ratpoints.Sturm.run\_countsAtInfinity\_eq\_signVariations},
  \lname{Ratpoints.Sturm.IsChain.run\_signsAt\_eq},
  \lname{Ratpoints.Sturm.IsChain.run\_countsAtInfinity\_eq}. Each row from the third on of a chain
  (Definition~\ref{def:sturmChain}) is a positive multiple of the negated remainder of the two rows
  before it, so that the polynomials of the rows of a chain that ends with a constant are, entry by
  entry, positive multiples of the Sturm sequence of $f$ and $f'$ over $\QQ$, and their images over
  $\mathbb{R}$ of that of $f$, read over $\mathbb{R}$ (the real polynomial of the coefficients).
  Horner's rule computes $2^{e \deg}$ times the value at $n/2^e$, so that its sign is that of the
  value; the count at the midpoint is the number of sign changes of the values of the rows there
  when $f$ or $f'$ does not vanish, which holds for a squarefree $f$; and the counts at $\pm\infty$
  are those of the signs of the leading coefficients. So, for a chain that ends with a constant and
  $d < 2^{57}$: for a squarefree $f$ and the depth of the midpoint at most $55$, the signs and the
  count of the step at a midpoint $m$ are $\mathrm{sign}\,f(m)$, $\mathrm{sign}\,f'(m)$ and the
  count $V_f(m)$ of the Sturm sequence of $f$ and $f'$ (Lemma~\ref{lem:sturmCounts}); and its counts
  at $\pm\infty$ are those of the Sturm sequence at $\pm\infty$.
\end{lemma}

\begin{proof}
  \uses{lem:polySigns}
  The relation $A s_{k-2} = Q s_{k-1} - g s_k$ with $A, g \ge 1$ of the chain (stage~D) makes $s_k$
  a positive multiple of $-s_{k-2} \bmod s_{k-1}$; Lemma~\ref{lem:posAssoc} gives the rest. A
  squarefree $f$ and $f'$ do not vanish together (Lemma~\ref{lem:polySigns}). The bound $d < 2^{57}$
  keeps the exponents of Horner's rule, $e \deg$ with $e \le 55$, \verb|long|s.
\end{proof}

\begin{lemma}[The signs of a real polynomial]\label{lem:polySigns}
  Proposed names: \lname{Polynomial.eventually\_atTop\_sign\_eval},
  \lname{Polynomial.eventually\_atBot\_sign\_eval},
  \lname{Polynomial.eventually\_atTop\_signVariations},
  \lname{Polynomial.eventually\_atBot\_signVariations},
  \lname{Polynomial.eventually\_sign\_eval\_eq},
  \lname{Polynomial.exists\_eq\_X\_sub\_C\_mul\_of\_isRoot},
  \lname{Polynomial.eventually\_sign\_eval\_nhdsGT},
  \lname{Polynomial.eventually\_sign\_eval\_nhdsLT},
  \lname{Polynomial.exists\_isRoot\_of\_eval\_mul\_eval\_neg},
  \lname{Polynomial.eval\_neg\_of\_forall\_not\_isRoot},
  \lname{Polynomial.length\_le\_natDegree\_add\_one\_of\_isChain},
  \lname{Squarefree.eval\_ne\_zero\_or\_eval\_derivative\_ne\_zero}. A nonzero real polynomial has,
  near $+\infty$, the sign of its leading coefficient and, near $-\infty$, that sign times
  $(-1)^{\deg}$, and a list of nonzero polynomials the corresponding sign changes; near a point
  where $p$ does not vanish $p$ keeps its sign; at a point $x$ with $p(x) \ne 0$ or $p'(x) \ne 0$,
  the sign of $p$ just right of $x$ is that of $p(x)$, or of $p'(x)$ when $p(x) = 0$, and just left
  of $x$ that of $p(x)$, or of $-p'(x)$ when $p(x) = 0$, each lemma stating the sign by this case
  distinction in its conclusion; values of opposite signs at $a < b$ give a root in between; a
  polynomial without a root in $(l, r)$ that is negative just right of $l$ is negative on $(l, r)$;
  and a chain of points with values of alternating signs has at most $\deg p + 1$ members. A
  squarefree polynomial over a perfect field and its derivative never vanish together.
\end{lemma}

\begin{proof}
  By the dominant term at $\pm\infty$, continuity, the factor $X - x$ at a root, the intermediate
  value theorem, and the count of the roots between the points of a chain.
\end{proof}

\begin{lemma}[The counts at the ends of an interval]\label{lem:sturmCounts}
  \uses{def:sturmTheorem, lem:polySigns}
  Proposed names: \lname{Ratpoints.sturmV}, \lname{Ratpoints.vBot}, \lname{Ratpoints.vTop},
  \lname{Ratpoints.vLeft}, \lname{Ratpoints.vRight}, \lname{Ratpoints.sgnLeftEnd},
  \lname{Ratpoints.sgnRightEnd}, \lname{Ratpoints.Inside}, \lname{Ratpoints.sturmV\_eq\_iff},
  \lname{Ratpoints.vLeft\_eq\_vRight\_iff}, \lname{Ratpoints.eval\_neg\_of\_vLeft\_eq\_vRight},
  \lname{Ratpoints.exists\_inside\_eval\_pos}. For a real polynomial $p$: the count $V_p(x)$ of
  Definition~\ref{def:sturmTheorem}, the counts at $\pm\infty$ from the leading coefficients, the
  count at the lower end $l$ of an interval ($V_p(l)$, or that at $-\infty$) and at its upper end
  $r$ ($V_p(r)$ plus $1$ at a root of $p$, or that at $+\infty$), the signs of $p$ just inside the
  ends, at a finite end that of $p$ there or, at a root, of $p'$ (just right of $l$) or of $-p'$
  (just left of $r$), the case distinction of Lemma~\ref{lem:polySigns}, and the points strictly
  between the ends. Assuming Sturm's theorem, for a squarefree $p$: the counts at the ends of an
  interval agree exactly when $p$ has no root strictly between them; when they agree and $p$ is
  negative just inside the lower end, $p$ is negative between the ends and just inside the upper
  end; and when they differ, or $p$ is positive just inside the lower end, $p$ is positive somewhere
  between the ends.
\end{lemma}

\begin{proof}
  Sturm's theorem on $(l, r]$ counts the roots, the root at $r$ itself giving the $1$; the counts at
  $\pm\infty$ are those at points far enough out (Lemma~\ref{lem:polySigns}); the signs between the
  ends follow from the absence of roots, and a root between them gives a point of positive value by
  its sign change, $p$ being squarefree.
\end{proof}

\begin{lemma}[The Sturm step keeps the points]\label{lem:coverage}
  \uses{lem:sturmCounts, lem:chainSturm, lem:bisection, lem:intersect, def:computeSturm,
    def:binaryForm}
  Proposed names: \lname{Ratpoints.Sturm.lowerEnd}, \lname{Ratpoints.Sturm.upperEnd},
  \lname{Ratpoints.Sturm.HalvingInv}, \lname{Ratpoints.Sturm.Kept}, \lname{Ratpoints.Sturm.SignsOf},
  \lname{Ratpoints.Sturm.HalvingInv.not\_kept}, \lname{Ratpoints.Sturm.HalvingInv.halves},
  \lname{Ratpoints.Sturm.kept\_halves}, \lname{Ratpoints.Sturm.covers\_addInterval},
  \lname{Ratpoints.Sturm.covers\_iterate}, \lname{Ratpoints.Sturm.squarefree\_chainPoly},
  \lname{Ratpoints.Sturm.halvingInv\_start}, \lname{Ratpoints.Sturm.kept\_top},
  \lname{Ratpoints.Sturm.eval\_neg\_of\_vBot\_eq\_vTop}, \lname{Ratpoints.Sturm.IsChain.signsOf},
  \lname{Ratpoints.Sturm.isChain\_of\_run\_chain\_eq\_some},
  \lname{Ratpoints.Sturm.covers\_computeSturm}, \lname{Ratpoints.Sturm.binaryForm\_eq\_mul\_eval},
  \lname{Ratpoints.Sturm.eval\_chainPoly\_nonneg}. The invariant of the bisection for $f$, read over
  $\mathbb{R}$: a well-formed interval (Lemma~\ref{lem:bisection}) whose counts and signs are those
  of $f$ at its ends (Lemma~\ref{lem:sturmCounts}); the values that the bisection of an interval
  keeps, between the values of the doubles of its ends, at a finite end only when $f$ is positive
  just inside it; and the agreement of the signs and the count of the step with those of $f$ at the
  midpoints of depth at most $55$ (Lemma~\ref{lem:chainSturm}). Assuming Sturm's theorem and for a
  squarefree $f$: a dropped interval keeps no $x$ with $f(x) \ge 0$; the halves of an interval with
  the invariant have it, and keep what it keeps; so the bisection to a depth of at most $55$ keeps
  what the intervals kept before hold and every $x$ with $f(x) \ge 0$ that the interval keeps. The
  real polynomial of a chain that does not fail is squarefree, the interval $[-\infty, +\infty]$ has
  the invariant and keeps every $|x| \le 2^{63}$, and at the exit of the step $f < 0$ everywhere.
  So, assuming Sturm's theorem, the Sturm step at a depth of at most $55$ on a polynomial of degree
  $1 \le d < 2^{57}$ with $c_d \ne 0$, for well-formed intervals of the caller, leaves intervals
  that hold every $x \in \QQ$ with $|x| \le 2^{63}$ and $f(x) \ge 0$ that lies in the caller's
  intervals. For $b > 0$, $F(a, b) = b^{d + (d \bmod 2)} f(a/b)$, so that a point gives
  $f(a/b) \ge 0$.
\end{lemma}

\begin{proof}
  \uses{lem:polySigns}
  The bisection by the induction principle of its recursion: a dropped interval by
  Lemma~\ref{lem:sturmCounts}; a split interval by the counts and the signs at the midpoint, which
  the step evaluates as those of $f$, the one-sided ones at a root of $f$ by
  Lemma~\ref{lem:polySigns}; a kept interval holds what it keeps by the values of the doubles of its
  ends (Lemma~\ref{lem:bisection}). The intersection keeps what lies in both lists
  (Lemma~\ref{lem:intersect}). The bound $|x| \le 2^{63}$ is needed: the bisection covers nothing
  beyond $\pm 2^{63}$, the doubles of $\pm\mathrm{LONG\_MAX}$.
\end{proof}

\begin{lemma}[The arrays of the Sturm step]\label{lem:sturmBuffers}
  \uses{lem:coverage, lem:polySigns, lem:intersect}
  Proposed names: \lname{Ratpoints.Sturm.Witnessed}, \lname{Ratpoints.Sturm.witnessed\_iterate},
  \lname{Ratpoints.Sturm.length\_le\_of\_witnessed},
  \lname{Ratpoints.Sturm.size\_run\_iterate\_le\_of\_chain\_eq\_some},
  \lname{Ratpoints.size\_run\_iterate\_le}, \lname{Ratpoints.size\_run\_computeSturm\_le}. An
  interval of doubles is \emph{witnessed} for a real polynomial $p$ when its ends have values in
  $[-2^{63}, 2^{63}]$ and $p$ is positive at a point that lies at or above the lower end unless that
  end has the value $-2^{63}$, and at or below the upper end unless it has the value $2^{63}$, the
  values of the ends at infinity of the bisection. Each interval that the bisection keeps is
  witnessed for $f$, read over $\mathbb{R}$; witnessed intervals that are increasing and strictly
  apart and hold every value of $[-2^{63}, 2^{63}]$ with $p \ge 0$ are at most
  $1 + \lfloor n/2 \rfloor$ for a nonzero $p$ of degree $n$, since their positive points and the
  negative points between them alternate in sign. So, assuming Sturm's theorem, for a polynomial of
  degree $1 \le d < 2^{57}$ with $c_d \ne 0$ whose chain succeeds, the bisection to a depth of at
  most $55$ keeps at most $1 + \lfloor d/2 \rfloor$ intervals, the size of the array \verb|ivlocal|
  on the stack, and the step leaves at most $\mathrm{num\_inter} + d$ intervals, the room that the
  manual asks of the caller.
\end{lemma}

\begin{proof}
  The witnesses by the induction principle of the bisection, a kept interval having a positive point
  by Lemma~\ref{lem:sturmCounts}; the count by
  \verb|Polynomial.length_le_natDegree_add_one_of_isChain| on the chain of the witnesses and the
  negative points between them; the intersection by its bound on the length
  (Lemma~\ref{lem:intersect}).
\end{proof}

\section{The search in the search region}\label{sec:regionSearch}

The search with a callback (Definition~\ref{def:findPointsWith}) runs in the order of
\verb|find_points_work_1|: the refusals of the length and the height and of the leading zeros, as
before; the normalization of the options (Definition~\ref{def:normalize}); the reversal check with
the normalized switch (Definition~\ref{def:reversal}), which may refuse the form or reverse the
polynomial; the search intervals (Definition~\ref{def:searchIntervals}), with the refusal
$\mathrm{badArgs}$; the Sturm step (Definition~\ref{def:sturmRegion}), which, when it is on
($\mathrm{sturm} \ge 0$), may refuse the form, return at once, or narrow the search intervals; and
the search of the polynomial handed on (Definition~\ref{def:searchPoints}) over the normalized range
of the denominators and the intervals that the step leaves, with its flag \verb|sturm_empty|. The
test of the \verb|long|s, which the C code makes between the reversal check and the search
intervals, stays with the search, which uses it. The search returns after the points at infinity
when the step has left no interval (line 5106, which Definition~\ref{def:searchPoints} owns), and
what a search gives back (Definition~\ref{def:Outcome}) gains the search intervals as the search
leaves them, for message~8: $[-h, h]$ when the caller gives none, narrowed by the Sturm step, in the
coordinates of the polynomial searched, and none after a return of the step. The search with a
callback gives the C range of the Sturm step, \csource{find_points.c}{4857}{4892}, to the new
component that models it.

\begin{definition}[The Sturm step of the search]\label{def:sturmRegion}
  \uses{def:computeSturm, def:reversal, def:Double}
  Proposed name: \lname{Ratpoints.sturmRegion}. The Sturm step of \verb|find_points_work_1| on the
  polynomial $f$ of degree $d$ that the reversal check hands on, with the depth of the step, for the
  search intervals: with the step off, the intervals as they are; with it on, the refusal
  $\mathrm{notSquarefree}$ when the remainder chain of $f$ is none
  (Definition~\ref{def:computeSturm}), and otherwise the intervals that the step leaves when there
  are some; when there are none, $f$ being negative on all the search intervals, none for an even
  $d$ whose leading coefficient is not a square, when there is no point at infinity either, so that
  the search returns at once with the count $0$ (the return of line 4889), and otherwise no
  intervals with the flag \verb|sturm_empty| set, so that the search delivers the points at infinity
  only (\csource{find_points.c}{4857}{4892}, from the search with a callback). The square test is
  that of the reversal check (Definition~\ref{def:reversal}), made on the polynomial searched, which
  gives the flag \verb|lcfsq| of the C code for an even degree, the only one that reads it.
\end{definition}

\begin{lemma}[The Sturm step of the search]\label{lem:sturmRegion}
  \uses{def:sturmRegion, lem:intersect, lem:coverage, lem:normalize, def:reversal, def:maxDegree}
  Proposed names: \lname{Ratpoints.run\_sturmRegion\_of\_neg},
  \lname{Ratpoints.run\_sturmRegion\_of\_nonneg},
  \lname{Ratpoints.run\_sturmRegion\_eq\_error\_iff},
  \lname{Ratpoints.run\_sturmRegion\_eq\_ok\_none}, \lname{Ratpoints.run\_sturmRegion\_subset},
  \lname{Ratpoints.run\_sturmRegion\_covers},
  \lname{Ratpoints.not\_nonneg\_and\_chain\_eq\_none\_of\_run\_sturmRegion\_eq\_ok},
  \lname{Ratpoints.lt\_two\_pow\_57\_of\_reversal\_eq\_ok},
  \lname{Ratpoints.Setup.degree\_le\_of\_reversal\_eq\_ok},
  \lname{Ratpoints.maxDegree\_lt\_two\_pow\_57}. With the step off it gives the intervals as they
  are; it refuses exactly when it is on and the chain is none, with $\mathrm{notSquarefree}$; it
  returns at once only for an even degree whose leading coefficient is not a square; at a depth of
  at most $55$, for well-formed intervals, the intervals that it leaves are well formed, none when
  \verb|sturm_empty| is set, and lie in the given intervals; and, assuming Sturm's theorem when the
  step is on, at a depth of at most $55$, for $1 \le d < 2^{57}$, $d$ the degree of $f$, and
  well-formed intervals, every $a/b$ with $b > 0$, $|a| \le 2^{63}$, $F(a, b) \ge 0$ and $a/b$ in
  the given intervals lies in the intervals that it leaves, the flag \verb|sturm_empty| clear. The
  degree that the reversal check hands on is at most the degree, so below $2^{57}$ for at most
  $\mathrm{maxDegree} + 1$ coefficients, the decided fact $\mathrm{maxDegree} < 2^{57}$ of
  \verb|Ratpoints/Model/Constants.lean|.
\end{lemma}

\begin{proof}
  The paths of the step, with Lemma~\ref{lem:intersect} for the intervals and
  Lemma~\ref{lem:coverage} for the coverage; the depth bound $55$ comes from the normalization
  (Lemma~\ref{lem:normalize}), and without it the ends would not all be doubles. At the return and
  with \verb|sturm_empty| no $a/b$ with $F(a, b) \ge 0$ lies in the given intervals, since the
  coverage would put it into an interval left.
\end{proof}

\begin{definition}[The search region that the search uses]\label{def:searchRegion}
  \uses{def:searchIntervals, def:sturmRegion, def:normalize, def:inIntervals}
  Proposed name: \lname{Ratpoints.searchRegion}. For the polynomial $f$ of degree $d$ that the
  reversal check hands on, the height bound $H$ and the options: none when the search intervals
  refuse the caller's intervals or the Sturm step, at the normalized depth, refuses or returns at
  once; no interval when it sets \verb|sturm_empty|; and otherwise the values of the intervals that
  it leaves. A definition of the proof files, not of the model; the theorems that do not assume
  Sturm's theorem are stated in its terms.
\end{definition}

\begin{lemma}[The search region that the search uses]\label{lem:searchRegion}
  \uses{def:searchRegion, lem:sturmRegion, lem:searchIntervals}
  Proposed names: \lname{Ratpoints.searchRegion\_subset}, \lname{Ratpoints.searchRegion\_covers},
  \lname{Ratpoints.searchRegion\_eq\_none},
  \lname{Ratpoints.searchRegion\_eq\_of\_run\_sturmRegion\_eq}. For well-formed intervals of the
  caller: every $x$ in the region that the search uses lies in the caller's intervals, or the caller
  gives none; for a chain that does not fail when the step is on, the region is none only when the
  Sturm step returns at once, for an even degree whose leading coefficient is not a square; and,
  assuming Sturm's theorem when the step is on, for $1 \le d < 2^{57}$, $1 \le H < 2^{63}$ and a
  chain that does not fail when the step is on, every $a/b$ with $b > 0$, $|a| \le H$ and $F(a, b)
  \ge 0$ that lies in the caller's intervals, or for none at all, lies in the region that the search
  uses, which is then not none.
\end{lemma}

\begin{proof}
  Lemma~\ref{lem:sturmRegion} on the intervals of Lemma~\ref{lem:searchIntervals}, the default
  interval holding every $a/b$ with $|a| \le H$.
\end{proof}

The list of a successful search (Lemma~\ref{lem:findPoints_searchList},
\lname{Ratpoints.exists\_searchList\_of\_findPoints\_eq\_ok}) takes the region: for well-formed
parameters, if $\mathrm{findPoints}(c, H, P, o)$ returns $L$, then the reversal check with the
normalized switch hands on a polynomial $c'$ of positive degree with the flag $\mathrm{rev}$,
$1 \le H < 2^{63}$, the caller's intervals are well formed, the chain of $c'$ does not fail when the
Sturm step is on, and there is a sound test $S$ over the normalized range such that $L$ is the list
that the search delivers for $S$ over the region that the search uses, the empty list when that is
none, and the same holds for every choice of the options of the output. The output of the search
(Theorem~\ref{thm:findPoints_order}, \lname{Ratpoints.exists\_eq\_of\_findPoints\_eq\_ok}) gives,
with the exact check on, $L$ as the points that the reference of $c'$ finds over the candidate
$(1, 0)$ and over the denominators of its search in the normalized range with $a/b$ in the region
that the search uses, mapped back after a reversal and filtered by $\mathrm{noY}$. Neither assumes
Sturm's theorem, and from them, with Lemma~\ref{lem:searchRegion}, follow the three statements that
keep their text or need no such assumption: no duplicates; every member under
\verb|RATPOINTS_NO_CHECK| lies in the search region of the caller, since the region that the search
uses lies in it; and the sublist under \verb|RATPOINTS_NO_CHECK|, since both searches use the same
region. The search region of the caller and the region that the search uses differ only by the
Sturm step, and on the points the difference vanishes under Sturm's theorem.

\begin{lemma}[The reversal and the search region]\label{lem:reversalRegion}
  \uses{def:reversal, lem:normalize, def:regionSpec, lem:reversal_swap}
  Proposed names: \lname{Ratpoints.exists\_reversal\_normalize\_eq\_error\_iff},
  \lname{Ratpoints.isFullBox\_of\_reversal\_normalize\_eq\_ok}. The reversal check with the
  normalized switch refuses exactly when the reversal is allowed, the search region is the whole
  box, $d$ is even and $c_0 = c_1 = 0$; and when it reverses, the search region is the whole box.
\end{lemma}

\begin{proof}
  The normalized switch is the switch and the whole box (Lemma~\ref{lem:normalize}), and the check
  refuses and reverses only with the switch on.
\end{proof}

\begin{lemma}[The filter by the search region]\label{lem:regionFilter}
  \uses{def:regionSpec, lem:normalize, def:swapPoint, lem:swapPoint_bijection, def:bruteForce,
    lem:searchDensRange, lem:reversalRegion, lem:sift_eq_checkDen, thm:findPoints_order,
    thm:bruteForce, lem:reversal_swap}
  Proposed names: \lname{Ratpoints.inRegion\_ite\_swapPoint\_of\_mem\_pointsOfDen\_zero},
  \lname{Ratpoints.inRegion\_ite\_swapPoint\_iff},
  \lname{Ratpoints.map\_append\_flatMap\_filter\_eq\_filter\_inRegion},
  \lname{Ratpoints.perm\_map\_append\_flatMap\_searchDens},
  \lname{Ratpoints.eq\_zero\_or\_le\_toNat\_bHigh\_normalize\_of\_inRegion},
  \lname{Ratpoints.filter\_filter\_bruteForce\_eq},
  \lname{Ratpoints.nodup\_map\_bruteForce\_of\_reversal\_eq\_ok}. When a reversal implies the whole
  box: the points at infinity, mapped, lie in the search region; a point over a denominator
  $1 \le b \le B$, $B$ the largest denominator of the search, with $|x| \le H$, mapped, lies in the
  search region exactly when $b$ is at least the normalized least denominator and $x/b$ lies in a
  caller's interval, or there are none; so the points over the candidate $(1, 0)$ and over
  denominators up to $B$ that are at least the least denominator, with the numerators filtered by
  the caller's intervals, mapped, are those over all the denominators up to $B$ and all the
  numerators, mapped and filtered by the search region. The points over the candidate $(1, 0)$ and
  the denominators of the search up to $B$, mapped, are a permutation of the points of the reference
  with $z = 0$ or $1 \le z \le B$; a point of height at most $H$ in the search region has $z = 0$ or
  $1 \le z \le B$; the filter by $1 \le z \le B$ before the filter by the search region changes
  nothing; and the reference of the polynomial that the reversal check hands on, mapped back after a
  reversal, has no duplicates, being a permutation of the reference of $c$.
\end{lemma}

\begin{proof}
  The search region of a mapped point: after a reversal the region is the whole box, which holds
  every point, and otherwise the map is the identity; the points of the reference over a denominator
  filtered by the numerators as in Lemma~\ref{lem:sift_eq_checkDen}; the permutation from that of
  the output of the search (Theorem~\ref{thm:findPoints_order}); the rest by the filters of
  Lemma~\ref{lem:searchDensRange} and the bijection of the reversal
  (Lemma~\ref{lem:swapPoint_bijection}), with the reference without duplicates
  (Theorem~\ref{thm:bruteForce}) and the reflected polynomial of the reversal
  (Lemma~\ref{lem:reversal_swap}).
\end{proof}

\begin{lemma}[The region of the caller under Sturm's theorem]\label{lem:elimSearchRegion}
  \uses{lem:searchRegion, lem:findPoints_searchList, thm:findPoints_order, lem:regionFilter,
    def:sturmTheorem, def:regionSpec}
  Proposed names: \lname{Ratpoints.elim\_searchRegion\_pointsOfDen\_zero},
  \lname{Ratpoints.elim\_searchRegion\_pointsOfDen\_filter},
  \lname{Ratpoints.exists\_eq\_filter\_inIntervals\_of\_findPoints\_eq\_ok},
  \lname{Ratpoints.exists\_eq\_filter\_inRegion\_of\_findPoints\_eq\_ok}. For well-formed intervals
  of the caller and a chain that does not fail when the step is on, the points at infinity over the
  region that the search uses are the points at infinity; and assuming Sturm's theorem when the step
  is on, for $0 < d < 2^{57}$, $d$ the degree of $c$, and $1 \le H < 2^{63}$, the points over a
  denominator $b > 0$ of the numerators with $a/b$ in the region that the search uses are those of
  the numerators with $a/b$ in the caller's intervals, or of all of them when there are none. So,
  assuming Sturm's theorem when the step is on, for well-formed parameters, the exact check on and
  $\mathrm{findPoints}(c, H, P, o) = L$: $L$ is the list of the points that the reference of $c'$
  finds over the candidate $(1, 0)$ and over the denominators of its search in the normalized range
  with the numerators in the caller's intervals, mapped back and filtered by $\mathrm{noY}$; and $L$
  is the list of the points over the candidate $(1, 0)$ and over the denominators of the search of
  $c'$ up to the largest denominator $B$, mapped back and filtered by $\mathrm{noY}$ and the search
  region of the caller (Definition~\ref{def:regionSpec}).
\end{lemma}

\begin{proof}
  A point over $b$ has $F(a, b) \ge 0$, so that the coverage puts $a/b$ into the region that the
  search uses exactly when it lies in the caller's intervals (Lemma~\ref{lem:searchRegion}, with
  $d < 2^{57}$ from Lemma~\ref{lem:sturmRegion}); the output of the search
  (Theorem~\ref{thm:findPoints_order}) and Lemma~\ref{lem:regionFilter} for the second form.
\end{proof}

From the last form come the theorems that assume Sturm's theorem: membership
(\lname{Ratpoints.mem\_iff\_isPoint\_of\_findPoints\_eq\_ok}): for well-formed parameters, the exact
check on and $\mathrm{findPoints}(c, H, P, o) = L$, a triple lies in $L$ exactly when it is a
normalized point of height at most $H$, with $y \ge 0$ under $\mathrm{noY}$, in the search region of
$o$; the set of elements (\lname{Ratpoints.toFinset\_eq\_filter\_points\_of\_findPoints\_eq\_ok},
for \verb|toFinset_eq_points_of_findPoints_eq_ok|), with both options of the output off, is
$\mathrm{points}(c, H)$ filtered by the search region; the permutation
(\lname{Ratpoints.perm\_filter\_bruteForce\_of\_findPoints\_eq\_ok}, for
\verb|perm_bruteForce_of_findPoints_eq_ok|), now for either setting of $\mathrm{noY}$, of the
reference filtered by $\mathrm{noY}$ and the search region; the reference itself
(\lname{Ratpoints.eq\_filter\_bruteForce\_of\_findPoints\_eq\_ok\_of\_reverse\_eq\_false}, for
\verb|eq_bruteForce_of_findPoints_eq_ok_of_reverse_eq_false|, now in the node of the permutation,
Theorem~\ref{thm:findPoints_perm}), filtered by the search region, when the reversal is forbidden
and the denominators of the search are those of the range; the order theorem
(\lname{Ratpoints.eq\_orderedPoints\_of\_findPoints\_eq\_ok}): $L$ is the list of the points in the
order of ratpoints for $\mathrm{maxDivisors}$, the switch of the reversal and the whole box, $c$,
$H$ and the largest denominator of the search, filtered by $\mathrm{noY}$ and the search region; and
the coverage under \verb|RATPOINTS_NO_CHECK|
(\lname{Ratpoints.mem\_of\_isPoint\_of\_findPoints\_eq\_ok\_of\_noCheck}): every normalized point of
height at most $H$ in the search region gives the survivor $(x, 0, z)$. The range statement under
\verb|RATPOINTS_NO_CHECK|
(\lname{Ratpoints.exists\_eq\_of\_mem\_of\_findPoints\_eq\_ok\_of\_noCheck}) adds that every
survivor lies in the search region, without the assumption. The refusal theorems
(Theorem~\ref{thm:refusals}) add the refusal of the search intervals, from
Lemmas~\ref{lem:searchIntervals}, \ref{lem:reversalRegion} and~\ref{lem:sturmRegion}: with the Sturm
step on, $\mathrm{findPoints}$ refuses exactly when $c$ has fewer than two or more than
$\mathrm{maxDegree} + 1$ entries, $H = 0$ or $H \ge 2^{63}$, $\Phi_c$ is not squarefree, or the
caller's intervals are not well formed; with it off, exactly when $c$ has fewer than two or more
than $\mathrm{maxDegree} + 1$ entries, $H = 0$ or $H \ge 2^{63}$, $d + 2 \le N$, the reversal is
allowed, the search region is the whole box, $d$ is even and $c_0 = c_1 = 0$, or the caller's
intervals are not well formed.

\begin{theorem}[The Sturm step changes only the speed]\label{thm:sturmSpeed}
  \uses{thm:order, thm:refusals, def:findPoints}
  Proposed name: \lname{Ratpoints.findPoints\_sturm\_neg\_eq\_ok}. Assuming Sturm's theorem when the
  Sturm step is on, for well-formed parameters and options with the exact check on: if
  $\mathrm{findPoints}(c, H, P, o)$ returns $L$, then the search with the Sturm step off returns $L$
  too.
\end{theorem}

\begin{proof}
  If the search with the step off refused, so would the search with $o$, its refusals being among
  those with the step on (Theorem~\ref{thm:refusals}). Otherwise both lists equal, by the order
  theorem (Theorem~\ref{thm:order}), the points in the order of ratpoints filtered by $\mathrm{noY}$
  and the search region of the caller, and that right-hand side does not read the depth.
\end{proof}

This is the promise of Chapter~\ref{chap:setup}, that the depth of the Sturm step changes only the
speed, now a theorem, with the exact check on; under \verb|RATPOINTS_NO_CHECK| the survivors depend
on the depth, as in the C program.

\section{What comparator certifies}\label{sec:certifiedH}

\verb|comparator/Challenge.lean| restates, besides the specification that it restated before, the
definitions of \verb|Ratpoints/Spec/Region.lean| and \verb|Ratpoints/Spec/Sturm.lean| with their
decidability instances, verbatim, and imports exactly the Mathlib modules among the public imports
of the files of the specification: those of the file of the region with H.5 and those of the file
of Sturm's theorem with H.11; from the model it takes, besides what it took before, the doubles
($\mathrm{Double}$), whose bits the statements read through the search intervals of the options. The
twelve statements, for all $c$, $H$, parameters $P$ and options $o$, with $D$ the values and $B_D$
the bits of the intervals of $o$, the search region that of $D$, $o.\mathrm{bLow}$ and
$o.\mathrm{bHigh}$, and $h_S$ the hypothesis $0 \le o.\mathrm{sturm} \to \mathrm{SturmTheorem}$:
\begin{enumerate}
  \item a triple lies in $\mathrm{bruteForce}(c, H)$ exactly when it is a normalized point of height
    at most $H$ (\verb|challenge_mem_bruteForce_iff|, unchanged);
  \item $\mathrm{bruteForce}(c, H)$ has no duplicates (\verb|challenge_nodup_bruteForce|,
    unchanged);
  \item under $h_S$, for $P$ well formed, $\mathrm{noCheck}$ false and
    $\mathrm{findPoints}(c, H, P, o) = L$: $p \in L$ exactly when $p$ is a normalized point of
    height at most $H$, with $y \ge 0$ if $\mathrm{noY}$ is set, in the search region
    (\verb|challenge_mem_iff_isPoint_of_findPoints_eq_ok|, with the search region from H.5 and with
    $h_S$ from H.11);
  \item for $P$ well formed and $\mathrm{findPoints}(c, H, P, o) = L$, $L$ has no duplicates
    (\verb|challenge_nodup_of_findPoints_eq_ok|, its text unchanged, now over the search region
    too);
  \item with the Sturm step on, $\mathrm{findPoints}$ refuses exactly when $c$ has fewer than two or
    more than $\mathrm{maxDegree} + 1$ entries, $H = 0$, $H \ge 2^{63}$, $\Phi_c$ is not squarefree
    or $B_D$ is not well formed (\verb|challenge_exists_findPoints_eq_error_iff_of_sturm_nonneg|,
    with the search intervals, H.5);
  \item with the Sturm step off, exactly when $c$ has fewer than two or more than
    $\mathrm{maxDegree} + 1$ entries, $H = 0$, $H \ge 2^{63}$, $d + 2 \le N$, or the reversal is
    allowed, the search region is the whole box, $d$ is even and $c_0 = c_1 = 0$, or $B_D$ is not
    well formed (\verb|challenge_exists_findPoints_eq_error_iff_of_sturm_neg|, with the search
    region, H.5);
  \item the parameters of every set of options of the program are well formed
    (\verb|challenge_wf_sieveParams|, unchanged);
  \item under $h_S$, for $P$ well formed, $\mathrm{noCheck}$ false and
    $\mathrm{findPoints}(c, H, P, o) = L$: $L$ is the list of the points in the order of ratpoints
    for $\mathrm{maxDivisors}$, the switch of the reversal of $o$ and the whole box, $c$, $H$ and
    the largest denominator $\mathrm{denominatorBound}(o.\mathrm{bHigh}, H)$, filtered by
    $\neg\mathrm{noY} \lor y \ge 0$ and the search region
    (\verb|challenge_eq_orderedPoints_of_findPoints_eq_ok|, with the search region from H.5 and with
    $h_S$ from H.11);
  \item[9.] for $P$ well formed, $\mathrm{noCheck}$ set and $\mathrm{findPoints}(c, H, P, o) = L$:
    every member of $L$ is $(x, 0, z)$ with $\gcd(x, z) = 1$, $z > 0$ or ($z = 0$ and $x = 1$),
    $|x| \le H$ and $z \le H$, in the search region
    (\verb|challenge_exists_eq_of_mem_of_findPoints_eq_ok_of_noCheck|, strengthened by the search
    region, H.5);
  \item[10.] under $h_S$, for $P$ well formed, $\mathrm{noCheck}$ set and
    $\mathrm{findPoints}(c, H, P, o) = L$: for every normalized point $(x, y, z)$ of height at most
    $H$ in the search region, $(x, 0, z) \in L$
    (\verb|challenge_mem_of_isPoint_of_findPoints_eq_ok_of_noCheck|, with the search region from H.5
    and with $h_S$ from H.11);
  \item[11.] for $P$ well formed, $\mathrm{noCheck}$ set, $\mathrm{findPoints}(c, H, P, o) = L$ and
    the search with $\mathrm{noCheck}$ cleared and $\mathrm{noY}$ set returning $L'$: $L'$ with $y$
    replaced by $0$ is a sublist of $L$
    (\verb|challenge_map_sublist_of_findPoints_eq_ok_of_noCheck|, unchanged);
  \item[12.] for $P$ well formed and every callback with its data, $\mathrm{findPointsWith}$ with
    them gives, mapped to its data and count, what the delivery of the list of
    $\mathrm{findPoints}(c, H, P, o)$ to the callback gives, and the same refusal, with the same
    error, exactly where $\mathrm{findPoints}$ refuses
    (\verb|challenge_map_findPointsWith_eq_map_findPoints|, unchanged).
\end{enumerate}
$\mathrm{SturmTheorem}$ is a hypothesis of statements 3, 8 and~10, which one reads with the
specification. It holds trivially when the depth is negative, so for a search with the Sturm step
off, the program's \verb|-s|, the three statements hold without it. For the default depth
$\mathrm{defaultSturm}$ of the options and of the program without \verb|-s| or \verb|-S|
(Definition~\ref{def:Options}), they assume Sturm's theorem, the textbook theorem of
Definition~\ref{def:sturmTheorem} that Mathlib does not have yet. With the exact check on,
Theorem~\ref{thm:sturmSpeed} gives under it the same list as the search with the step off. The
theorems of the library behind the twelve statements have guards on their axioms, which are the
standard ones. Statement 4 needs no such hypothesis, since the list that the search delivers has no
duplicates whatever region it uses; statement 9 neither, since the intervals that the Sturm step
leaves lie in the caller's; and statement 11 neither, since both searches use the same intervals.

\section{The program}\label{sec:regionCli}

The program reads the ends of its intervals as \verb|main.c| reads them, in a new model file
\verb|Ratpoints/Model/CliNumbers.lean|, which also takes the readers of the numbers of the command
line from \verb|Ratpoints/Model/Cli.lean| (Definition~\ref{def:cliNumbers}): those of the integers,
of the one word of an argument and of a \verb|long| unchanged, while \verb|Ratpoints.Cli.isDouble|,
which only says whether an argument is a real number, gives way to \lname{Ratpoints.Cli.scanDouble?}
of Definition~\ref{def:cliEnds}, which also gives the double that \verb|scan_double| reads.

\begin{definition}[The ends of the command line]\label{def:cliEnds}
  \uses{def:cliNumbers, def:Double}
  Proposed names: \lname{Ratpoints.Cli.isInfWord}, \lname{Ratpoints.Cli.digitValue},
  \lname{Ratpoints.Cli.natOfDigitsIn}, \lname{Ratpoints.Cli.mantissaDigits},
  \lname{Ratpoints.Cli.floatParts}, \lname{Ratpoints.Cli.exponentDigits},
  \lname{Ratpoints.Cli.nearestBinary}, \lname{Ratpoints.Cli.nearestDecimal},
  \lname{Ratpoints.Cli.strtod}, \lname{Ratpoints.Cli.scanDouble?}, \lname{Ratpoints.Cli.ExactEnd},
  \lname{Ratpoints.Cli.outward}, \lname{Ratpoints.Cli.scanEnd}, \lname{Ratpoints.Cli.heightEnd}. The
  parts of a floating constant as \verb|scan_end| and \verb|strtod| read them: the sign, whether it
  is hexadecimal, the digits of the mantissa without the point and the number of them after it, as
  the loop of \verb|scan_end| reads them (\csource{main.c}{822}{834}), and the exponent; the double
  that \verb|strtod| of glibc reads from a word, in the C locale: $\pm\infty$ for an infinity, and
  otherwise the double nearest to the exact value $\pm m\,10^{e - k}$, or $\pm m\,2^{e - 4k}$ for a
  hexadecimal word, with ties to even (Definition~\ref{def:Double}), a zero of the sign of the word
  for a value of at most $2^{-1075}$ and an infinity of its sign for one of at least $2^{1024} -
  2^{970}$, the size of the value decided by bounds before a power is expanded, so that a word with
  a huge exponent costs no more than another; the double that \verb|scan_double| reads, a word with
  only white space around it, read by \verb|strtod|, none for anything else, among them a NaN
  (\csource{main.c}{746}{752}, from Definition~\ref{def:cliNumbers}). The end \verb|exact_end|,
  $\mathrm{inf} = -1$, $1$ or $0$ for the rational number $q$ (\csource{main.c}{71}{71}); the double
  that \verb|outward| gives for an end $v$ from the double $x$ nearest to it: while $x$ is finite
  and lies above $v$ (a lower end) or below it (an upper end), the next double outward
  (\csource{main.c}{761}{763}, \csource{main.c}{766}{773}, \csource{main.c}{775}{776}); the end that
  \verb|scan_end| reads and its double, or none when the word is not a number: for an infinite
  double of \verb|strtod|, the infinite end of its sign and the double itself; otherwise the number
  of the word, $0$ when its digits are $0$, $\pm 2^{-1075}$ when \verb|strtod| reads $0$ although
  they are not, and else $m\,2^{e - 4k}$ or $m\,10^{e - k}$ with the exponent $e$ accumulated as
  \verb|scan_end| accumulates it, up to \verb|LONG_MAX / 20|, with its sign, and the double of
  \verb|outward| from the double of \verb|strtod| towards it (\csource{main.c}{798}{806},
  \csource{main.c}{808}{819}, \csource{main.c}{835}{847}, \csource{main.c}{849}{849},
  \csource{main.c}{851}{865}, \csource{main.c}{868}{873}, \csource{main.c}{875}{876},
  \csource{main.c}{878}{880}); and the end that \verb|height_end| gives for a missing first
  \verb|-l| (negated) and a missing last \verb|-u|, the double of $H$ moved up to the next double
  when it is below $2^{63}$ and below $H$, as the default interval of the library has it
  (\csource{main.c}{299}{305}).
\end{definition}

\begin{definition}[The test of the region and the join]\label{def:cliRegionTest}
  \uses{def:cliEnds, def:Double}
  Proposed names: \lname{Ratpoints.Cli.cmpEnd}, \lname{Ratpoints.Cli.cmpEndX},
  \lname{Ratpoints.Cli.inRegion}, \lname{Ratpoints.Cli.joinIntervals}, \lname{Ratpoints.Cli.printF}.
  The sign of one end minus another, as \verb|cmp_end| computes it,
  $\mathrm{inf}_e - \mathrm{inf}_f$ when one of them is infinite, so that two infinities of one sign
  are equal (\csource{main.c}{309}{313}); the sign of an end minus $a/b$, $b > 0$, as
  \verb|cmp_end_x| computes it (\csource{main.c}{116}{120}); whether $a/b$ lies in one of the $n$
  intervals of the command line, by the binary search of \verb|in_region| for the first interval
  that does not end below $a/b$, which must hold it (\csource{main.c}{126}{137}); the intervals that
  \verb|join_intervals| gives the library: from the second on, each interval is merged into the last
  interval formed when its lower double is at most the upper double of that interval (it becomes the
  smallest interval holding both: the lesser lower double and the greater upper double), and
  otherwise starts a new one, so that intervals whose doubles meet or overlap become one
  (\csource{main.c}{891}{905}); and what \verb|printf| of glibc prints for a double under \verb|%f|:
  the exact value rounded to six places, ties to even, in full however large it is, after a \verb|-|
  whenever the sign bit is set, also for $-0.0$ and a negative value that rounds to $0$, \verb|inf|
  and \verb|nan| with their sign. \verb|printF| models the C library and holds no C range.
\end{definition}

\begin{definition}[The intervals of the command line]\label{def:cliIntervals}
  \uses{def:cliEnds, def:cliRegionTest, def:cliFailure, def:maxDegree, def:Options}
  Proposed names: \lname{Ratpoints.Cli.Intervals}, \lname{Ratpoints.Cli.Intervals.lower},
  \lname{Ratpoints.Cli.Intervals.upper}, \lname{Ratpoints.Cli.Intervals.complete},
  \lname{Ratpoints.Cli.OptionArg.bound}, \lname{Ratpoints.Cli.Options.store}. The state of the
  reading of the search intervals: whether the last end read is a lower one, the number of the
  intervals read, and the arrays of the doubles and of the exact ends of $2\,\mathrm{maxDegree}$
  entries (\csource{main.c}{66}{66}, \csource{main.c}{73}{73}, \csource{main.c}{317}{317},
  \csource{main.c}{319}{319}). Reading a lower end (\verb|-l|): error~7 when the last end read is a
  lower one too (\verb|-l -l|), when $\mathrm{maxDegree}$ intervals have been read, when the word is
  not a number and when the end is not above the upper end before it, compared as numbers; otherwise
  the lower end of the next interval (\csource{main.c}{401}{413}). Reading an upper end (\verb|-u|):
  when the last end read is not a lower one, the first interval begins at $-\infty$ with the double
  $-h$ of the height end, and a later one fails with error~7 (\verb|-u -u|); error~7 when the word
  is not a number; the lower double $-\infty$ when the interval has no lower end and its upper
  double lies below $-h$; error~7 when the lower end is above the upper end; otherwise the interval
  is complete (\csource{main.c}{417}{436}). After the options, a pending lower end gets the upper
  end $+\infty$ with the double $h$, or $+\infty$ when $h$ lies below the lower double
  (\csource{main.c}{669}{669}, \csource{main.c}{673}{678}). The kind of argument of \verb|-l| and
  \verb|-u|, an end, a missing or malformed one failing with error~7; and the options of the search
  with the field that an option sets: the reversal for \verb|-k| and \verb|-K|, the depth for
  \verb|-s| and \verb|-S|, the least and the largest denominator for \verb|-dl| and \verb|-du|, and
  the options of the output for \verb|-x -X -y -Y|.
\end{definition}

The command line (Definition~\ref{def:cliInput}) reads \verb|-l| and \verb|-u| into the intervals,
\verb|-dl| and \verb|-du| into the options of the search, starts from the options with the height
bound as the largest denominator, completes a pending interval after the options, and gives the
search intervals of the options as the doubles read and the exact ends in the settings of the output
(\lname{Ratpoints.Cli.Output}, which gains them); its C ranges gain the cases of the four options
and the completion (\csource{main.c}{393}{400}, \csource{main.c}{414}{416},
\csource{main.c}{650}{668}, \csource{main.c}{333}{336}, \csource{main.c}{380}{380},
\csource{main.c}{679}{679}). The failures (Definition~\ref{def:cliFailure}) gain error~7 for the
intervals and error~9 for the refusal $\mathrm{badArgs}$, with their messages
(\csource{main.c}{1036}{1048} in place of the links of lines 1036-1043 and 1047, the exclusion of
line 1048 going), and lose the option not supported yet with its exit status~12, as
\verb|-l -u -dl -du|, the last options of ratpoints that the program did not read, are now read
(\verb|-t|, the number of threads, as before checked and ignored). The data of the callback
(Definition~\ref{def:cliProcess}) gain the number of the intervals of the command line, before the
join, and their exact ends (\csource{main.c}{246}{247}), and the callback passes over an affine
point whose $x/z$ lies in none of the exact intervals, with the weight $0$ and the flag clear
(\csource{main.c}{150}{152}). The report (Definition~\ref{def:cliReport}) gains message~3, the
search region of the command line before the join, printed only when the command line gives
intervals, after the height bound, or after the curve equation under \verb|-v|
(\csource{main.c}{224}{224}), and message~8, the search intervals that the search gives back, after
message~7 and not under \verb|-v|, each interval printed as \verb|[%f, %f]| and the intervals joined
by \verb| U | (\lname{Ratpoints.Cli.showIntervals}, which the node of the report names; the texts
\csource{main.c}{979}{991} and \csource{main.c}{999}{1033} in place of the links of lines 979-983,
999-1021 and 1031-1033). The program (Definition~\ref{def:cliRun}) builds the data of the callback
from the intervals before the join and runs the search on the joined intervals
(\csource{main.c}{248}{248}).

\begin{lemma}[The numbers of the command line]\label{lem:cliNumbers}
  \uses{def:cliEnds, def:cliRegionTest, lem:rounding, lem:Double, def:regionSpec}
  Proposed names: \lname{Ratpoints.Cli.ExactEnd.value}, \lname{Ratpoints.Cli.InExactRegion},
  \lname{Ratpoints.Cli.isNaN\_strtod}, \lname{Ratpoints.Cli.lt\_nextUp\_strtod},
  \lname{Ratpoints.Cli.wordValue}, \lname{Ratpoints.Cli.scanEndValue},
  \lname{Ratpoints.Cli.floatParts\_spec}, \lname{Ratpoints.Cli.mantissaDigits\_spec},
  \lname{Ratpoints.Cli.OutwardInvariant}, \lname{Ratpoints.Cli.outward\_spec},
  \lname{Ratpoints.Cli.run\_outward}, \lname{Ratpoints.Cli.run\_scanEnd\_props},
  \lname{Ratpoints.Cli.run\_scanEnd\_isMaxBelow}, \lname{Ratpoints.Cli.heightEnd\_props},
  \lname{Ratpoints.Cli.cmpEnd\_neg\_iff}, \lname{Ratpoints.Cli.cmpEnd\_pos\_iff},
  \lname{Ratpoints.Cli.cmpEndX\_neg\_iff}, \lname{Ratpoints.Cli.cmpEndX\_nonpos\_iff},
  \lname{Ratpoints.Cli.ChainEnds}, \lname{Ratpoints.Cli.InRegionInvariant},
  \lname{Ratpoints.Cli.inRegion\_spec}, \lname{Ratpoints.Cli.run\_inRegion},
  \lname{Ratpoints.Cli.NoNaN}, \lname{Ratpoints.Cli.JoinsWell}, \lname{Ratpoints.Cli.JoinInvariant},
  \lname{Ratpoints.Cli.joinIntervals\_spec}, \lname{Ratpoints.Cli.run\_joinIntervals},
  \lname{Char.toLower\_eq\_x\_iff}, \lname{WithBotTop.bot\_lt\_coe},
  \lname{WithBotTop.coe\_lt\_top}. The value of an exact end, $\bot$, $\top$ or its rational number,
  and the predicate that a rational number lies in one of the $n$ exact intervals; an element of
  \verb|WithBotTop ℚ| other than $\bot$ and $\top$ lies strictly between them (the mirror
  \verb|Order/WithBotTop|), and the characters whose lower case is \verb|x| are \verb|x| and
  \verb|X| (for the prefix \verb|0x| of a hexadecimal word). \verb|strtod| gives no NaN, and the
  value of a floating constant from which it reads a finite double other than a zero lies below the
  next double above that double: the double lies above the value or less than one step below it, the
  half of the faithful rounding of Lemma~\ref{lem:rounding} that the lower ends need. \verb|outward|
  moves a double to the side of the end that the direction asks, and for a lower end, from a double
  whose next double lies above the end, it lands on the largest double at most the end; so the
  double that \verb|scan_end| gives for a lower end of a word of fewer than $2^{56}$ characters is
  the largest double at most the end, and for an upper end a double at least the end. The height end
  is finite and at least $H$, for $1 \le H < 2^{63}$. The signs of \verb|cmp_end| and
  \verb|cmp_end_x| are those of the differences of the values; the binary search decides membership
  in increasing and disjoint exact intervals; and \verb|join_intervals| gives none exactly when
  there are none, its intervals hold the values of the intervals it joins, and for intervals without
  a NaN, with $l \le u$, whose lower ends increase from the second on, its intervals have $l \le u$
  and are increasing and strictly apart.
\end{lemma}

\begin{proof}
  The parts of the word and its value, with the capped exponent equal to the true one below
  $2^{56}$ characters; the rounding of \verb|strtod| by Lemma~\ref{lem:rounding}; the loops of
  \verb|outward|, \verb|in_region| and \verb|join_intervals| by named invariants with the program
  logic of Lean core.
\end{proof}

From the double that \verb|strtod| reads, nearest to the value of the word, \verb|outward| lands on
the largest double at most a lower end and on the smallest double at least an upper end, as in the C
program; the proofs show the first for a faithful rounding, and of the second only that the double
is at least the end, which is all that the theorems need; the differential tests compare the doubles
of the ends with the C program, in messages~3 and~8.

\begin{theorem}[The search region of the command line]\label{thm:cliRegion}
  \uses{def:cliIntervals, def:cliInput, lem:cliNumbers, def:cliProcess, thm:findPoints, thm:cliRun,
    def:regionSpec, lem:searchIntervals}
  Proposed names: \lname{Ratpoints.Cli.IntervalOf}, \lname{Ratpoints.Cli.Intervals.ReadsWF},
  \lname{Ratpoints.Cli.Intervals.LowerMax}, \lname{Ratpoints.Cli.checkOptions\_preserves},
  \lname{Ratpoints.Cli.input\_eq\_ok}, \lname{Ratpoints.Cli.reads\_of\_input\_eq\_ok},
  \lname{Ratpoints.Cli.joinsWell\_of\_reads}, \lname{Ratpoints.Cli.intervalsWF\_take\_of\_reads},
  \lname{Ratpoints.Cli.keeps\_ofOutput\_iff}, \lname{Ratpoints.Cli.inRegion\_joinIntervals\_iff},
  \lname{Ratpoints.Cli.inRegion\_iff\_of\_input\_eq\_ok},
  \lname{Ratpoints.Cli.intervalsWF\_joinIntervals\_of\_input\_eq\_ok},
  \lname{Ratpoints.Cli.inIntervals\_joinIntervals\_of\_inExactRegion},
  \lname{Ratpoints.Cli.mem\_filter\_keeps\_iff\_of\_input\_eq\_ok},
  \lname{Ratpoints.zero\_lt\_maxDegree}, \lname{Array.set!\_set!},
  \lname{Array.getElem!\_set!\_eq\_or}, \lname{Array.getElem!\_take\_toList},
  \lname{Array.mem\_take\_toList\_of\_lt}, \lname{Array.exists\_eq\_of\_mem\_take\_toList},
  \lname{Except.bind\_eq\_ok}. For a command line that reads as $c$, $H$, the options $s$, $o$ and
  the settings $u$: the exact ends of its intervals are increasing and disjoint, by the checks of
  error~7, so that the test of the callback is membership in the exact intervals; the joined
  intervals hold every $x$ of the exact region with $|x| \le H$, their ends being the doubles
  outward of the exact ends, or $\mp h$ for a missing first lower or last upper end; when every word
  has fewer than $2^{56}$ characters, the joined intervals are well formed, so that error~9 is
  unreachable from the command line; and, assuming Sturm's theorem when the Sturm step is on,
  without \verb|-x|, if $\mathrm{findPoints}$ on the joined intervals returns $L$, the points of $L$
  that the callback keeps (Theorem~\ref{thm:cliRun}) are exactly the normalized points of height at
  most $H$ with $y \ge 0$ under \verb|-y|: the points at infinity unless \verb|-i|, and the affine
  points with $z$ in the range of \verb|-dl| and \verb|-du| and $x/z$ in the exact region of
  \verb|-l| and \verb|-u| when there is one. Under \verb|-q| the program prints these points
  (Theorem~\ref{thm:cliRun}), and under \verb|-1| the first of them.
\end{theorem}

\begin{proof}
  An invariant of the reading of the intervals, preserved by each option and by the completion: the
  exact ends increase and are disjoint, each complete interval of doubles holds its exact interval
  as far as the search goes, and the doubles of the lower ends after the first are the largest
  doubles at most their exact ends (Lemma~\ref{lem:cliNumbers}), which increase, so that the join
  makes intervals that are strictly apart. The last statement is membership
  (Theorem~\ref{thm:findPoints}) for the joined intervals with the test of the callback: a point in
  the joined intervals with $|x| \le H$ is kept exactly when it lies in the exact region. The bound
  $|x| \le H$ is needed, since the doubles $\mp h$ of a missing end stand for $\mp\infty$; the bound
  on the length of the words too, since \verb|scan_end| stops accumulating the digits of an exponent
  at \verb|LONG_MAX / 20|, so that a word of about $4.6 \cdot 10^{17}$ digits, far beyond what a
  command line holds, could give an exact end other than the value that \verb|strtod| rounds.
\end{proof}

The theorem on the program (Theorem~\ref{thm:cliRun}) keeps its statements with the joined
intervals: the counted points are those of the list that the callback keeps, counting them
(\lname{Ratpoints.Cli.Data.keeps}: a point at infinity unless \verb|-i|, and an affine point when
the command line gives no intervals or its $x/z$ lies in one of them by \verb|in_region|; the
theorem's node names it next to \lname{Ratpoints.Cli.countedPoints}, which takes this test in place
of the switch of the points at infinity), and the program prints them, under \verb|-q| only them and
the strings, under \verb|-1| the first; with Theorem~\ref{thm:cliRegion} it prints exactly the
points of the search region of its command line. Theorem~\ref{thm:cliRun} takes no hypothesis of
Sturm's theorem, since it is stated over the list of $\mathrm{findPoints}$. The last statement of
Theorem~\ref{thm:cliRegion}, and with it the statement that the program prints exactly the points of
the search region of its command line, takes the hypothesis from H.11 on, as membership does: it is
needed for the command lines without \verb|-s|, and under \verb|-s| the theorem holds without it.

\section{The differential tests}\label{sec:regionTests}

The harness compares the search region and the Sturm step with the C program at the pin:
\begin{itemize}
  \item every run of \verb|test4.sh| with \verb|-l|, \verb|-u|, \verb|-dl| or \verb|-du|, $165$ of
    them, and the one of \verb|test3.sh| are compared, by their points in their order, their
    reports and their errors; the option status~12 and its check go;
  \item a search above the largest height of the default suites, $30000$, is compared when it has
    a search region, the model lists at most $10^5$ denominators for it (by a bound that mirrors the
    divisor set-up of the model), and it sifts at most $2 \cdot 10^{10}$ numerators, from the widths
    of the intervals and, under \verb|-1|, from the first point of the expected output, where the
    search stops; such searches run with a timeout and a cap of $4$~GB on the memory: the searches
    of \verb|test4.sh| at large heights in a narrow region, among them three under \verb|-1| at the
    heights $10^{12}$ and $2^{63} - 1$ whose stop lies in the sieve;
  \item messages~3 and~8 are compared byte for byte in every report: message~3 from H.7, message~8
    from H.11 (until then dropped from the report of the C program, as the model does not narrow the
    search intervals before H.11); only message~4, on the moduli used, stays dropped (stage~I);
  \item under \verb|-x| the lines are the survivors of the sieve over the intervals that the Sturm
    step keeps. From H.7 on, the weaker check runs on every run under \verb|-x| without a format or
    a string, now with the exact region of the command line and the range of the denominators, as
    \verb|process| of \verb|main.c| tests them; it is the only check where the pin changes the
    output of \verb|ratpoints-lean|, or where the C program, as its report under \verb|-v -z| says,
    sieves with other moduli under the pin (a subset of the $30$ primes, or composite moduli such as
    $9$ in place of $3$, or $55$ and $63$ in place of $3$, $5$, $7$ and $11$), which stage~I models,
    and the reason is counted. Otherwise the outputs are also compared exactly with the C program on
    the command line with the moduli pinned (\verb|-p 30 -N 30| appended), on which it sieves with
    the first $30$ odd primes in its first two stages and no third stage: with \verb|-s| appended as
    well until H.11, as in stage~G, and without \verb|-s| from H.11 on. The moduli are always read
    from the command line that is compared, since the choice of the moduli of the C program reads
    the search intervals;
  \item a new default suite \verb|regions| runs $400$ random search regions against the C program:
    intervals with ends of every kind (decimals, hexadecimal constants, ends that are no doubles,
    infinities, numbers beyond the range of the doubles, tiny ones, $-0$, intervals that touch or
    whose doubles meet, a missing first \verb|-l| or last \verb|-u|), ranges of the denominators of
    every shape and random options, at small heights, and in one run in ten at a large height with a
    narrow region, its ends exact doubles at times, also around points of large height and at the
    ends of the range $[-H, H]$ of the numerators, at times under \verb|-x| with the moduli pinned;
  \item a new default suite \verb|depths| runs $300$ random depths of the Sturm step, \verb|-S 0| to
    \verb|-S 60| and \verb|-s|, with \verb|-q| or the report (message~8), at times with intervals,
    \verb|-k| or \verb|-x|;
  \item the errors suite has the cases of the search region: errors~7 and~6, the join, the test of
    the callback, the suppression of the reversal, and the order with a range of the denominators;
    the fuzz draws ends for \verb|-l| and \verb|-u| and \verb|long|s for \verb|-dl| and \verb|-du|.
\end{itemize}
In the prototype the default suites compare $4530$ times with none failed: \verb|test4| $565$,
\verb|test3| $16$, \verb|rptest| $1206$, \verb|random| $300$, \verb|outputs| $299$, the errors suite
$360$, the fuzz $800$, the parameters $272$, the high suite $12$, \verb|regions| $400$ and
\verb|depths| $300$; under \verb|-x| $142$ runs are compared exactly and $7$ by the weaker check
alone, all for the C program's choice of other moduli under the pin: in six of them a subset of the
$30$ primes, in one (of \verb|test4.sh|, with \verb|-n 15 -C 0|) the moduli $55$ and $63$; and
\verb|test4high| compares $13$, none failed. Every point list is compared in its order.

The \verb|#guard| examples of the model check the doubles against a driver of gcc and glibc (the
conversions of \verb|long|s, the comparisons, \verb|nextafter|, the negation, \verb|fabs|,
\verb|frexp|, \verb|ldexp| and the truncations), the ends, the scaled ends and the forms against the
lines of the C code run verbatim, the bits of intervals against the build of ratpoints under
\verb|DEBUG|, \verb|sift| over regions against the plain sieve, with the store of the rows built
lazily, the Sturm step against a driver of \verb|_ratpoints_compute_sturm| (few real intervals, a
form negative everywhere, a form that is not squarefree, dyadic roots met by midpoints, the depths
$0$ to $55$, $-0.0$ kept, an empty intersection, coefficients of $70$ bits) and the bound $1 +
\lfloor d/2 \rfloor$ of its array, the search with a callback with its outcome and its refusal
$\mathrm{badArgs}$ against \verb|find_points| of the library, and, for the program, \verb|strtod|
and \verb|%f| against glibc, the ends of \verb|scan_end|, the height end, \verb|in_region|, the
join, and whole runs with regions, errors~6 and~7, $100$ intervals accepted and $101$ refused,
against the C program.

\section{Deviations}\label{sec:regionDeviations}

Each component differs from the C code in a stated way, recorded in its \verb|Deviation:| paragraph:
\begin{itemize}
  \item the doubles as a structure over their bits, with functions of the model for the operations
    of the C language and of the C library on them, which model the platform (IEEE~754 binary64,
    rounding to nearest, gcc and glibc) and hold no C range; \verb|strtod| and \verb|printf| under
    \verb|%f| as what glibc computes, not how it computes it: the nearest double to the exact value
    of the word, its size decided by bounds before a power is expanded, and the exact decimal
    expansion rounded to six places;
  \item the 128-bit product, sum and shift of the scaled end on two $64$-bit words, the high word of
    the product by $\mathrm{mulHi}$, for the \verb|unsigned __int128| of the arm that the builds of
    the model compile, which Lean core does not have, the other arm excluded; the \verb|long| $b$
    converted to a word, as the C code converts it to an \verb|unsigned long|;
  \item natural numbers for the \verb|int| exponents and shifts of the ends and the forms, which are
    from $0$ to $127$ and from $0$ to $63$, and Booleans for their \verb|int| flags; the mantissa
    and the exponent of an end, the end, its form and the bounds of the run returned, the form as an
    option, for the fields, the pointers and the \verb|int| through which the C code writes them,
    and the array of the bounds that \verb|run_setup| allocates; the \verb|&&| of
    \verb|exact_bounds|, which makes the form of the upper end only when the lower end has one, as a
    bind of options; \verb|Int64| for the \verb|long|s, with the division of \verb|Int64|, which
    truncates as C's does, and \verb|>>>| on \verb|Int64| the arithmetic shift, for the \verb|>>| of
    a negative \verb|long|, which gcc defines so, the arithmetic of \verb|Int64| wrapping where that
    of the \verb|long|s overflows, which it does not here, by the bound of the forms; the code and
    the bits of an interval returned as a triple, the bits $(0, 0)$ with the codes $1$ and $2$, for
    which the C code writes nothing; no output under \verb|DEBUG|;
  \item the bounds of the search intervals an array passed from the run to each denominator, for the
    struct \verb|rp_intervals| that \verb|sift| reaches through the field \verb|iv| of the sieving
    thread, which \verb|worker_init| clears and \verb|run_block| sets for each block;
  \item the options of the search a structure: a list of pairs of doubles for the array
    \verb|domain| with its number \verb|num_inter|, the empty list for $\mathrm{num\_inter} = 0$, so
    that no number of intervals is negative and the room in the array for the intervals of the Sturm
    step is no concern of the caller, and integers for the \verb|long|s \verb|b_low| and
    \verb|b_high|; of the fields that share their lines in \verb|ratpoints.h|, \verb|sp1| and
    \verb|sp2| are the parameters', \verb|sp3| is stage~I's, \verb|num_primes| and
    \verb|max_forbidden| are the options of the program's;
  \item the normalized options returned, for the fields that the C code normalizes in place and
    \verb|find_points_work| restores; no clamp of the number of the intervals, which a list cannot
    make negative; the bit \verb|RATPOINTS_NO_REVERSE_AUTO|, which the C code sets in the flags and
    the reversal check tests with \verb|RATPOINTS_NO_REVERSE|, as the switch of the reversal
    cleared;
  \item the search intervals returned as a list, none for \verb|RATPOINTS_BAD_ARGS|, for the array
    that the C code sets to the default interval in place; the test of a null pointer excluded;
  \item the Sturm step of the search a function of its own, its result a refusal, none, or the
    intervals with the flag \verb|sturm_empty|, for the codes that the C code returns at lines 4879
    and 4889, the intervals that it leaves in \verb|domain| and its local flag; the square test made
    on the polynomial searched, which gives the C code's \verb|lcfsq| for an even degree, the only
    one that reads it; at the return of line 4889 the outcome with the data, the count $0$, the flag
    of the reversal and no interval, the flag being in fact clear there; no failure for lack of
    memory; no output under \verb|RATPOINTS_VERBOSE|;
  \item what a search gives back with the search intervals as a list, for \verb|args->domain| and
    \verb|args->num_inter|, which the search writes back and message~8 reads; the range, the
    intervals and the flag \verb|sturm_empty| arguments of the search and of the sieving, for the
    fields of the arguments and the local flag of \verb|find_points_work_1|; the test of the
    \verb|long|s, which the C code makes between the reversal check and the search intervals
    (\verb|find_points.c| lines 4814-4823), made by the search, after the search intervals and the
    Sturm step, where it is used;
  \item in the Sturm step, the rows as arrays and integers for the GMP integers, with
    \verb|Int.sign| for \verb|mpz_sgn| and no scratch integers; \verb|Int64| for the \verb|long|
    numerators and exponents, with the arithmetic shift, wrapping where the \verb|long|s overflow,
    which they do not up to the depth $55$; natural numbers for the \verb|long| counts, the depths
    and the depth of the step, which is from $0$ to $55$; the steps of Horner's rule counted up; the
    signs and the count at the midpoint, the midpoint, the kept interval and the body of the loop of
    the intersection functions of their own, their results returned, for the locals of
    \verb|iterate| and of \verb|_ratpoints_compute_sturm|; one constant for \verb|max| and
    \verb|min|;
  \item the bisection a recursive function, for the nested function \verb|iterate| of gcc, with the
    intervals kept as an array passed and returned, for the array \verb|ivlocal| of
    $1 + \lfloor d/2 \rfloor$ intervals on the stack and its pointer, which the C code reaches from
    the nested function, with no bound on their number, which the C code relies on and the proofs
    show under Sturm's theorem (Lemma~\ref{lem:sturmBuffers}), and the chain, its last index and the
    depth of the step as arguments; the doubles of the ends as bits;
  \item the intersection read from the caller's array, where the C code reads it from a copy on the
    heap and writes over the caller's array, and returned as a new array, with no bound on its
    length, which the manual asks the caller to provide for and the proofs show under Sturm's
    theorem;
  \item the Sturm step with its arguments for the fields of the arguments, its result an option of
    an array for the \verb|long| it returns and the intervals it writes, the table of the chain made
    by the chain of stage~D and the array of the intervals kept by the bisection; no \verb|mpz_init|
    and \verb|mpz_clear|, no copy on the heap and no failure for lack of memory;
  \item in the program, the reading of the intervals a structure for the two locals of
    \verb|read_input| and the two global arrays, zero and $0$ at the start, as \verb|mpq_init| sets
    the exact ends; the end and its double returned by \verb|scan_end|, or none, the end of an
    infinity with $q = 0$, which \verb|scan_end| leaves as it is and nothing reads; rational numbers
    of core for the GMP integers and rationals, without their initialization and release, and
    without the allocation of the digits and its error~10; the exact value of a double compared as a
    rational number in \verb|outward|; the signs of \verb|cmp_end| and \verb|cmp_end_x| as $-1$, $0$
    or $1$; lists and arrays for the arrays of \verb|main.c|, natural numbers for its \verb|long|
    indices; the NaN refused by the syntax of \verb|scan_double|, for its test \verb|isnan|; the
    largest denominator $-1$ of line 336, which line 380 replaces by the height bound, not set; the
    loops of messages~3 and~8 one function, \verb|%lf| being \verb|%f|.
\end{itemize}
The deviations that stage~H retires: ``every numerator from $-H$ to $H$'' of one denominator
(Chapter~\ref{chap:words}); the reversal check without \verb|RATPOINTS_NO_REVERSE_AUTO|, no error~9,
and of the Sturm step only its test for a repeated factor, with the survivors outside the region
that the step keeps under \verb|RATPOINTS_NO_CHECK| (Chapter~\ref{chap:setup}); no search intervals
and no range of the denominators other than $1$ to $H$ (Chapter~\ref{chap:denominators}); no search
intervals of the sieving thread (Chapter~\ref{chap:stages}); the form of the bits of the whole line
decided in lines of the search intervals, and the sieve specs made before the bits of the interval
(Chapter~\ref{chap:packing}); and the search intervals not written back, the survivors under
\verb|RATPOINTS_NO_CHECK| outside the region of the Sturm step, the return of line 4889 and
\verb|sturm_empty| not met, messages~3 and~8 left out, the test of the search region in
\verb|process|, and \verb|-l -u -dl -du| answered with the exit status~12
(Chapter~\ref{chap:output}).

The lines of this part of the C code that are neither modelled nor left to a later stage are
excluded in \verb|correspondence/ratpoints.toml|, each by the item named in parentheses: the
includes of the headers of the C library and of the header of the library in \verb|rp-private.h|
(\csource{rp-private.h}{30}{35}, \csource{rp-private.h}{60}{60}) and of \verb|gmp.h| in
\verb|ratpoints.h| (\csource{ratpoints.h}{31}{31}), H.1, and the includes of \verb|sturm.c|
(\csource{sturm.c}{29}{32}, H.8), other, since the model has no headers (\verb|float.h| and
\verb|math.h| are the headers that stage~H needs); the arm of \verb|exact_scale| for a build without
a 128-bit type or with \verb|RP_NO_INT128| defined (\csource{find_points.c}{3239}{3256}, H.2),
other, since the builds of the model compile the arm with \verb|unsigned __int128|; the copy of the
caller's intervals on the heap, the failure of its allocation with the clears of the table, and its
release (\csource{sturm.c}{274}{275}, \csource{sturm.c}{278}{284}, \csource{sturm.c}{324}{324},
H.8), memory; the prototype of \verb|_ratpoints_compute_sturm| (\csource{rp-private.h}{617}{617},
H.9), interface; and in \verb|main.c| the initialization and release of the GMP rational of
\verb|outward| and of the GMP integers of \verb|scan_end| and the allocation and release of its
digits (\csource{main.c}{765}{765}, \csource{main.c}{774}{774}, \csource{main.c}{820}{821},
\csource{main.c}{848}{848}, \csource{main.c}{850}{850}, \csource{main.c}{867}{867},
\csource{main.c}{874}{874}, \csource{main.c}{877}{877}, H.6), memory. The exclusion of the message
of error~9 (\verb|main.c:1048|) goes (H.5): the program gives that error for the refusal
$\mathrm{badArgs}$. After stage~H every line of \verb|sturm.c| is recorded; the lines of
\verb|main.c| left free are the defaults of the tuning options (lines 339-343 and 345-346) and
message~4 (lines 992-998), and those of \verb|find_points.c| the estimate of the work and the choice
of the moduli by the curve, the blocks and the threads; in \verb|ratpoints.h| the tuning macros
(lines 92-94 and 120-331), the tuning fields and the number of threads of \verb|ratpoints_args|
(lines 365-367 and 370), the numbers of the moduli used (line 378) and the private data of
\verb|ratpoints_args| (lines 384-400), and in \verb|rp-private.h| the macros of the threads (lines
76-84) and, in \verb|rp_worker|, the counters of the block with the decision of the correction
(lines 579-582) and the fields of the pool with the name of the type (lines 584-585), all of them
stage~I's.

\section{What changes in the earlier chapters}\label{sec:regionEarlier}

Each change is made by the item of Section~\ref{sec:regionFiles} named in parentheses, so that the
blueprint stays true between the commits; the new uses are those that the dependency audit asks for
once the names have moved as stated.
\begin{itemize}
  \item Chapter~\ref{chap:spec}: the paragraph after Definition~\ref{def:bruteForce} says that the
    order and the membership hold for the points in the search region (H.5); the section on what
    comparator certifies (Section~\ref{sec:comparator-ab}) states membership, the order and the
    coverage under \verb|RATPOINTS_NO_CHECK| with the search region (H.5) and the hypothesis of
    Sturm's theorem (H.11), the refusal theorems with the search intervals (H.5), the restated files
    of the region (H.5) and of Sturm's theorem (H.11) with their public imports, and
    $\mathrm{Double}$ among what the challenge takes from the model (H.5), the record with H.11; the
    list of Section~\ref{sec:cli}: the program prints messages~3 (H.7) and~8 (H.11), leaving out
    only message~4, reads \verb|-l -u -dl -du| (H.7), and the item on the options that it does not
    support and the exit status~12 goes (H.7); the readers of numbers
    (Definition~\ref{def:cliNumbers}) lose \verb|Ratpoints.Cli.isDouble|, whose range
    \csource{main.c}{746}{752} goes with \lname{Ratpoints.Cli.scanDouble?} to
    Definition~\ref{def:cliEnds}, and live in \verb|Ratpoints/Model/CliNumbers.lean| (H.6); the
    failures (Definition~\ref{def:cliFailure}) gain error~9 for the refusal $\mathrm{badArgs}$, with
    \csource{main.c}{1048}{1048}, their uses gaining Definition~\ref{def:searchIntervals} (H.5), and
    error~7, their link of the messages becoming \csource{main.c}{1036}{1048}, as they lose the
    option not supported yet with its status~12 (H.7); the command line
    (Definition~\ref{def:cliInput}) reads \verb|-l -u -dl -du| as in Section~\ref{sec:regionCli},
    its links \csource{main.c}{393}{400}, \csource{main.c}{414}{416}, \csource{main.c}{650}{668},
    \csource{main.c}{333}{336}, \csource{main.c}{380}{380} and \csource{main.c}{679}{679} in place
    of \csource{main.c}{393}{397} and \csource{main.c}{667}{668}, and uses
    Definition~\ref{def:cliIntervals} (H.7); the program (Definition~\ref{def:cliRun}) runs the
    search on the joined intervals, gains \csource{main.c}{248}{248}, prints the points of the
    search region in the order of Theorem~\ref{thm:order}, and uses
    Definition~\ref{def:cliRegionTest} (H.7); the section on the differential tests: its list of
    what the later stages added gains stage~H, with the suites \verb|regions| and \verb|depths| and
    the region runs of \verb|test4.sh| and \verb|test3.sh| (H.7, the depths with H.11), its closing
    paragraph no options without support, the messages~3 (H.7) and~8 (H.11) compared, and the counts
    of the last item (H.12); the files section names \verb|Ratpoints/Spec/Region.lean| (H.1),
    \verb|Ratpoints/Spec/Sturm.lean| (H.10) and \verb|Ratpoints/Model/CliNumbers.lean| (H.6).
  \item Chapter~\ref{chap:sieve}: the skeleton of a run refers to this chapter for the search region
    (H.12); the sentence that only \verb|findPoints| runs a component with \verb|Id.run| also names
    $\mathrm{findPointsWith}$, the program's run and its pure callback, \verb|Main.lean| and the
    \verb|#guard| lines (H.12); the search for the points (Definition~\ref{def:findPoints}) searches
    the search region of the options, normalized, with the search intervals and the refusal
    $\mathrm{badArgs}$ (H.5), and the Sturm step narrows the intervals (H.11); the run
    (Definition~\ref{def:sieveRun}) takes the bounds of the search intervals, gains
    \csource{find_points.c}{3913}{3913}, its contract quantifying the ranges of each denominator and
    its invariants taking the bounds (H.3) and the range (H.4), and uses
    Definition~\ref{def:exactBounds}; the reduction of the run to the reference
    (Theorem~\ref{thm:run_reference}) gives the points of the numerators in the ranges (H.3);
    Theorem~\ref{thm:findPoints} names
    \lname{Ratpoints.toFinset\_eq\_filter\_points\_of\_findPoints\_eq\_ok} for
    \verb|toFinset_eq_points_of_findPoints_eq_ok|, states membership and the set of elements in the
    search region, its uses gaining Definition~\ref{def:regionSpec} (H.5), under the hypothesis of
    Sturm's theorem, its uses gaining Definition~\ref{def:sturmTheorem} (H.11), and its proof uses
    Lemma~\ref{lem:regionFilter} (H.5) and Lemma~\ref{lem:elimSearchRegion} (H.11); the paragraph
    after it gives the twelve statements with their changes and the record (H.12).
  \item Chapter~\ref{chap:words}: one denominator (Definition~\ref{def:sift}) runs over the bounds
    of the search intervals with the lazy set-up, its links gaining
    \csource{find_points.c}{3129}{3130} (for \csource{find_points.c}{3129}{3129}),
    \csource{find_points.c}{3435}{3440} and \csource{rp-private.h}{574}{574}, its contract on the
    ranges, its \verb|\lean| list giving
    \lname{Ratpoints.Sieve.flatMap\_acceptedPoints\_eq\_flatMap\_checkNumWith}, which only the bits
    of one interval use now, to Lemma~\ref{lem:siftIntervals}, its uses gaining
    Definitions~\ref{def:intervalBits} and~\ref{def:numeratorRanges} and
    Lemma~\ref{lem:siftIntervals}, its uses of Lemmas~\ref{lem:accepted} and~\ref{lem:staged} going
    with that lemma to Lemma~\ref{lem:siftIntervals}, and its deviation ``every numerator from $-H$
    to $H$'' goes (H.3); Lemma~\ref{lem:sift_eq_checkDen}, whose equation names the numerators,
    takes those in the ranges (H.3); the summary of the C code of Section~\ref{sec:wordsC} refers to
    this chapter for the search intervals (H.12).
  \item Chapter~\ref{chap:setup}: the introduction and the paragraphs after the options and after
    Theorem~\ref{thm:refusals} say that the refusals add the search intervals, that membership and
    the order hold in the search region (H.5), and that the depth of the Sturm step changes only the
    speed by Theorem~\ref{thm:sturmSpeed} (H.11); the options (Definition~\ref{def:Options}) gain
    the search intervals and the range of the denominators with \csource{ratpoints.h}{359}{359},
    \csource{ratpoints.h}{363}{364} and \csource{ratpoints.h}{369}{369}, and use
    Definition~\ref{def:Double} (H.5); the remark after the remainder chain says that stage~H reads
    its signs (H.12); the reversal check (Definition~\ref{def:reversal}) takes the normalized
    switch, which folds in \verb|RATPOINTS_NO_REVERSE_AUTO| (H.5); the reversal exchanging $x$ and
    $z$ (Lemma~\ref{lem:reversal_swap}) takes
    \lname{Ratpoints.Setup.exists\_reversal\_eq\_error\_iff} from the refusal theorems, its
    statement gaining that the check refuses exactly when its switch is on, $d$ is even and
    $c_0 = c_1 = 0$ (H.5); the refusal theorems (Theorem~\ref{thm:refusals}) add the search
    intervals that are not well formed and restrict the refusal of $c_0 = c_1 = 0$ to the whole box,
    and use Lemmas~\ref{lem:searchIntervals} and~\ref{lem:reversalRegion} (H.5) and
    Lemma~\ref{lem:sturmRegion} (H.11); the permutation (Theorem~\ref{thm:findPoints_perm}) names
    \lname{Ratpoints.perm\_filter\_bruteForce\_of\_findPoints\_eq\_ok} for
    \verb|perm_bruteForce_of_findPoints_eq_ok|, a permutation of the reference filtered by
    $\mathrm{noY}$ and the search region, and
    \lname{Ratpoints.eq\_filter\_bruteForce\_of\_findPoints\_eq\_ok\_of\_reverse\_eq\_false} for
    \verb|eq_bruteForce_of_findPoints_eq_ok_of_reverse_eq_false|, which moves to it from the output
    of the search, its statement gaining that $L$ is the reference filtered by the search region
    when also the options forbid the reversal and the search of $c$ does not take the divisors times
    squares, and its uses gaining Lemmas~\ref{lem:regionFilter}, \ref{lem:searchDensRange}
    and~\ref{lem:normalize} (H.5), and both theorems of the node hold under Sturm's theorem, its
    uses gaining Lemma~\ref{lem:elimSearchRegion} (H.11); the search
    (Definition~\ref{def:searchPoints}) takes the range (H.4), the search intervals (H.5) and the
    flag \verb|sturm_empty| (H.11); the paragraph on the search with a callback describes its flow
    with the normalization and the search intervals (H.5) and the Sturm step (H.11), the Sturm
    step's range going to Definition~\ref{def:sturmRegion} (H.11); Section~\ref{sec:setupCli} lists
    no option as not supported (H.7); the deviations lose the reversal check without
    \verb|RATPOINTS_NO_REVERSE_AUTO| and ``no error~9'' (H.5) and the Sturm step with only its test
    for a repeated factor (H.11); the sentence on the macros of \verb|ratpoints.h| records
    \verb|RATPOINTS_NO_REVERSE_AUTO| (H.5); the files section names the Sturm step in
    \verb|Ratpoints/Model/Sturm.lean| (H.8).
  \item Chapter~\ref{chap:denominators}: the summary of the C code says that $b_{\mathrm{high}}$ is
    the normalized largest denominator (H.4); the shapes (Definition~\ref{def:Shape}) are chosen by
    \verb|run_setup|, which returns the bounds of the search intervals with the shape and gains
    \csource{find_points.c}{3869}{3870}, the node naming \lname{Ratpoints.Sieve.fst\_run\_runSetup}
    and \lname{Ratpoints.Sieve.snd\_run\_runSetup}, by which the shape that \verb|run_setup| chooses
    and the bounds that it makes are those of the flags and of the search intervals, and using
    Definition~\ref{def:exactBounds} (H.3), and reads the largest denominator for the Jacobi tables
    (H.4); the denominators of the search (Definition~\ref{def:searchDens}) take the range, in
    \verb|Ratpoints/Proof/SearchDens.lean| (H.4); the output of the search
    (Theorem~\ref{thm:findPoints_order}) is stated over the region that the search uses (H.5 over
    the search intervals, H.11 over the region of Definition~\ref{def:searchRegion}, which it then
    uses), names \lname{Ratpoints.perm\_filter\_bruteForce\_append\_flatMap\_searchDens} for
    \verb|perm_bruteForce_append_flatMap_searchDens| and uses Lemma~\ref{lem:searchDensRange} (H.4),
    and gives \verb|eq_bruteForce_of_findPoints_eq_ok_of_reverse_eq_false| to the permutation
    (Theorem~\ref{thm:findPoints_perm}), its statement losing the clause that $L$ is
    $\mathrm{bruteForce}(c, H)$ when also the options forbid the reversal and the search of $c$ does
    not take the divisors times squares (H.5); the deviations lose ``no search intervals and no
    range of the denominators other than $1$ to $H$'' (H.5) and the exclusions leave the search
    intervals to no later stage (H.3); the files section names
    \verb|Ratpoints/Proof/SearchDens.lean| (H.4).
  \item Chapter~\ref{chap:tables}: what \verb|sieving_info| computes is sound for every height and
    takes the largest denominator (Lemma~\ref{lem:sievingInfo}, H.4).
  \item Chapter~\ref{chap:stages}: the deviations lose the search intervals of the sieving thread,
    and the exclusions leave the search intervals (H.3) and the Sturm step (H.11) to no later stage.
  \item Chapter~\ref{chap:packing}: the summary of the C code and the deviation of the bits of an
    interval say that the form of the conversion to bits is chosen by the bounds
    (Definition~\ref{def:intervalBits}), and the deviation of the bits of one denominator
    (Definition~\ref{def:siftBits}) that the sieve specs are made at the first interval with the
    code $0$, as the C code makes them (H.3); the node of Definition~\ref{def:siftBits} says that
    one denominator turns each interval into bits (H.3); the exclusions leave nothing to stage~H
    (H.3).
  \item Chapter~\ref{chap:rows}: the exclusions leave the lazy set-up to no later stage (H.3).
  \item Chapter~\ref{chap:output}: the introduction says that the report prints the messages on the
    search region (message~3, H.7) and on the search intervals (message~8, H.11) and leaves out only
    that on the moduli used (message~4, stage~I); what a search gives back
    (Definition~\ref{def:Outcome}) writes back the search intervals and uses
    Definition~\ref{def:Double} (H.11); the search with a callback
    (Definition~\ref{def:findPointsWith}) normalizes the options, validates the search intervals
    (H.5) and runs the Sturm step of Definition~\ref{def:sturmRegion}, which takes
    \csource{find_points.c}{4857}{4892} from it (H.11), its uses gaining
    Definitions~\ref{def:normalize}, \ref{def:searchIntervals} and~\ref{def:sturmRegion}; the
    sieving (Definition~\ref{def:sieving}) takes the range (H.4) and the search intervals (H.5); in
    Section~\ref{sec:searchDelivery} the return when the Sturm step leaves nothing to search is met
    (H.11); the sentence of the contracts of the sieve that attaches
    \verb|Ratpoints.Sieve.flatMap_acceptedPoints_eq_flatMap_checkNumWith| to the contract of one
    denominator says that it serves the bits of one search interval (H.3); the points that the
    search delivers (Definition~\ref{def:searchList}) take the intervals, their uses gaining
    Definition~\ref{def:inIntervals} (H.3), and the range (H.4); so does Lemma~\ref{lem:searchList},
    which uses Lemmas~\ref{lem:sift_eq_checkDen}, \ref{lem:exactBounds}
    and~\ref{lem:numeratorRanges} (H.3) and Lemma~\ref{lem:searchDensRange} (H.4), gains
    \lname{Ratpoints.exists\_isSoundTest\_run\_searchPoints\_appendPoint}, with the sentence that
    for a sound test of the sieve the search with the appending callback from the empty array gives
    the list of Definition~\ref{def:searchList}, for every choice of the flags (H.5), that lemma
    taking the argument \verb|sturm_empty| and the lemma using Lemma~\ref{lem:searchListRegion},
    through \lname{Ratpoints.searchList\_nil}, at H.11; Lemma~\ref{lem:findPoints_searchList} is
    stated over the search intervals, its proof using Lemmas~\ref{lem:normalize}
    and~\ref{lem:searchIntervals} (H.5), and over the region of Definition~\ref{def:searchRegion},
    which its statement uses, its proof using Lemmas~\ref{lem:searchRegion}
    and~\ref{lem:sturmRegion} (H.11); the survivors under \verb|RATPOINTS_NO_CHECK|
    (Lemma~\ref{lem:noCheckList}) hold for any intervals of the values, their uses gaining
    Definition~\ref{def:inIntervals} (H.3), and for the range (H.4): the equation and the sublist
    only gain the parameters, membership needs $1 \le b_{\mathrm{low}}$ and
    $b_{\mathrm{high}} \le H$, no duplicates also $1 \le b_{\mathrm{high}}$, and the denominators of
    the search have no duplicates for $1 \le b_{\mathrm{high}}$, with \verb|nodup_searchDens| moving
    to \verb|Ratpoints/Proof/SearchDens.lean|, the proof using Lemma~\ref{lem:searchDensRange} for
    the range of the denominators in place of $1, \dots, H$ (H.4); the protocol
    (Theorem~\ref{thm:protocol}) keeps its statement, its proof using Lemmas~\ref{lem:normalize}
    and~\ref{lem:searchIntervals} (H.5) and Lemma~\ref{lem:sturmRegion},
    \verb|map_findPointsWith_flags| covering the search intervals written back (H.11); the
    statements under \verb|RATPOINTS_NO_CHECK| (Theorem~\ref{thm:noCheck}) are those of
    Section~\ref{sec:certifiedH} (H.5, H.11); the paragraph of Section~\ref{sec:outputOptions} on
    the main theorems says that membership and the set of elements hold in the search region, names
    \lname{Ratpoints.toFinset\_eq\_filter\_points\_of\_findPoints\_eq\_ok} and
    \lname{Ratpoints.perm\_filter\_bruteForce\_of\_findPoints\_eq\_ok}, states the theorem of the
    output of the search over the search intervals and the permutation for either setting of
    $\mathrm{noY}$, filtered by $\mathrm{noY}$ and the search region, only the reference itself
    taking $\mathrm{noY}$ off, and the refusal theorems with the search intervals and the whole box
    (H.5), and adds the hypothesis of Sturm's theorem and the region of
    Definition~\ref{def:searchRegion} (H.11); the points in the order of ratpoints
    (Definition~\ref{def:orderedPoints}) take the largest denominator (H.5); the denominators and
    the points of the search (Lemma~\ref{lem:searchDenominators}) name
    \lname{Ratpoints.searchDens\_eq\_filter\_searchDenominators} for
    \verb|searchDens_eq_searchDenominators|, by which the denominators of the search from
    $b_{\mathrm{low}}$ on are those of the specification up to $b_{\mathrm{high}}$ that are at least
    $b_{\mathrm{low}}$, and gain \lname{Ratpoints.Sieve.squares1Dens\_eq\_filter} and
    \lname{List.range'\_eq\_filter\_range'}, by which the divisors times squares and the numbers
    from $b_{\mathrm{low}} \ge 1$ on are those from $1$ without the denominators below
    $b_{\mathrm{low}}$ (H.4); the order theorem (Theorem~\ref{thm:order}) takes the search region
    (H.5) and Sturm's theorem (H.11) and uses Lemmas~\ref{lem:regionFilter}
    and~\ref{lem:elimSearchRegion}; Section~\ref{sec:certifiedG} refers to
    Section~\ref{sec:certifiedH} (H.12); the program: the paragraph after
    Definition~\ref{def:cliOutput} loses the clause that \verb|-l -u -dl -du| keep the exit
    status~12, and the settings of the output (Definition~\ref{def:cliOutput}) gain the exact ends,
    their uses gaining Definition~\ref{def:cliEnds} (H.7); the callback of the program
    (Definition~\ref{def:cliProcess}) tests the search region, with \csource{main.c}{150}{156} in
    place of \csource{main.c}{153}{156}, its data gain \csource{main.c}{246}{247}, and it uses
    Definition~\ref{def:cliRegionTest} (H.7); the report (Definition~\ref{def:cliReport}) gains
    message~3, printed only when the command line gives intervals, after the height bound, or after
    the curve equation under \verb|-v|, with \csource{main.c}{224}{224} and
    \lname{Ratpoints.Cli.showIntervals} in its \verb|\lean| list, its uses gaining
    Definitions~\ref{def:Double} and~\ref{def:cliRegionTest} (the doubles of messages~3 and~8, and
    \verb|printF|) (H.7), and message~8 (H.11), with the texts \csource{main.c}{979}{991} and
    \csource{main.c}{999}{1033}; the theorem on the program (Theorem~\ref{thm:cliRun}) counts the
    points that the callback keeps, its \verb|\lean| list gaining \lname{Ratpoints.Cli.Data.keeps}
    and \lname{Ratpoints.Cli.Data.keeps\_printPoint} and its statement saying which points the
    callback keeps (H.7); the section on the differential tests refers to
    Section~\ref{sec:regionTests}, its item on the report drops only messages~4 and~8 and compares
    message~3 (H.7) and drops only message~4 (H.11), its item on \verb|-l -u -dl -du| and the exit
    status~12 goes (H.7), and its two runs of \verb|test3.sh| that need messages~4 and~8 still need
    message~4, of stage~I (H.11); the deviations lose those that stage~H retires
    (Section~\ref{sec:regionDeviations}), each with its item, and the lines left free for stage~H
    are recorded (H.12).
  \item The stage list of the introduction refers to this chapter (Chapter~\ref{chap:region}),
    which is input after Chapter~\ref{chap:output} (H.12). The generated table follows each item.
\end{itemize}

\section{Files, names and order of work}\label{sec:regionFiles}

The model, in the folder \verb|Ratpoints/Model|: the doubles in a new file without imports
\verb|Double.lean| (about $350$ lines with its guards); the ends, the scaled ends, the forms and the
bounds in a new file \verb|Bounds.lean| (about $420$ lines); the bits of an interval and the loop of
one denominator over the intervals in \verb|Sieve.lean| (about $1330$ lines); \verb|run_setup| and
the run in \verb|Denominators.lean|; the Sturm step in \verb|Sturm.lean| (about $490$ lines); the
options, their normalization, the search intervals, the Sturm step of the search, what a search
gives back, the sieving, the search and the search with a callback in \verb|FindPoints.lean| (about
$1170$ lines); the refusal $\mathrm{badArgs}$ in \verb|Setup.lean|; the decided facts
$\mathrm{maxDegree} < 2^{57}$ and $0 < \mathrm{maxDegree}$ in \verb|Constants.lean|, which the build
with other values of the constants checks at $\mathrm{maxDegree} = 30$ and $200$; the readers of the
numbers, \verb|strtod|, the ends of the command line, the test of the region, the join and \verb|%f|
in a new file \verb|CliNumbers.lean| (about $530$ lines); and the options of the region, the reading
of the intervals, the callback with the test of the region and messages~3 and~8 in \verb|Cli.lean|
(about $1060$ lines). The specification: \verb|Ratpoints/Spec/Region.lean| (about $140$ lines) and
\verb|Ratpoints/Spec/Sturm.lean| (about $50$ lines), new, and the largest denominator in
\verb|Ratpoints/Spec/Order.lean|. The proofs, in the folder \verb|Ratpoints/Proof|, the contracts of
the sieve and of the loops of the Sturm step still in files without Mathlib: the doubles in a new
file \verb|Double.lean| (about $1140$ lines) and the faithful roundings in a new file
\verb|Rounding.lean| (about $770$ lines); the values and the bits of the intervals, the ranges of
the numerators and the normalized options in a new file \verb|Region.lean| (about $280$ lines); the
search intervals and the bounds in a new file \verb|Bounds.lean| (about $1380$ lines); the
representation of a range by the bounds and the ranges of a run in the core file
\verb|BitRange.lean|; the loop over the intervals and the contract of one denominator in the core
file \verb|Sift.lean| (about $1270$ lines); the contract of the run and what \verb|run_setup|
returns in \verb|Denominators.lean|; the denominators of the search in a new file
\verb|SearchDens.lean|, split off \verb|FindPoints.lean| (about $240$ and $1310$ lines); the list of
the brute-force reference split at the candidate $(1, 0)$ and filtered to a range of denominators in
\verb|BruteForce.lean| (Lemma~\ref{lem:searchDensRange}), the divisors times squares in a range as a
filter in \verb|Shapes.lean| (Lemma~\ref{lem:searchDenominators}), the points at infinity and the
points of a numerator in the search intervals in \verb|Stages.lean|
(Lemma~\ref{lem:searchListRegion}), and the degree after the reversal check in \verb|Setup.lean|
(Lemma~\ref{lem:sturmRegion}); the contracts of the signs, the counts and the intersection, and the
failure of the step, in the core file \verb|Sturm.lean| (about $920$ lines); the bisection in a new
file \verb|Bisection.lean| (about $1060$ lines); the chain as a Sturm sequence in a new file
\verb|SturmSigns.lean| (about $400$ lines); the counts at the ends of an interval in a new file
\verb|SturmCounts.lean| (about $270$ lines); the coverage and the bound of the bisection in a new
file \verb|SturmCoverage.lean| (about $670$ lines); the Sturm step of the search and the region that
the search uses in a new file \verb|SturmRegion.lean| (about $300$ lines); the theorems with the
search region in \verb|Output.lean| (about $1090$ lines); the numbers of the command line in a new
file \verb|CliNumbers.lean| (about $1120$ lines) and the theorem on the region of the program in
\verb|Cli.lean| (about $950$ lines); and the general facts on core types in \verb|Containers.lean|
and \verb|Bits.lean|, under their names in the namespaces of core. The candidates for Mathlib in the
mirrors \verb|Ratpoints/Mathlib/Algebra/Polynomial/FieldDivision.lean|,
\verb|Ratpoints/Mathlib/Algebra/Polynomial/Sturm/Sequence.lean|,
\verb|Ratpoints/Mathlib/Analysis/Polynomial/Basic.lean|,
\verb|Ratpoints/Mathlib/FieldTheory/Perfect.lean| and
\verb|Ratpoints/Mathlib/Order/WithBotTop.lean|, new, and
\verb|Ratpoints/Mathlib/Data/List/Pairwise.lean|; Mathlib after the version that the project pins
has \verb|Polynomial.signRight x p| (pull request \#44248,
\verb|Mathlib/Analysis/Polynomial/SignRight.lean|), the sign just right of $x$ defined as a limit
for every $p$, which the project can take over at its next update of Mathlib, so that the mirror
\verb|Ratpoints/Mathlib/Analysis/Polynomial/Basic.lean| defines no sign of its own just right or
just left of a point. Every Lean file stays below $1500$ lines (the largest,
\verb|Ratpoints/Proof/Bounds.lean|, about $1380$). The new modules join their part roots and
\verb|Ratpoints.lean|. One commit per item, each passing the project's checks, the build with other
values of the constants and the differential tests, with the nodes and the passages of the earlier
chapters that it changes, the C paragraphs of the model and the \verb|\csource| links of its nodes,
locked, the C links moved from their old owners and the exclusions in the same commit, the
\verb|\lname| of its declarations turned into \verb|\lean|, and the dependency audit of the
blueprint after each item that changes nodes; a lemma of a node comes with the item that first uses
it. The output of \verb|ratpoints-lean| does not change up to H.6, changes at H.7, where the options
of the region take effect and message~3 enters the reports, and at H.11, where the Sturm step
narrows the search intervals, which changes the survivors under \verb|-x|, and message~8 enters the
reports. The certified statements change twice: at H.5, with the search region and without the
hypothesis of Sturm's theorem, since the model does not narrow the intervals by the Sturm step
before H.11, and at H.11, where statements 3, 8 and~10 take that hypothesis.
\begin{enumerate}
  \item[H.1] The doubles (Definition~\ref{def:Double}) in \verb|Ratpoints/Model/Double.lean|, with
    \verb|#guard| examples against gcc and glibc, but for $\mathrm{ofRatio}$, which comes with H.6;
    the beginning of \verb|Ratpoints/Spec/Region.lean|, $\mathrm{IsNaN}$ and $\mathrm{doubleValue}$
    with their examples, not yet restated by the challenge; their lemmas (Lemma~\ref{lem:Double}) in
    \verb|Ratpoints/Proof/Double.lean|; the includes \csource{rp-private.h}{30}{35},
    \csource{rp-private.h}{60}{60} and \csource{ratpoints.h}{31}{31} excluded; the files section of
    Chapter~\ref{chap:spec}; and the nodes and passages of the earlier chapters that
    Section~\ref{sec:regionEarlier} assigns to it. Nothing uses them yet.
  \item[H.2] The exact bounds: the ends, the scaled ends, the forms and the bounds
    (Definitions~\ref{def:exactEnd}, \ref{def:exactScale} and~\ref{def:exactBounds}) in
    \verb|Ratpoints/Model/Bounds.lean|, the bits of an interval (Definition~\ref{def:intervalBits})
    in \verb|Ratpoints/Model/Sieve.lean|, with their guards against the C code run verbatim and the
    build under \verb|DEBUG|; $\mathrm{AreSearchIntervals}$ and $\mathrm{IntervalsWF}$ of the
    specification; the values and the bits of the intervals and the ranges
    (Definitions~\ref{def:inIntervals} and~\ref{def:numeratorRanges}) in
    \verb|Ratpoints/Proof/Region.lean| and \verb|Ratpoints/Proof/BitRange.lean|; the lemmas of the
    ranges (Lemma~\ref{lem:numeratorRanges}) in \verb|Ratpoints/Proof/Region.lean|, with the mirror
    lemma on chains of pairs in \verb|Ratpoints/Mathlib/Data/List/Pairwise.lean|; the lemmas of the
    ends, the bits and the bounds (Lemmas~\ref{lem:exactEnd}, \ref{lem:intervalBits}
    and~\ref{lem:exactBounds}) in \verb|Ratpoints/Proof/Bounds.lean|, without its part on the search
    intervals; the links \csource{find_points.c}{3217}{3221}, \csource{find_points.c}{3231}{3238},
    \csource{find_points.c}{3257}{3257}, \csource{find_points.c}{3264}{3340},
    \csource{find_points.c}{3381}{3406} and \csource{find_points.c}{3410}{3434}, as the C paragraphs
    of the components cut them, and the exclusion of \csource{find_points.c}{3239}{3256}; and the
    nodes and passages of the earlier chapters that Section~\ref{sec:regionEarlier} assigns to it.
    Nothing calls them yet.
  \item[H.3] One denominator over the intervals with the lazy set-up, for the default interval only,
    so that the output does not change: \verb|sift| over the bounds, \verb|run_setup| returning
    them, with \lname{Ratpoints.Sieve.fst\_run\_runSetup} and
    \lname{Ratpoints.Sieve.snd\_run\_runSetup}, the run taking them, with their links
    (Section~\ref{sec:regionSieve}); $\mathrm{searchIntervals}$ with the default interval only
    (\csource{find_points.c}{4829}{4829} and lines 4831-4841, the default interval, as the link of
    Definition~\ref{def:searchIntervals} until H.5 extends it to 4831-4855), which the sieving calls
    with no intervals, and its lemmas on the default interval (Lemma~\ref{lem:searchIntervals}, but
    for those of the validation); the loop over the intervals (Lemma~\ref{lem:siftIntervals}, with
    \lname{Ratpoints.Sieve.flatMap\_acceptedPoints\_eq\_flatMap\_checkNumWith} from the node of one
    denominator) in \verb|Ratpoints/Proof/Sift.lean|, the contracts of one denominator and of the
    run on the ranges, the lemmas on lists in \verb|Ratpoints/Proof/Containers.lean| and
    \verb|Ratpoints/Proof/BitRange.lean|, and the points that the search delivers over the intervals
    of its values, their statements kept by the default interval
    (\lname{Ratpoints.exists\_run\_searchIntervals\_nil},
    \lname{Ratpoints.intervalsWF\_and\_inIntervals\_neg}); comparator, since the proofs of the
    certified statements change; the nodes of Definitions~\ref{def:sift}, \ref{def:sieveRun},
    \ref{def:Shape}, \ref{def:siftBits} and~\ref{def:searchList}, Lemmas~\ref{lem:sift_eq_checkDen},
    \ref{lem:searchList} and~\ref{lem:noCheckList} and Theorem~\ref{thm:run_reference}, the sentence
    of Chapter~\ref{chap:output} on the contracts of the sieve, the summary of the C code of
    Chapter~\ref{chap:packing}, and the deviations and exclusions of Chapters~\ref{chap:words},
    \ref{chap:denominators}, \ref{chap:stages}, \ref{chap:packing} and~\ref{chap:rows} that leave
    the search intervals or the lazy set-up to stage~H; and the nodes and passages of the earlier
    chapters that Section~\ref{sec:regionEarlier} assigns to it.
  \item[H.4] The range of the denominators through the search: the sieving and the search take
    $b_{\mathrm{low}}$ and $b_{\mathrm{high}}$, \verb|run_setup| the largest denominator, and
    $\mathrm{findPoints}$ passes $1$ and $H$; the denominators of the search with the range in a new
    file \verb|Ratpoints/Proof/SearchDens.lean|, split off \verb|Proof/FindPoints.lean| with the
    lemmas on them (Lemma~\ref{lem:searchDensRange}), the renames
    \lname{Ratpoints.searchDens\_eq\_filter\_searchDenominators} and
    \lname{Ratpoints.perm\_filter\_bruteForce\_append\_flatMap\_searchDens},
    \lname{Ratpoints.Sieve.squares1Dens\_eq\_filter} and \lname{List.range'\_eq\_filter\_range'}
    (Lemma~\ref{lem:searchDenominators}), and what \verb|sieving_info| computes sound for every
    height (Lemma~\ref{lem:sievingInfo}); comparator, as in H.3; the nodes of
    Definitions~\ref{def:searchDens}, \ref{def:sieving}, \ref{def:searchPoints}, \ref{def:sieveRun},
    \ref{def:Shape} and~\ref{def:searchList}, Lemmas~\ref{lem:sievingInfo}, \ref{lem:searchList},
    \ref{lem:searchDenominators} and~\ref{lem:noCheckList}, Theorem~\ref{thm:findPoints_order}, the
    summary of the C code of Chapter~\ref{chap:denominators}, and the files sections of
    Chapters~\ref{chap:denominators} and~\ref{chap:output}; and the nodes and passages of the
    earlier chapters that Section~\ref{sec:regionEarlier} assigns to it.
  \item[H.5] The search region of the options, without the Sturm step yet: the fields of the options
    with \csource{ratpoints.h}{359}{359}, \csource{ratpoints.h}{363}{364} and
    \csource{ratpoints.h}{369}{369}; their normalization (Definition~\ref{def:normalize}) and its
    lemmas (Lemma~\ref{lem:normalize}); the validation of the search intervals with the refusal
    $\mathrm{badArgs}$ (Definition~\ref{def:searchIntervals}, whose link becomes
    \csource{find_points.c}{4831}{4855}), the rest of Lemma~\ref{lem:searchIntervals}, and error~9
    of the program, with \csource{main.c}{1048}{1048} and its exclusion gone; the search with a
    callback in the order of the C code, the reversal check with the normalized switch, the search
    over the search intervals of the options; the rest of \verb|Ratpoints/Spec/Region.lean|
    (Definition~\ref{def:regionSpec}) and the largest denominator of
    \verb|Ratpoints/Spec/Order.lean|, with its examples, the example of the interval of
    \verb|-l -0.2 -u 0.6| changed to the doubles that the program passes the library, the upper end
    \verb|0x3fe3333333333334| in place of the double \verb|0x3fe3333333333333| nearest to $0.6$, in
    the list of $\mathrm{IsFullBox}$ and in the filter by $\mathrm{InRegion}$, checked by
    \verb|decide +kernel| as before; Lemmas~\ref{lem:reversalRegion} and~\ref{lem:regionFilter},
    \lname{Ratpoints.Setup.exists\_reversal\_eq\_error\_iff} moving to
    Lemma~\ref{lem:reversal_swap}, Lemma~\ref{lem:searchListRegion} but for
    \lname{Ratpoints.searchList\_nil}, and
    \lname{Ratpoints.exists\_isSoundTest\_run\_searchPoints\_appendPoint}
    (Lemma~\ref{lem:searchList}), lifted out of the list of a successful search when that lemma is
    restated with the region, without the argument \verb|sturm_empty| yet; statements 3, 5, 6, 8, 9
    and~10 with the search region, without the hypothesis of Sturm's theorem, proved from the search
    intervals of the options; the corollaries with the search region and their renames, the theorem
    on the reference itself moving to the node of the permutation; the challenge restates the file
    of the region and imports its Mathlib modules, takes $\mathrm{Double}$ from the model and states
    the six statements as changed, \verb|Solution| proves them, and comparator; the nodes of
    Definitions~\ref{def:Options}, \ref{def:reversal}, \ref{def:findPointsWith},
    \ref{def:findPoints}, \ref{def:orderedPoints}, \ref{def:cliFailure}, \ref{def:sieving}
    and~\ref{def:searchPoints}, Lemmas~\ref{lem:reversal_swap}, \ref{lem:searchList}
    and~\ref{lem:findPoints_searchList}, Theorems~\ref{thm:findPoints}, \ref{thm:findPoints_perm},
    \ref{thm:refusals}, \ref{thm:findPoints_order}, \ref{thm:order}, \ref{thm:noCheck}
    and~\ref{thm:protocol}, Section~\ref{sec:comparator-ab}, the paragraph after
    Definition~\ref{def:bruteForce}, the introduction of Chapter~\ref{chap:setup} with its
    paragraphs after the options and after Theorem~\ref{thm:refusals}, its paragraph on the search
    with a callback and its sentence on the macros of \verb|ratpoints.h|, the paragraph of
    Section~\ref{sec:outputOptions} on the main theorems, and the deviations of
    Chapters~\ref{chap:setup} and~\ref{chap:denominators} that H.5 retires; and the nodes and
    passages of the earlier chapters that Section~\ref{sec:regionEarlier} assigns to it.
  \item[H.6] The numbers of the program: $\mathrm{ofRatio}$ of the doubles, and
    \verb|Ratpoints/Model/CliNumbers.lean| with the readers moved from \verb|Cli.lean|,
    \verb|strtod|, the ends, the test of the region, the join and \verb|%f|
    (Definitions~\ref{def:cliEnds} and~\ref{def:cliRegionTest}), with their links,
    \csource{main.c}{746}{752} moving from \verb|isDouble| to $\mathrm{scanDouble?}$, and the
    exclusions of the memory of \verb|outward| and \verb|scan_end|, and guards against glibc; the
    faithful roundings (Lemma~\ref{lem:rounding}) in \verb|Ratpoints/Proof/Rounding.lean|, the
    lemmas of the numbers (Lemma~\ref{lem:cliNumbers}) in \verb|Ratpoints/Proof/CliNumbers.lean|,
    with \lname{Char.toLower\_eq\_x\_iff} in \verb|Ratpoints/Proof/Containers.lean| and the mirror
    \verb|Order/WithBotTop|; the node of Definition~\ref{def:cliNumbers} and the files section of
    Chapter~\ref{chap:spec}; and the nodes and passages of the earlier chapters that
    Section~\ref{sec:regionEarlier} assigns to it. Nothing of the program calls them yet.
  \item[H.7] The region of the program: \verb|-l -u -dl -du| read as \verb|main.c| reads them, the
    intervals with the rules of error~7 (Definition~\ref{def:cliIntervals}), the callback that tests
    the exact region, message~3 when the command line gives intervals, the search on the joined
    intervals, and the exit status~12 and \verb|Ratpoints.Cli.Failure.unsupportedOption| deleted,
    with their links (Section~\ref{sec:regionCli}); $0 < \mathrm{maxDegree}$; the theorem on the
    region of the program (Theorem~\ref{thm:cliRegion}) in \verb|Ratpoints/Proof/Cli.lean|, its last
    statement without the hypothesis of Sturm's theorem, and the theorem on the program with the
    test of the callback (Theorem~\ref{thm:cliRun}, with \lname{Ratpoints.Cli.Data.keeps} and
    \lname{Ratpoints.Cli.Data.keeps\_printPoint}); the harness compares the region runs of
    \verb|test4.sh| and \verb|test3.sh|, admits large heights with a narrow region, has the suite
    \verb|regions| and the cases of the region in the errors suite and the fuzz, compares message~3
    and drops message~8 from the report of the C program, its option status~12 gone, and under
    \verb|-x| the weaker check also tests $z$ against the range of the denominators and $x/z$
    against the exact region of the command line, as \verb|process| of \verb|main.c| does, and the
    exact comparison (still with \verb|-p 30 -N 30 -s| appended) runs only when the C program on
    that same command line sieves with the first $30$ odd primes in its first two stages and no
    third stage (or does not sieve), as its report under \verb|-v -z| says, otherwise the weaker
    check alone running and its reason (stage~I) counted; the output changes; the nodes of
    Definitions~\ref{def:cliInput}, \ref{def:cliFailure}, \ref{def:cliRun}, \ref{def:cliOutput},
    \ref{def:cliProcess} and~\ref{def:cliReport}, Theorem~\ref{thm:cliRun}, Sections~\ref{sec:cli}
    and~\ref{sec:setupCli}, the introduction of Chapter~\ref{chap:output} and its paragraph after
    Definition~\ref{def:cliOutput}, the sections on the differential tests of
    Chapters~\ref{chap:spec} and~\ref{chap:output}, and the deviations of Chapter~\ref{chap:output}
    on the options of the region, the status~12, message~3 and the test of the callback; and the
    nodes and passages of the earlier chapters that Section~\ref{sec:regionEarlier} assigns to it.
  \item[H.8] The Sturm step, first part: the signs and the counts and the intersection
    (Definitions~\ref{def:evalSign} and~\ref{def:intersect}) with their links, and the exclusions of
    the includes of \verb|sturm.c| and of the copy on the heap; their contracts
    (Lemma~\ref{lem:evalSign}, and of Lemma~\ref{lem:intersect} the recursion of the intersection,
    its invariant and its contract) in \verb|Ratpoints/Proof/Sturm.lean|; the half of
    \verb|Ratpoints/Proof/SturmSigns.lean| that relates the signs and the counts to Mathlib's sign
    changes (Lemma~\ref{lem:chainSturm} but for the chain); the files section of
    Chapter~\ref{chap:setup}; and the nodes and passages of the earlier chapters that
    Section~\ref{sec:regionEarlier} assigns to it. Nothing calls them yet.
  \item[H.9] The Sturm step, second part: $\mathrm{LONG\_MAX}$, the midpoint, the kept interval, the
    bisection and the step (Definitions~\ref{def:iterate} and~\ref{def:computeSturm}) with their
    links, the step not yet called, and the exclusion of its prototype
    (\csource{rp-private.h}{617}{617}); the bisection and the rest of the intersection and the step
    (Lemmas~\ref{lem:bisection} and~\ref{lem:intersect}) in \verb|Ratpoints/Proof/Bisection.lean|,
    with the lemmas on \verb|Int64| of Lemma~\ref{lem:bisection} in \verb|Ratpoints/Proof/Bits.lean|
    and \lname{Ratpoints.Sturm.run\_computeSturm\_eq\_none\_iff} in
    \verb|Ratpoints/Proof/Sturm.lean|; and the nodes and passages of the earlier chapters that
    Section~\ref{sec:regionEarlier} assigns to it.
  \item[H.10] Sturm's theorem and the coverage: \verb|Ratpoints/Spec/Sturm.lean|
    (Definition~\ref{def:sturmTheorem}), not yet restated by the challenge; the mirrors of
    Lemmas~\ref{lem:posAssoc} and~\ref{lem:polySigns}, the latter, a change against the prototype,
    without its definitions of the signs just right and just left of a point, its two lemmas on
    these signs stating the sign by the case distinction on $p(x) = 0$ in their conclusions, and
    $\mathrm{sgnLeftEnd}$ and $\mathrm{sgnRightEnd}$ of \verb|Ratpoints/Proof/SturmCounts.lean|
    reading the same case distinction; the counts at the ends of an interval
    (Lemma~\ref{lem:sturmCounts}) in \verb|Ratpoints/Proof/SturmCounts.lean|; the chain as a Sturm
    sequence, the rest of Lemma~\ref{lem:chainSturm}; the coverage and the bound of the bisection
    (Lemma~\ref{lem:coverage}, and of Lemma~\ref{lem:sturmBuffers} all but the two bounds of the
    search) in \verb|Ratpoints/Proof/SturmCoverage.lean|; the files section of
    Chapter~\ref{chap:spec}; and the nodes and passages of the earlier chapters that
    Section~\ref{sec:regionEarlier} assigns to it.
  \item[H.11] The switch: the Sturm step of the search (Definition~\ref{def:sturmRegion}) called by
    the search with a callback, which gives it \csource{find_points.c}{4857}{4892}, with the return
    of line 4889 and the flag \verb|sturm_empty| of the search; what a search gives back with the
    search intervals; message~8 in the report; the lemmas of the step of the search, the region that
    the search uses and the bounds of the arrays (Lemmas~\ref{lem:sturmRegion}
    and~\ref{lem:searchRegion}, Definition~\ref{def:searchRegion}, the two bounds of
    Lemma~\ref{lem:sturmBuffers}) in \verb|Ratpoints/Proof/SturmRegion.lean|, with the decided fact
    $\mathrm{maxDegree} < 2^{57}$ in \verb|Ratpoints/Model/Constants.lean| and
    \lname{Ratpoints.Setup.degree\_le\_of\_reversal\_eq\_ok} in \verb|Ratpoints/Proof/Setup.lean|;
    the list of the search and its output over the region that the search uses, with the rest of
    Lemma~\ref{lem:searchListRegion}, \lname{Ratpoints.searchList\_nil}, and the argument
    \verb|sturm_empty| of \lname{Ratpoints.exists\_isSoundTest\_run\_searchPoints\_appendPoint};
    Lemma~\ref{lem:elimSearchRegion}; statements 3, 8 and~10 with the hypothesis of Sturm's theorem,
    and the corollaries with it; Theorem~\ref{thm:sturmSpeed}; the theorem on the region of the
    program with the hypothesis; the challenge restates the file of Sturm's theorem and imports its
    Mathlib modules, states the three statements with the hypothesis, \verb|Solution| proves them,
    and comparator with \verb|--paranoid|; the harness compares \verb|-x| exactly without \verb|-s|
    (the check of the moduli now on the command line without \verb|-s|), has the suite \verb|depths|
    and compares message~8; the output changes; the nodes of Definitions~\ref{def:findPointsWith},
    \ref{def:Outcome}, \ref{def:searchPoints}, \ref{def:findPoints} and~\ref{def:cliReport},
    Lemmas~\ref{lem:searchList} and~\ref{lem:findPoints_searchList}, Theorems~\ref{thm:findPoints},
    \ref{thm:findPoints_perm}, \ref{thm:order}, \ref{thm:noCheck}, \ref{thm:refusals},
    \ref{thm:findPoints_order}, \ref{thm:protocol} and~\ref{thm:cliRegion},
    Sections~\ref{sec:comparator-ab}, \ref{sec:cli} and~\ref{sec:searchDelivery}, the sections on
    the differential tests of Chapters~\ref{chap:spec} and~\ref{chap:output}, the introduction of
    Chapter~\ref{chap:output}, the paragraph of Section~\ref{sec:outputOptions} on the main
    theorems, the introduction of Chapter~\ref{chap:setup} with its paragraphs and its paragraph on
    the search with a callback, the exclusions of Chapter~\ref{chap:stages}, and the deviations of
    Chapters~\ref{chap:setup} and~\ref{chap:output} on the Sturm step, the search intervals written
    back and message~8; and the nodes and passages of the earlier chapters that
    Section~\ref{sec:regionEarlier} assigns to it.
  \item[H.12] The records and the texts: the description of the repository,
    \verb|formalization.yaml| and \verb|comparator/README.md| with the twelve statements and the
    record of comparator; the passages of the earlier chapters that stage~H makes untrue beyond
    those of the items (Section~\ref{sec:regionEarlier}), among them the stage list of the
    introduction, the skeleton of a run, the summaries of the C code, the remark after the remainder
    chain of Chapter~\ref{chap:setup}, the paragraph after Theorem~\ref{thm:findPoints},
    Section~\ref{sec:certifiedG}, the counts of the differential tests of Chapter~\ref{chap:spec}
    and the record of the lines that Chapter~\ref{chap:output} left free for stage~H; the convention
    that \verb|Id.run| appears only in $\mathrm{findPoints}$, $\mathrm{findPointsWith}$, the
    program's run and its pure callback, \verb|Main.lean| and \verb|#guard| lines; and the generated
    table.
\end{enumerate}
H.5 to H.10 keep the exact comparison of stage~G under \verb|-x|, with \verb|-s| appended (from H.7
on with the weaker check on the exact region and the range of the denominators and with the check of
the moduli of the C program), and drop message~8 from the report of the C program, since the model
narrows the search intervals by the Sturm step only from H.11 on. The alternative order, the Sturm
step before the search region of the options, would change the certified statements once, in one
large item; the order above keeps the items of the region and of the step apart and brings each
certified change with the theorem that proves it.

Names that the prototype has and the audit of the stage is to review: the lemmas named \verb|spec_*|
that are not contracts of the program logic (\lname{Ratpoints.Double.spec\_ofDyadic},
\lname{Ratpoints.Sieve.spec\_intervalBitsExact} and their kin), \verb|m_run_exactEnd|,
\verb|s_run_exactEnd|, \verb|toNat_rounding|, \verb|mod_add_carry|, \verb|gap_tail|,
\verb|exists_left| and \verb|exists_right|, $\mathrm{PosAssoc}$, $\mathrm{Kept}$,
$\mathrm{Witnessed}$, the step lemmas of the invariants of the program, and the two declarations
named \verb|intervalsWF_domainBits_iff|, \lname{Ratpoints.intervalsWF\_domainBits\_iff} (the indexed
form, in \verb|Proof/Bounds.lean|) and \lname{Ratpoints.Sturm.intervalsWF\_domainBits\_iff} (the
form of the bisection, in \verb|Proof/Bisection.lean|); \lname{Ratpoints.Cli.Options.store}, which
acts on $\mathrm{Options}$ but lives in the namespace \verb|Ratpoints.Cli|, so that dot notation
does not reach it; and the general facts of \verb|Proof/Bisection.lean| and \verb|Proof/Output.lean|
that may belong in a general file. \lname{Ratpoints.searchDens\_eq\_range'} and
\lname{Ratpoints.append\_flatMap\_searchDens\_eq\_bruteForce} stay, since
Chapter~\ref{chap:denominators} cites them. The names of this chapter are those of the prototype;
the audit may rename them.

The items also give the docstrings of the prototype the precision of this chapter: those of
\verb|Ratpoints.Cli.Failure| and \verb|Ratpoints.Cli.Failure.ofRefusal| say that error~9 is for the
codes of the library other than \verb|RATPOINTS_NON_SQUAREFREE| and \verb|RATPOINTS_NO_MEMORY|,
error~10, which the model cannot return (H.5); the docstrings that say ``closed'' for $l \le u$ (in
\verb|Ratpoints/Spec/Region.lean| and its restatement in the challenge,
\verb|Ratpoints/Model/FindPoints.lean|, \verb|Ratpoints/Proof/Bisection.lean|,
\verb|Ratpoints/Proof/Cli.lean| and \verb|Ratpoints/Proof/CliNumbers.lean|) say $l \le u$, the
refusal of the search intervals reading ``without $l \le u$, not increasing and strictly apart, or
with a NaN end'', with the item that brings each; the module docstring of
\verb|Ratpoints/Proof/Double.lean| states the exactness of the nearest double with the hypotheses of
Lemma~\ref{lem:Double} (H.1); the docstring of
\verb|Ratpoints.Cli.mem_filter_keeps_iff_of_input_eq_ok| names the points at infinity unless
\verb|-i| and the affine points in the range of \verb|-dl| and \verb|-du| and in the intervals of
\verb|-l| and \verb|-u| when there are some (H.7); and the harness gives, beside the composite
modulus $9$ for $3$, the subset of the $30$ primes and the moduli $55$ and $63$ as the cases where
the C program chooses other moduli under the pin (H.7).

\section{The evidence}\label{sec:region-evidence}

\begin{itemize}
  \item A reading of the C code of the search region, the Sturm step, the bounds of the numerators,
    the normalizations and the refusals, and of the options of the region, \verb|scan_end|,
    \verb|outward|, \verb|join_intervals|, \verb|in_region| and messages~3 and~8 of \verb|main.c|,
    at the pin, by three readers, each with a second reader who checked the findings, with Python
    models and drivers of the C code run against the pinned program, $0$ mismatches in all: the
    reading of \verb|-l| and \verb|-u|, the doubles of \verb|outward| and message~3 on $9000$
    command lines (ties, $-0$, the boundaries $2^{-1075}$ and $2^{1024} - 2^{970}$, hexadecimal
    words, infinities and NaNs, malformed words, heights up to $2^{63} - 1$); \verb|sturm.c| with
    the default interval, the clamp, \verb|RATPOINTS_NO_REVERSE_AUTO| with the reversal and
    message~8 on $10500$ runs (degrees $1$ to $12$, coefficients up to $10^{25}$, depths $0$ to
    $1000$ and \verb|-s|), the bound $1 + \lfloor d/2 \rfloor$ of the array of the bisection
    asserted on all of them; the bits of intervals against the build under \verb|DEBUG| on $2823$
    denominators; the points in their order with regions against brute force on $2985$ runs and in
    the divisor order with \verb|-dl| and \verb|-du| on $1495$; and no defect of the C code. The
    reading recorded the behaviours that the model follows: $-0$ printed as \verb|-0.000000|, an
    interval beyond the height bound kept as \verb|[-inf, -20]|, \verb|[inf, inf]| accepted, the
    order of the refusals, the survivors under \verb|-x| depending on the depth, a form negative
    everywhere printing nothing, the suppression of the reversal by \verb|-dl| and \verb|-du|, the
    path of \verb|sturm_empty|, and \verb|-S 9223372036854775807| clamped to $55$. A reading of the
    Lean side proved, before the stage, that the rows of the chain are positive multiples of
    Mathlib's Sturm sequence, checked the proposition of Sturm's theorem exactly on $30000$
    instances with ends at roots, and checked a model of the bisection and the intersection with
    exact dyadic ends against a driver of \verb|_ratpoints_compute_sturm| on $500$ and $600$ runs.
  \item Before the stage, a sketch of the model in Lean in two parts, the library and the program,
    checked against the C program: the default suites with the new suites \verb|regions| and
    \verb|depths|, $162$ of the $165$ region runs of \verb|test4.sh| and the one of \verb|test3.sh|,
    the random suites on two more seeds, $11000$ whole outcomes of the search with a callback
    against \verb|find_points| of the library, with the count, the flags, the points in their order
    and the bits of the search intervals, $2000$ Sturm steps bit for bit, $4558$ comparisons of
    message~8, $5998$ ends, scaled ends and forms of random doubles, $465$ denominators of
    intervals, and \verb|strtod| and \verb|%f| against glibc on $46000$ lines, none failing. It
    found a mismatch that stage~H did not cause and stage~I resolves: under \verb|-x -p 30 -N 30|
    the C program can choose a composite modulus, so that the pin does not fix its moduli.
  \item Four reviews of the sketch, of the library against the C code, of the program, of the
    harness and of the specification, each finding checked by a second reader: $42$ findings, $37$
    confirmed, $4$ in part and $1$ refuted, no defect of the behaviour of the model. They made the
    Sturm step own the clears of its table, \verb|outward| and \verb|scan_end| own the lines of
    their GMP numbers that carry meaning, the specification speak about the bits and the values of
    the doubles, the theorems without Sturm's theorem a filter by the intervals that the step
    leaves, and gave the harness the check of the moduli that the C program uses under the pin, the
    exact region of the command line in the weaker check, the bound of the numerators under
    \verb|-1| and the admission of large heights by the number of the denominators, by which the
    three runs of \verb|test4.sh| under \verb|-1| at $10^{12}$ and $2^{63} - 1$ are compared.
  \item A prototype of all the proofs on the sketch: the whole library, the program, the challenge
    and the solution build, with no \verb|sorry| outside the challenge, every proof at most $20$
    lines with no new exception of the style lint, standard axioms, $\mathrm{SturmTheorem}$ a
    hypothesis; comparator with \verb|--paranoid| accepts the twelve statements with all six
    kernels, all of them stated at once in the prototype, where the order of work brings each change
    with its theorem; the build with other values of the constants passes its nine variants, the
    decided facts on $\mathrm{maxDegree}$ among them; the differential tests compare $4530$ times
    with none failed and $13$ times in \verb|test4high|; it adds $823$ public names to those of the
    sketch. It changed the model of the sketch, as the reviews left it, only by the two decided
    facts on $\mathrm{maxDegree}$, and it showed which hypotheses the statements need: the depth at
    most $55$ for the intervals of the step, $d < 2^{57}$ and $|a| \le 2^{63}$ for its coverage,
    $|x| \le H$ and words of fewer than $2^{56}$ characters for the program; it proved the roundings
    of \verb|strtod| faithful, not nearest, and the doubles of upper ends at or above the end, not
    the least such, which is all that the theorems need.
\end{itemize}
